\documentclass[11pt,reqno]{amsart}
%version submitted to arXiv Nov 6 2026
\usepackage[T1]{fontenc}
\usepackage[utf8]{inputenc}
\usepackage{lmodern}
\usepackage{amsmath,amssymb,amsthm,mathtools,mathrsfs}
\usepackage[expansion=false]{microtype}
\usepackage[a4paper,margin=28mm]{geometry}
\usepackage{xcolor}
\usepackage{enumitem}
\usepackage{booktabs}
\usepackage[hidelinks]{hyperref}
% cleveref must be loaded last, after hyperref; moving it breaks the build.
\usepackage[capitalise,nameinlink]{cleveref}
\makeatletter
\AtBeginDocument{%
  \@ifundefined{firstaid@cref@updatelabeldata}{}{%
    \gdef\cref@fix@ctr{}\gdef\cref@fix@type{}%
    \let\cref@fix@orig@optarg\refstepcounter@optarg
    \def\refstepcounter@optarg[#1]#2{%
      \gdef\cref@fix@ctr{#2}\gdef\cref@fix@type{#1}%
      \cref@fix@orig@optarg[#1]{#2}}%
    \let\cref@fix@orig@noarg\refstepcounter@noarg
    \def\refstepcounter@noarg#1{%
      \gdef\cref@fix@ctr{#1}\gdef\cref@fix@type{#1}%
      \cref@fix@orig@noarg{#1}}%
    \renewcommand\firstaid@cref@updatelabeldata[1]{%
      \cref@constructprefix{#1}{\cref@result}%
      \edef\cref@fix@tmp{#1}%
      \ifx\cref@fix@tmp\cref@fix@ctr
        \let\@tempa\cref@fix@type
      \else
        \@ifundefined{cref@#1@alias}%
          {\def\@tempa{#1}}%
          {\edef\@tempa{\csname cref@#1@alias\endcsname}}%
      \fi
      \protected@xdef\cref@gcurrentlabel@temp{%
        [\@tempa][\arabic{#1}][\cref@result]%
        \csname p@#1\expandafter\endcsname\csname the#1\endcsname}%
      \aftergroup\firstaid@cref@smugglelabel}%
  }%
}
\makeatother
\crefformat{enumi}{(#2#1#3)}
\Crefformat{enumi}{(#2#1#3)}
\crefrangeformat{enumi}{(#3#1#4)--(#5#2#6)}
\crefmultiformat{enumi}{(#2#1#3)}{ and~(#2#1#3)}{, (#2#1#3)}{ and~(#2#1#3)}
\Crefmultiformat{enumi}{(#2#1#3)}{ and~(#2#1#3)}{, (#2#1#3)}{ and~(#2#1#3)}

\setlength{\emergencystretch}{2em}
\setlist[enumerate]{itemsep=2pt,topsep=4pt}

% Legacy AMS macro used by some MathSciNet entries and not defined by
% any @preamble block.
\providecommand{\germ}{\mathfrak}

\numberwithin{equation}{section}
\theoremstyle{plain}
\newtheorem{theorem}{Theorem}[section]
\newtheorem{proposition}[theorem]{Proposition}
\newtheorem{lemma}[theorem]{Lemma}
\newtheorem{corollary}[theorem]{Corollary}
\theoremstyle{definition}
\newtheorem{definition}[theorem]{Definition}
\theoremstyle{remark}
\newtheorem{remark}[theorem]{Remark}
\newtheorem{example}[theorem]{Example}
% Informal statements in the introduction, numbered A, B, ...
\theoremstyle{plain}
\newtheorem{introtheorem}{Theorem}
\renewcommand{\theintrotheorem}{\Alph{introtheorem}}
\crefname{introtheorem}{Theorem}{Theorems}


% Number systems. Never write \mathbb{...} for these in the body.
\newcommand{\N}{\mathbb{N}}      % the natural numbers
\newcommand{\Z}{\mathbb{Z}}      % the integers
\newcommand{\Q}{\mathbb{Q}}      % the rationals
\newcommand{\R}{\mathbb{R}}      % the reals
\newcommand{\C}{\mathbb{C}}      % the complex numbers
\newcommand{\Gm}{\mathbb{G}_m}   % the multiplicative group
\newcommand{\Fq}[1]{\mathbb{F}_{#1}}  % the finite field with #1 elements

% Standing notation. Redefine here to change it throughout.
\newcommand{\OO}{\mathcal{O}}    % ring of integers of F
\newcommand{\pp}{\mathfrak{p}}   % its maximal ideal
% The two standing lattices (\cref{sec:standardgroups}): a lattice in \g and
% its dual lattice in \gd. Redefine to change the symbols throughout; decorated
% lattices (\Lat_P, \Lat_\pm, \dLat_i, ...) are built from these macros.
\newcommand{\Lat}{\mathcal{L}}   % lattice in \g
\newcommand{\dLat}{\Lambda}      % lattice in \gd (the dual of \Lat)
\newcommand{\Aff}{\mathbb{A}}    % affine space
\newcommand{\bG}{\mathbf{G}}     % the connected reductive group over F
\newcommand{\bA}{\mathbf{A}}     % a split torus
\newcommand{\bP}{\mathbf{P}}     % a parabolic subgroup of \bG
\newcommand{\bQ}{\mathbf{Q}}     % another parabolic subgroup of \bG
\newcommand{\bM}{\mathbf{M}}     % a Levi subgroup
\newcommand{\bU}{\mathbf{U}}     % a unipotent radical
\newcommand{\bT}{\mathbf{T}}     % a maximal torus
\newcommand{\bH}{\mathbf{H}}     % a group of type B in characteristic 2 (\cref{lem:oddorthogonal})
% The odd orthogonal and ordinary parts of \bG (\cref{lem:split}), \bG=\Godd\times\Gord,
% and their Lie algebras.
\newcommand{\Godd}{\bG_{\mathrm{odd}}}
\newcommand{\Gord}{\bG_{\mathrm{ord}}}
\newcommand{\godd}{\g_{\mathrm{odd}}}
\newcommand{\gord}{\g_{\mathrm{ord}}}
\newcommand{\g}{\mathfrak{g}}    % the Lie algebra of G
\newcommand{\gd}{\mathfrak{g}^*} % its dual
\newcommand{\uu}{\mathfrak{u}}   % Lie algebra of a unipotent radical
\newcommand{\m}{\mathfrak{m}}    % Lie algebra of a Levi subgroup
\newcommand{\liet}{\mathfrak{t}}  % Lie algebra of a torus
\newcommand{\lieh}{\mathfrak{h}}  % Lie algebra of \bH
% Elliptic linear forms on a Lie algebra #1, and their complement (the
% forms vanishing on the Lie algebra of the unipotent radical of some proper
% parabolic subgroup).
% The optional argument is an index written before "ell": \Ell[0,]{\g^*}.
\newcommand{\Ell}[2][]{#2_{#1\mathrm{ell}}}   % use with the dual Lie algebra: \Ell{\gd}
\newcommand{\Par}[2][]{#2_{#1\mathrm{par}}}
% Densities: supercuspidal density of a compact open subgroup, elliptic
% density of a lattice in the dual Lie algebra.
\newcommand{\scd}[1]{\delta_{\mathrm{sc}}(#1)}
\newcommand{\fdeg}[1]{d_{#1}}  % formal degree of a discrete series #1, a Haar measure on G
\newcommand{\elld}[2][]{\delta_{\mathrm{ell}}^{#1}(#2)}  % optional argument: the group
\newcommand{\Bad}{\mathcal{B}}   % the boundary nets \Bad_h (\cref{thm:sard})
% The names of the sets \Bad_h and of the sequence (\Bad_h)_h (\cref{thm:sard}).
% Redefine these to rename the notion throughout; always write them with braces.
\newcommand{\bnet}{boundary net}
\newcommand{\bnets}{boundary nets}
\newcommand{\Bnets}{Boundary nets}
\newcommand{\Hk}{\mathcal{H}}    % Hecke algebra
\newcommand{\Pc}{\mathscr{P}^{\mathrm{st}}_{\mathrm{max}}}    % representatives of maximal parabolics
\newcommand{\Pcss}{\mathscr{P}^{\mathrm{ss}}_{\mathrm{max}}}  % their semistandard conjugates (\cref{sec:ellipticforms})
\newcommand{\Nss}{N_{\bG}}  % the number |\Pcss| of maximal semistandard parabolics (\cref{sec:parabolics})
\newcommand{\Fl}{\mathcal{F}}    % flag manifold \Fl_P=\bFl_P(F)
\newcommand{\bFl}{\boldsymbol{\Fl}}  % flag variety \bFl_P=\bG/\bP
\newcommand{\sS}{\mathfrak{s}^+}   % enlarged positive lattice
\newcommand{\sT}{\mathfrak{s}^-}   % enlarged negative lattice
\newcommand{\Gsc}{\widehat{G}_{\mathrm{sc}}}    % supercuspidal unitary dual
\newcommand{\Gtemp}{\widehat{G}_{\mathrm{temp}}} % tempered unitary dual
\newcommand{\Pps}{\delta_{\mathrm{ps}}(K_n)}       % principal series density (\cref{sec:gl2exact})
\newcommand{\Pst}{\delta_{\mathrm{St}}(K_n)}
% Normalization constants for principal-series and Steinberg mass in \cref{sec:gl2exact}.
% Redefine these to rename the constants throughout the calculations.
\newcommand{\Cps}{c_{\mathrm{ps}}}
\newcommand{\Cst}{c_{\mathrm{St}}}
\newcommand{\Cmu}{c_\mu}
% Probability of a nonsquare discriminant in the elementary GL2 calculation.
\newcommand{\Pns}{p_{\mathrm{ns}}}
\newcommand{\St}{\mathrm{St}}                   % Steinberg representation
\newcommand{\Iw}{\mathcal{I}}                   % Iwahori subgroup of GL_2
\newcommand{\Xc}{\mathcal{X}}                   % unitary characters of F^\times
\newcommand{\cG}{\mathcal{G}}                   % smooth (in the unramified case, \cref{sec:integraldata}: reductive) group scheme over O
\newcommand{\cP}{\mathcal{P}}                   % parabolic subgroup scheme
\newcommand{\cM}{\mathcal{M}}                   % Levi subgroup scheme
\newcommand{\cU}{\mathcal{U}}                   % unipotent radical scheme
\newcommand{\cT}{\mathcal{T}}                   % maximal torus scheme
\newcommand{\cB}{\mathcal{B}}                   % Borel subgroup scheme
\newcommand{\Bldg}{\mathcal{B}}                 % Bruhat--Tits building
\newcommand{\cS}{\mathcal{S}}                   % maximal split torus scheme
\newcommand{\one}{\mathbf{1}}
\newcommand{\der}{\operatorname{der}}

\newcommand{\abs}[1]{\lvert#1\rvert}
\newcommand{\norm}[1]{\lVert#1\rVert}
% Compatibility with the defining cocharacter; redefine to rename throughout.
\newcommand{\etacompat}{$\eta$-compatible}

\DeclareMathOperator{\vol}{vol}
\DeclareMathOperator{\tr}{tr}
\DeclareMathOperator{\diag}{diag}
\DeclareMathOperator{\Ad}{Ad}
\DeclareMathOperator{\ad}{ad}
\DeclareMathOperator{\Hom}{Hom}
\DeclareMathOperator{\val}{val}
\DeclareMathOperator{\Lie}{Lie}
\DeclareMathOperator{\GL}{GL}
\DeclareMathOperator{\PGL}{PGL}
\DeclareMathOperator{\SL}{SL}
\DeclareMathOperator{\Gal}{Gal}
\DeclareMathOperator{\SO}{SO}
\DeclareMathOperator{\Mat}{M}
\DeclareMathOperator{\Crit}{Crit}
\DeclareMathOperator{\Ind}{Ind}
\DeclareMathOperator{\rad}{rad}
\DeclareMathOperator{\id}{id}
\DeclareMathOperator{\chr}{char}   % characteristic of a field
\DeclareMathOperator{\Irr}{Irr}
\DeclareMathOperator{\End}{End}
\DeclareMathOperator{\disc}{disc}
\DeclareMathOperator{\Spec}{Spec}
\DeclareMathOperator{\cInd}{c-Ind}  % compact induction
\DeclareMathOperator{\cond}{cond}   % conductor exponent of a character of F^\times

% The condition of \cref{def:peculiar} (some incidence map is everywhere
% critical), which governs the exponent \eqref{eq:tau}. Redefine to rename
% the notion throughout; write \peculiar{}.
\newcommand{\peculiar}{peculiar}
\newcommand{\Peculiar}{Peculiar}
% The condition of \cref{def:ordinary} (Chi's ``coroot-unramified''); \Gord
% in \cref{lem:split} is the \ordinary{} part. Redefine to rename the notion
% throughout; write \ordinary{}.
\newcommand{\ordinary}{ordinary}
\newcommand{\Ordinary}{Ordinary}
% The factor \Godd of \cref{lem:split}. Redefine to rename it throughout;
% write \oddorth{}.
\newcommand{\oddorth}{odd orthogonal}
\newcommand{\Oddorth}{Odd orthogonal}
% The error exponent of \eqref{eq:tau} and the dimension d_G entering it.
% Redefine to change the symbols throughout.
\newcommand{\eexp}{\tau}
\newcommand{\dG}{d_{\bG}}
% Ellipticity radius of \lf in \dLat (defined before \cref{prop:elliptic}).
% Redefine to change the symbol throughout.
\newcommand{\erad}{R_{\mathrm{ell}}}
% Standing notation: a typical linear form in \gd, and a typical cell, i.e., a
% coset of \pp^r\dLat in \dLat (or a union of such cosets). Redefine to rename
% throughout.
\newcommand{\lf}{\lambda}
\newcommand{\cell}{\omega}
% The bilinear form \Bf_\lf(X,Y)=\lf([X,Y]) on \uu\times\uu^- of \eqref{eq:levipairing}.
\newcommand{\Bf}{\mathsf{B}}
% Averaging projectors in the Jacquet injectivity criterion. Redefine to
% change their notation throughout that argument.
\newcommand{\Avg}{\mathsf{P}}
% Normalized charts (\cref{sec:standardgroups}). Redefine to rename the notion
% throughout; write \nchart{} with braces.
\newcommand{\nchart}{normalized chart}
\newcommand{\npar}{normalized parametrization}  % the inverse of a normalized chart
% A linear form in normal position (vanishing on \uu_P), cf. \cref{sec:bad}.
\newcommand{\nf}{\zeta}
% The incidence manifolds \Inc_P and their disjoint union \Inc (\cref{sec:incidencemaps}),
% and the incidence varieties \bInc_P, \bInc of which they are the F-points.
\newcommand{\Inc}{\mathcal{Y}}
\newcommand{\bInc}{\boldsymbol{\Inc}}
% The critical values of the incidence map p (\cref{sec:incidencemaps}), or of the
% map f of \cref{sec:sard}. Redefine to change the symbol throughout.
\newcommand{\Crv}{\Sigma}
% The image f(X) and the compact open set f^{-1}(\dLat) of \cref{sec:bnets}.
\newcommand{\Vf}{Y}
\newcommand{\XL}{X_{\dLat}}
% Closed polydiscs in F^D (\cref{sec:sard}).
\newcommand{\Incc}[1]{\Inc^{\circ}_{#1}}  % big cell of \Inc_P (\cref{sec:levi})
% Families of charts as before \cref{thm:sard}; write \iatlas{}.
\newcommand{\iatlas}{integral atlas}
\newcommand{\Iatlas}{Integral atlas}
\begin{document}

\title[Supercuspidal and elliptic densities]
{Supercuspidal and elliptic densities for reductive groups over non-archimedean local fields}

\author{Erez Lapid}
\address{Department of Mathematics, Weizmann Institute of Science, Rehovot 7610001, Israel}
\email{erez.lapid@weizmann.ac.il}

\author{Mark Shusterman}
\address{Department of Mathematics, Weizmann Institute of Science, Rehovot 7610001, Israel}
\email{mark.shusterman@weizmann.ac.il}

\date{November 6, 2026}
\subjclass[2020]{22E50, 22E35}
\keywords{Supercuspidal representations, Plancherel measure,
congruence subgroups, coadjoint orbits}

\begin{abstract}
Let $\bG$ be a connected reductive group over a non-archimedean local
field $F$ of arbitrary characteristic, with residue field of cardinality
$q$, and let $(K_n)$ be a tower of principal congruence subgroups attached
to a lattice $\Lat$ in the Lie algebra. We show that the proportion of the
Plancherel mass of the $K_n$-fixed vectors carried by supercuspidal
representations is $\elld{\dLat}+O(q^{-n/4})$, where the elliptic density
$\elld{\dLat}>0$ is the probability that a random element of the dual
lattice $\dLat$ is nonzero on the Lie algebra of the unipotent radical
of every proper parabolic subgroup. For $\GL_d$ this limit was anticipated by
Howe. The error improves to $O(q^{-\eexp n/4})$ with an explicit
$\eexp\in(1,2]$, except when $\chr F=2$ and $\bG$ has an isotropic direct
factor which over $\bar F$ is a product of adjoint groups of type
$\mathrm B_d$. We determine when $\elld{\dLat}=1$, give an explicit bound
uniform in $F$ for reductive models, and deduce a bound uniform over the
completions of a global field. The proof combines finite Fourier analysis
on congruence quotients, a quantitative Sard theorem for parabolic
incidence maps and a Jacquet-module argument; it uses neither a
classification of supercuspidal representations nor the exponential map.
\end{abstract}

\maketitle

% Table of contents at the section level only; raise tocdepth to 2 to
% list subsections as well.
\setcounter{tocdepth}{1}
\tableofcontents

\section{Introduction}\label{sec:intro}

Let $\bG$ be a connected reductive group over a non-archimedean local
field $F$ of arbitrary characteristic, with ring of integers $\OO$,
maximal ideal $\pp$ and residue field of cardinality $q$, and let
$G=\bG(F)$. For a compact open subgroup $K$ of $G$, the Plancherel
formula distributes the mass of the $K$-fixed vectors over the tempered
dual of $G$. We show that as $K$ runs through a tower of principal
congruence subgroups, the proportion of this mass carried by
supercuspidal representations converges, at a rate which is a
power of $q^{-n}$, to a purely geometric quantity: the probability that a
random linear form in a lattice of the dual Lie algebra is elliptic. For
$\GL_d$ the question goes back to Howe \cite{MR0492088} (see
\cref{sec:related}). Our argument uses no construction or classification of
supercuspidal representations, and it applies to every $\bG$ and $F$,
including wildly ramified groups and small residual characteristic.

\subsection{The main result}

For simplicity of presentation, assume in this introduction that the
center of $\bG$ is anisotropic, so that the supercuspidal representations
are discrete in the tempered dual and have formal degrees. Let $\cG$ be a
smooth affine group scheme over $\OO$ with generic fiber $\bG$, let
\[
K_n=\ker\bigl(\cG(\OO)\to\cG(\OO/\pp^n)\bigr),\qquad n\ge1,
\]
be its principal congruence subgroups, let $\Lat=\Lie(\cG)(\OO)$, a
lattice in $\g=\Lie(\bG)(F)$, and let $\dLat\subset\gd$ be the dual
lattice. A linear form $\lf\in\gd$ is \emph{elliptic} if
$\lf|_{\uu_P}\ne0$ for every proper parabolic $F$-subgroup $\bP$ of $\bG$,
where $\uu_P$ is the Lie algebra of the unipotent radical of $\bP$. The
\emph{elliptic density} $\elld{\dLat}$ of $\dLat$ is the proportion of
elliptic forms in $\dLat$, with respect to Haar measure. For a
supercuspidal representation $\pi$ let $\fdeg{\pi}$ be its formal degree,
viewed as a Haar measure on $G$.

\begin{introtheorem}\label{thm:A}
We have $\elld{\dLat}>0$ and, as $n\to\infty$,
\[
\sum_\pi\fdeg{\pi}(K_n)\dim\pi^{K_n}=\elld{\dLat}+O(q^{-n/4}),
\]
where $\pi$ runs over the isomorphism classes of irreducible supercuspidal
representations of $G$.
\end{introtheorem}

By the Plancherel formula, the left-hand side is the proportion
$\scd{K_n}$ of the Plancherel mass of the $K_n$-fixed vectors carried by
supercuspidal representations. Thus, asymptotically, the supercuspidal
part of $L^2(G/K_n)$ is measured by the probability that a random element
of $\dLat$ is elliptic.

\Cref{thm:main} is more precise. It allows an arbitrary center, in which
case $\scd{K_n}$ is defined through the Plancherel measure, and arbitrary
towers $K_n=\kappa^{-1}(\pp^n\Lat)$ defined by a lattice $\Lat\subset\g$
and a chart $\kappa$ of $G$ near $1$ (\cref{sec:standardgroups}), for $n$
large. These include towers attached to lattices which are not Lie
algebras of group schemes, and parahoric towers (\cref{rem:badplaces}). The
error term is $O(q^{-\eexp n/4})$ with $\eexp=1+1/\dG{}\in(1,2]$, where
$\dG{}$ is the largest dimension of the unipotent radical of a maximal
proper parabolic subgroup ($\dG{}=1$ if there is none), except for the \peculiar{} groups of
\cref{def:peculiar}, which occur only in characteristic $2$ and for which
$\eexp=1$. Finally, a localized version holds for the supercuspidal mass
of the part of $e_{K_n}$ attached to a subset of $\dLat$
(\cref{cor:localized}).

\subsection{When is the elliptic density equal to one?}

Outside characteristic $2$, the supercuspidal representations carry
asymptotically all the Plancherel mass only in the trivial case where there are
no proper parabolic $F$-subgroups. In
characteristic $2$ there is a genuine exception, coming from adjoint groups
of type $\mathrm B_d$ (with $\mathrm B_1=\mathrm A_1$). Every $\bG$ is
uniquely a product $\Godd\times\Gord$, where $\Godd$ becomes a product of
such groups over $\bar F$ and $\Gord$ is \ordinary{} in the sense of
\cref{def:ordinary}; if $\chr F\ne2$ then $\Godd=1$ (\cref{lem:split}).

\begin{introtheorem}\label{thm:B}
The elliptic density is $1$ (for one, equivalently for every, lattice) if
and only if $\Gord$ is anisotropic modulo its center. In particular, if
$\chr F\ne2$, this happens if and only if $\bG$ is anisotropic modulo its
center.
\end{introtheorem}

This is \cref{thm:endpoint}. The simplest example is $\PGL_2$ in
characteristic $2$: there $[E_{12},E_{21}]=E_{11}-E_{22}$ is the identity,
hence zero in $\mathfrak{pgl}_2$, the incidence map (see below) is
everywhere critical, and the non-elliptic forms are a null set although
$\PGL_2$ is split. The error does not vanish: by an exact computation,
$1-\scd{K_n}\sim\frac12q^{-\lfloor n/2\rfloor}$ (\cref{prop:pgl2}).

\subsection{Examples}
The limit depends on the lattice and not only on the group. For
$\bG=\GL_2$ and $\Lat=\Mat_2(\OO)$ the elliptic density is $\frac12$ in
every characteristic (\cref{prop:gl2}), while the Plancherel formula gives
\[
\scd{K_n}=\frac12-q^{1-2n}+\frac1{(q+1)q^{3n-2}}
\]
for every $F$ and $n\ge1$ (\cref{prop:gl2exact}). Both computations are
independent of the proof of \cref{thm:A}. For the Lie algebra of an
Iwahori subgroup the limit is $\frac q{q+1}$, and for its dual with respect
to the trace pairing it is $\frac1{q+1}$ (\cref{prop:iwahori}).

\subsection{Complements}
When $\cG$ is a reductive model of $\bG$, the constants can be made
explicit in terms of the root datum, uniformly in $F$ and its
characteristic (\cref{thm:uniform}). For instance, for $\bG=\GL_d$,
$\Lat=\Mat_d(\OO)$, every non-archimedean local field $F$ and $n\ge13$,
\[
\bigl|\scd{K_n}-\elld{\dLat}\bigr|\le(2^d-2)\,q^{3-n/4}.
\]
Consequently, for a reductive group over a global field with a fixed
linear embedding, the estimate holds with a constant independent of the
place (\cref{thm:global}).

Under the compact-center assumption above, the proof also bounds the size of the supercuspidal part of the Hecke
algebra of $K_n$-biinvariant functions from below:
\[
[\cG(\OO):K_n]^{-1}\sum_{\pi}(\dim\pi^{K_n})^2\ge\elld{\dLat}-O(q^{-n/4})
\]
(\cref{thm:squares}). This does not follow formally from \cref{thm:A}, which gives
no information on individual formal degrees.

\subsection{Idea of the proof}
The proof has a geometric input and a representation-theoretic input,
organized by Howe's method of analyzing
representations through the dual blobs of shallow characters
\cite[\S2]{MR0492088}. For $n\ge2r$, the characters of the abelian group
$K_{n-r}/K_n$ are parametrized by the cosets, or \emph{cells}, in
$\dLat/\pp^r\dLat$, and the Plancherel mass of the $K_n$-fixed vectors
decomposes accordingly (\cref{sec:cells}). Let $f_{n,r}$ be the function on
$\dLat$, constant on cells, which records the supercuspidal proportion of
the mass of each cell; then $\scd{K_n}$ is the integral of $f_{n,r}$ over
$\dLat$, and \cref{thm:A} says that $f_{n,r}$ is close to the indicator
function of the elliptic set in $L^1$.

The geometric input is the \emph{incidence map} $p$: for each maximal
proper parabolic subgroup $\bP$, the proper map
$G\times^P\uu_P^\perp\to\gd$, $[g,\nf]\mapsto\Ad^*(g)\nf$, whose images form
the complement of the elliptic set. A rudimentary quantitative Sard theorem
for $F$-analytic maps (\cref{sec:sard}) produces \bnets{} $\Bad_h\subset\dLat$,
unions of cosets of $\pp^{2h}\dLat$ of measure $O(q^{-\eexp h})$ shrinking to
the critical values of $p$, outside of which every coset of
$\pp^{2h}\dLat$ is either elliptic or has a parabolic lift with
quantitatively non-degenerate pairing (\cref{prop:geometric}). The exponent
$\eexp$ comes from the degrees of the Jacobian polynomials.

The representation-theoretic input treats the two alternatives. A cell well
inside the elliptic set occurs only in supercuspidal representations, by a
standard argument due to Moy and Prasad \cite[Corollary~4.9]{MR1371680},
of which Beuzart-Plessis \cite{MR3463423}, \cite[Appendix~D]{MR4325284} gave a version in
characteristic zero using the exponential map
(\cref{prop:elliptic}). Conversely, a cell with a quantitatively regular
parabolic lift occurs in no supercuspidal representation
(\cref{thm:regular}); this is the technical heart of the paper and is
closely related to the theory of covers of Bushnell and Kutzko
\cite{MR1643417}. We extend the cell character to a larger compact open
subgroup, normalize the extension by conjugation in the opposite unipotent
radical, and show, by averaging along a contracting exhaustion of the
unipotent radical, that its isotypic vectors survive in the Jacquet module.
Neither part uses the exponential map. The depth conditions needed in the
two parts force $h\le n/4+O(1)$, which explains the exponent $n/4$.
Positivity of the elliptic density and \cref{thm:B} are proved in
\cref{sec:positivity}, using elliptic tori and an explicit computation for
odd orthogonal groups in characteristic $2$.

\subsection{Relation to earlier work}\label{sec:related}

The point of view of this paper goes back to Howe's work on $\GL_d$
over a non-archimedean local field of arbitrary characteristic
\cite{MR0492088}, whose dual blobs of shallow characters are essentially
our cells (\cref{sec:cells}). For a sufficiently regular
shallow character $\psi$, Howe's matching formula identifies the
Plancherel integral of the multiplicity of another shallow character
$\delta$ over the representations containing $\psi$ with the Haar
measure of the intersection of the dual blob of $\delta$ with the
$\Ad^*(G)$-orbit of that of $\psi$, up to normalization
\cite[Theorem~4.1]{MR0492088}. These representations are supercuspidal
when the attached Cartan subalgebra is anisotropic, and induced from a
proper parabolic otherwise \cite[\S3]{MR0492088}.
Howe indicated, without proof, that these families exhaust most of the
Plancherel mass as the level grows, which would yield the limit in
\eqref{eq:main} for $\GL_d$ with the same limit $\elld{\dLat}$
\cite[\S4, (1)--(3)]{MR0492088}. His question about a partition of the
unitary dual matched with invariant subsets of $\gd$
\cite[end of \S4]{MR0492088} motivates the quantitative comparison in
\cref{prop:elliptic,thm:regular} and \eqref{eq:localized}.

For general $\bG$, supercuspidal representations are known to be exhausted
by Yu's construction \cite{MR1824988} (see also \cite{MR4357723}) under
tameness assumptions on $\bG$ and $p$, by work of Kim \cite{MR2276772} and
Fintzen \cite{MR4199732}. In characteristic zero, for large $p$ and
tamely ramified $\bG$, Kim's formulas matching the Plancherel measure of
families defined by good minimal $K$-types
\cite[Theorem~3.3.1]{MR2075916} or Yu's data
\cite[Theorem~10.1]{MR2276772} with the Haar measure of their dual
blobs near the identity, together with Kim's exhaustion theorems
\cite[Theorems~11.1 and~19.1]{MR2276772}, suggest the qualitative limit
$\scd{K_n}\to\elld{\dLat}$.\footnote{Under the hypotheses (Hk), (HB),
(HGT) and (HN) of \cite[\S3.4]{MR2276772}.} We have not carried out this
deduction, and a rate of convergence would require additional uniformity
not contained in these statements.
Fintzen's exhaustion theorem \cite[Theorem~8.1]{MR4199732} applies in
every characteristic when $\bG$ splits tamely and $p$ does not divide
the order of the Weyl group; Kaletha's regular supercuspidal
representations, which in general form a proper subclass of those
exhausted there, are parametrized by tame elliptic regular pairs
\cite[Proposition~3.7.14]{MR4013740}. The matching formulas still use
character expansions, going back for $\GL_d$ to Howe
\cite{MR0342645}, and the exponential map, and require characteristic
zero and large $p$.

The contribution of \cref{thm:main} is thus twofold. It imposes no
additional restrictions on $F$, $p$, $\bG$ or $\Lat$; in particular,
it includes the cases where $p$ divides the order of the Weyl group
or $\bG$ is wildly ramified. Its proof combines finite Fourier analysis
on $K_{n-r}/K_n$ with a direct argument using Jacquet modules, and
avoids character expansions and the exponential map. For $\GL_d$, Howe's
arguments already work in every characteristic without the exponential
map \cite{MR0492088}; the new points are the extension to general $\bG$
and the quantitative error term.

In a different direction, Kim, Shin and Templier
study supercuspidal representations of increasing formal degree in families
and their limit multiplicities \cite{MR4046074}.

Comparisons with more specific earlier results are made where they are
used: for cells in \cref{sec:cells}, for elliptic cells in
\cref{sec:elliptic}, for the Jacquet argument and the theory of covers in
\cref{sec:groups,sec:jacquet}, and for Sard's theorem in \cref{sec:sard}.


The paper is organized as follows.
\Cref{sec:notation} collects basic preliminaries on parabolic subgroups,
principal congruence subgroups, cells and elliptic forms, and proves
that elliptic cells occur only in supercuspidal representations.
\Cref{sec:incidence} studies the incidence maps and constructs the
\bnets{}. \Cref{sec:regular} shows that cells with a quantitatively
regular parabolic lift occur in no supercuspidal representation.
\Cref{sec:proofs} contains the precise statements and their proofs,
including the unramified and global cases. \Cref{sec:positivity}
proves positivity of the elliptic density and characterizes full
density. The appendices contain the quantitative Sard theorem, the
technical details of the Jacquet argument, and the examples.

\subsection*{AI disclosure}
In preparing this paper the authors used the AI systems Claude (Anthropic)
and ChatGPT (OpenAI), for drafting and revising text, checking
computations and searching the literature. The authors have reviewed the
resulting material and take full responsibility for the content of the
paper.

\subsection*{Acknowledgments}
The work was supported in part by a grant of access to OpenAI models
through the ChatGPT for Academic Researchers Program.
The second named author is supported by the European Union (ERC, Function
Fields, 101161909) and is the Dr.~A.~Edward Friedmann Career Development
Chair in Mathematics.


\section{Preliminaries}\label{sec:notation}

\subsection{Notation}\label{sec:notationbasic}

Throughout, $F$ is a non-archimedean local field of arbitrary
characteristic, with ring of integers $\OO$, maximal ideal $\pp$, uniformizer $\varpi$,
residue field $k=\OO/\pp$ of cardinality $q$ and normalized valuation $\val$.
Let $\bG$ be a connected reductive group over $F$,
$G=\bG(F)$, $\g=\Lie(\bG)(F)$, $\gd=\Hom_F(\g,F)$ its dual and $D=\dim\bG=\dim_F\g$.


We freely use standard terminology and results on algebraic groups over
$F$ (which is imperfect when $\chr F>0$) from
\cite{MR1102012,MR3729270,MR3362817}, on reductive group schemes over
$\OO$ from \cite{MR3362641,MR4520154}, on the Plancherel formula from
\cite{MR1989693}, and on analytic manifolds and groups over $F$ from
\cite{MR2179691}.

Denote by $\Ad$ the adjoint action of $G$ on $\g$ and by $\Ad^*$ the coadjoint action on $\gd$, so that $\Ad^*(g)\lf=\lf\circ\Ad(g^{-1})$.

For any $\lf\in\gd$ let $\tilde\Bf_\lf$ be the alternating form on $\g$ given by
\[
\tilde\Bf_\lf(X,Y)=\lf([X,Y]).
\]
The corresponding linear map $\g\to\gd$, $Y\mapsto\tilde\Bf_\lf(\cdot,Y)$, is
$Y\mapsto\ad^*(Y)\lf$, where $\ad^*(Y)\lf=-\lf\circ\ad(Y)$ is the coadjoint
action of $\g$; it is the differential at $1$ of $g\mapsto\Ad^*(g)\lf$.

Let $\Hk(G)$ be the Hecke algebra of $G$, the algebra of smooth compactly
supported distributions on $G$ under convolution. 
We fix a Haar measure $dg$ on $G$. Then $f\mapsto f\,dg$ is a linear isomorphism
$C_c^\infty(G)\to\Hk(G)$; we write $\varphi/dg\in C_c^\infty(G)$ for the preimage
of $\varphi\in\Hk(G)$.
For any compact open subgroup $J$ of $G$ and a character $\chi$ of $J$, the element
\[
e_J^\chi=\vol(J)^{-1}\chi^{-1}\one_J\,dg
\]
is an idempotent of $\Hk(G)$; it does not depend on the choice of $dg$. For simplicity $e_J=e_J^1$.

We fix once and for all a character $\psi$ of the additive group of $F$ which is trivial on $\OO$ and nontrivial on $\pp^{-1}$.

For a lattice $\Lat_0$ in a finite-dimensional $F$-vector space and a
vector $X$ let $\val_{\Lat_0}(X)=\sup\{a\in\Z:X\in\pp^a\Lat_0\}$ be
the depth of $X$, and $\norm{X}_{\Lat_0}=q^{-\val_{\Lat_0}(X)}$ its norm (the
maximum norm of its coordinates in any basis of $\Lat_0$); we drop
the subscript $\Lat_0$ when the lattice is clear from the context.
Linear maps between spaces with lattices are given the operator norms.
The level and depth parameters $n,r,h,b,a,s$ used below are integers.

For a subset $X$ of a topological space, $\partial X$ denotes its boundary.

By an \emph{analytic manifold} we mean a Hausdorff analytic manifold
over $F$, and a \emph{ball} in such a manifold of dimension $m$ is an open subset analytically
isomorphic to the unit ball $\{y\in F^m:\norm{y}\le1\}$ of $F^m$, for the maximum norm.
Only this elementary notion is used (charts given by convergent power series on
compact open balls); no rigid analytic geometry is needed.
For a morphism $f:X\to Z$ of analytic manifolds of the same dimension, the \emph{critical set} $\Crit f$ is the closed set
of points at which $f$ is not \'etale, i.e.\ at which the differential
$df_x:T_xX\to T_{f(x)}Z$ is not bijective, and
$\Crv_f=f(\Crit f)$ is the set of \emph{critical values}.

\subsection{Parabolic subgroups}\label{sec:parabolics}

Parabolic subgroups of $\bG$ are always assumed defined over $F$, and for such a $\bP$ we write $\uu_P$ for the Lie algebra of its unipotent radical.


For a cocharacter $\eta:\Gm\to\bG$ defined over $F$ we write $P_\bG(\eta)$ for the parabolic subgroup
\[
P_\bG(\eta)=\{g:\lim_{t\to0}\eta(t)g\eta(t)^{-1}\text{ exists}\}
\]
with unipotent radical
\[
U_\bG(\eta)=\{g:\lim_{t\to0}\eta(t)g\eta(t)^{-1}=1\}
\]
and Levi decomposition $P_\bG(\eta)=Z_\bG(\eta)\ltimes U_\bG(\eta)$.
Every parabolic subgroup is of this form.
The multiplication map $U_\bG(-\eta)\times P_\bG(\eta)\to\bG$ is an open immersion whose image is the big cell.

We fix a maximal $F$-split torus $\bA_0$ of $\bG$ and a minimal parabolic subgroup $\bP_0\supset\bA_0$.
A standard (resp., semistandard) parabolic subgroup is one containing $\bP_0$ (resp., $\bA_0$).
There are finitely many semistandard parabolic subgroups.
Denote by $\Pc$ (resp., $\Pcss$) the set of maximal proper standard (resp., semistandard) parabolic subgroups.
Every semistandard parabolic subgroup $\bP$ has a unique Levi subgroup $\bM_P$ containing $\bA_0$:
we can write $\bP=P_\bG(\eta)$ for a cocharacter $\eta\in X_*(\bA_0)$ which is positive on its unipotent roots and zero on its Levi roots.
Then $\bM_P=Z_\bG(\eta)$ and we write $U_P^-=U_\bG(-\eta)$ for the opposite unipotent radical; it depends only on $\bP$.
Any parabolic subgroup in $\Pcss$ is of the form $\dot w\bP\dot w^{-1}$ with $\bP\in\Pc$ and $\dot w$ in the normalizer $N_G(A_0)$ of $\bA_0$ in $G$.

We write $\Nss=\abs{\Pcss}$ for the number of maximal semistandard parabolic
subgroups; thus $\Nss=0$ if and only if $\Pc=\emptyset$, i.e., if and only if
$\bG^{\der}$ is anisotropic over $F$. We have
\[
\Nss=\sum_{\bP\in\Pc}[W_F:W_{F,M_P}],
\]
where $W_F=N_G(A_0)/Z_G(A_0)$ (resp., $W_{F,M_P}$) is the relative Weyl group of
$G$ (resp., of $M_P$). For example, for $\bG=\GL_d$, $\Nss=2^d-2$.


\subsection{Representations}\label{sec:representations}
By a representation of $G$ we tacitly mean a smooth representation on a complex vector space.

For a representation $(\pi,V)$ of $G$ and $\varphi\in\Hk(G)$ let $\pi(\varphi)=\int_G\pi(g)\,d\varphi(g)$.
For a parabolic subgroup $\bP=\bM\bU$ let $V_U=V/\langle\pi(u)v-v\rangle$ be the (unnormalized) Jacquet module.
For a compact open subgroup $J$ and a character $\chi$ of $J$, the operator $\pi(e_J^\chi)$ is the projection onto the
$\chi$-isotypic subspace $V^{J,\chi}$.

Let $\Irr(G)$ be the set of isomorphism classes of
irreducible representations of $G$, $\Gtemp\subset\Irr(G)$
the subset of unitarizable tempered ones, and $\Gsc\subset\Gtemp$ the subset of unitarizable supercuspidal ones.
Recall that every irreducible representation is admissible.

We use Harish-Chandra's Plancherel formula 
\begin{equation}\label{eq:plancherel}
(\varphi/dg)(1)=\int_{\Gtemp}\tr\pi(\varphi)\,d\mu_G(\pi),\quad\varphi\in\Hk(G),
\end{equation}
where $\mu_G$ is the Plancherel measure.
Rescaling the Haar measure $dg$ by $c$ rescales $\mu_G$ by $c^{-1}$.
If $e\in\Hk(G)$ is an idempotent, then $\tr\pi(e)=\dim\pi(e)V$, and
\[
0\le \int_A\tr\pi(e)\,d\mu_G(\pi)\le(e/dg)(1)
\]
for any measurable subset $A$ of $\Gtemp$.
In particular, for a compact open subgroup $K$ of $G$, its \emph{supercuspidal density} is by definition
\[
\scd{K}=\vol(K)\int_{\Gsc}\dim\pi^K\,d\mu_G(\pi)\in[0,1],
\]
which does not depend on $dg$. This is the proportion of the Plancherel mass of $\vol(K)e_K$ carried by
supercuspidal representations.


\subsection{Elliptic linear forms}\label{sec:ellipticforms}

Let
\[
\Par{\gd}=\bigcup_{\bP}\uu_P^\perp,
\]
the union over all proper parabolic subgroups $\bP$ of $\bG$. Thus,
\begin{equation}\label{eq:N}
\Par{\gd}=\bigcup_{\bP\in\Pc}\Ad^*(G)\uu_P^\perp.
\end{equation}
The complement
\begin{equation}\label{eq:E}
\Ell{\gd}=\gd\setminus\Par{\gd}
\end{equation}
is the \emph{elliptic} part of $\gd$: it consists of the linear forms whose restriction to $\uu_P$ is nonzero
for every proper $\bP$. For a lattice $\dLat\subset\gd$, the
\emph{elliptic density} of $\dLat$ is
\[
\elld{\dLat}=\elld[\bG]{\dLat}=\nu_{\dLat}(\dLat\cap\Ell{\gd}),
\]
where $\nu_{\dLat}$ is the probability Haar measure on $\dLat$.
Null sets in $\gd$ are taken with respect to a Haar measure.

In characteristic zero, a non-degenerate invariant bilinear form on $\g$ identifies $\gd$ with $\g$.
Under this identification, $\uu_P^\perp$ is $\Lie(\bP)(F)$, and up to a null set, $\Ell{\gd}$
is the set of regular semisimple elements whose centralizer is an elliptic maximal torus.
Indeed, the regular semisimple elements form a conull set, a regular
semisimple element of $\Lie(\bP)$ is conjugate under $\bU_P(F)$ to an
element of $\m$, and a maximal torus lies in a proper Levi subgroup if and
only if it is not elliptic.
However, in characteristic two, $\Par{\gd}$ is a null set for $\PGL_2$ even though $\PGL_2$ is split (\cref{ex:pgl2}).




\subsection{The incidence maps}\label{sec:incidencemaps}
For $\bP\in\Pc$ let $P=\bP(F)$ and let $\bFl_P=\bG/\bP$ be the (projective)
variety of parabolic subgroups conjugate to $\bP$.
Recall that the map $G\to\bFl_P(F)$ is surjective and that
parabolic $F$-subgroups which are conjugate over $\bar F$ are conjugate under $G$.
Thus $\Fl_P=\bFl_P(F)=G/P$ is the (compact, $F$-analytic) manifold of parabolic
$F$-subgroups of $\bG$ conjugate to $\bP$.
Let $\bInc_P$ be the \emph{incidence variety}, the closed subvariety of
$\bFl_P\times\Lie(\bG)^*$ consisting of the pairs $(\bQ,\lf)$ such that $\lf$
vanishes on the Lie algebra of the unipotent radical of $\bQ$, and let
$p_P:\bInc_P\to\Lie(\bG)^*$ be the projection. Via the first projection, the
incidence variety $\bInc_P$ is the vector bundle
$\bG\times^{\bP}\Lie(\bU_P)^\perp$ over $\bFl_P$, where $\bU_P$ is the unipotent
radical of $\bP$; in particular it is smooth. Its $F$-points are
\begin{equation}\label{eq:Y}
\Inc_P=\bInc_P(F)=\{(\bQ,\lf)\in\Fl_P\times\gd:\lf|_{\uu_Q}=0\},\qquad p_P(\bQ,\lf)=\lf.
\end{equation}
This is the associated bundle $G\times^P\uu_P^\perp$, i.e., the
coadjoint version of the partial Grothendieck--Springer resolution
(cf.~\cite{MR737927} for the adjoint version $G\times^P\Lie(\bP)(F)\to\g$ over $\C$,
and \cite[\S2.4]{Chi26} for the coadjoint version for Borel subgroups);
its subbundle $G\times^P\Lie(\bP)(F)^\perp$, with fibre
$\Lie(\bP)(F)^\perp\subset\uu_P^\perp$, is the cotangent bundle of $\Fl_P$
\cite[Lemma~1.4.9]{MR1433132}.
In particular, $G$ acts on $\Inc_P$ on the left by
$g\cdot(\bQ,\lf)=(g\bQ g^{-1},\Ad^*(g)\lf)$, and $p_P$ is $G$-equivariant.

$\Inc_P$ is a vector subbundle of $\Fl_P\times\gd$ of rank $D-\dim\uu_P$
over the manifold $\Fl_P$ of dimension $\dim\uu_P$.
Hence, $\Inc_P$ is an analytic manifold of dimension $D$ and $p_P$ is proper.
Thus, $p_P(\Inc_P)=\Ad^*(G)\uu_P^\perp$ is closed.
Taking the disjoint union over parabolic subgroups in $\Pc$ we obtain an analytic manifold
\[
\Inc=\bInc(F),\qquad \bInc=\bigsqcup_{\bP\in\Pc}\bInc_P,
\]
of dimension $D$ and a proper map $p:\Inc\to\gd$ with image $\Par{\gd}$ by \eqref{eq:N}.
In particular, $\Ell{\gd}$ is open. Let
\begin{equation}\label{eq:Sigma}
\Crv=\Crv_p=\bigcup_{\bP\in\Pc}p_P(\Crit p_P)
\end{equation}
be the set of critical values of $p$. Since $p$ is proper, $\Crv$ is closed, and
it is a null set by Sard's theorem (see \cref{thm:sard}).

If $\lf\in\Par{\gd}\setminus\Crv$, then $\lf=p(y)$ with $dp$ invertible at $y$,
so $\lf$ is an interior point of $\Par{\gd}$ by the inverse function theorem. Hence
\begin{equation}\label{eq:boundary}
\partial\Ell{\gd}\subset\Crv;
\end{equation}
in particular, $\partial\Ell{\gd}$ is a null set.



\begin{definition}[Peculiar groups]\label{def:peculiar}
We say that $\bG$ is \emph{\peculiar{}} if for some $\bP\in\Pc$ the incidence
map $p_P$ is everywhere critical, i.e., $\Crit p_P=\Inc_P$.
\end{definition}

\Peculiar{} groups exist only when $\chr F=2$.
We give a precise characterization in \cref{prop:peculiar}.

Let $\dG{}$ be the maximal dimension of the unipotent radical of a
maximal proper parabolic subgroup of $\bG$, with the convention
$\dG{}=1$ if there is none, and put
\begin{equation}\label{eq:tau}
\eexp=\eexp_{\bG}=
\begin{cases}
1+1/\dG{}&\text{if $\bG$ is not \peculiar{},}\\
1&\text{if $\bG$ is \peculiar{}.}
\end{cases}
\end{equation}



\subsection{Standard groups}\label{sec:standardgroups}
By a \emph{\nchart{}} of $G$ we mean an $F$-analytic isomorphism $\kappa$ from a
neighborhood of $1$ in $G$ onto a neighborhood of $0$ in $\g$ with
$\kappa(1)=0$ and $d\kappa_1=\id$; \nchart{}s exist since $G$ is an
$F$-analytic manifold ($\bG$ being smooth over $F$).
We call the inverse of a \nchart{} a \emph{\npar{}}.
We use similar terminology for smooth $F$-subgroups of $\bG$.
If $\bU$ is the unipotent radical of a parabolic $F$-subgroup of $\bG$, then $\bU$ is $F$-isomorphic, as a variety, to an affine space.
Composing such an isomorphism with a translation and a linear change of variables we
obtain a \npar{} $\xi:\Lie(\bU)(F)\to\bU(F)$, defined on all of $\Lie(\bU)(F)$, which is
an isomorphism of $F$-varieties.

Let $(\kappa,\Lat)$ be a pair consisting of a \nchart{} $\kappa$ and a lattice $\Lat$ of $\g$.
We say that $(\kappa,\Lat)$ is \emph{standard} if the image of $\kappa$ contains $\pp\Lat$,
$K_1=\kappa^{-1}(\pp\Lat)$ is a (compact open) subgroup of $G$ and the group law in $K_1$ is given by a formal group law with coefficients in $\OO$:
\[
\kappa(xy)=\Phi(\kappa(x),\kappa(y)),\ x,y\in K_1
\]
where, in the coordinates $X=(X_1,\dots,X_D)$ of a basis of $\Lat$,
the components $\Phi_1,\dots,\Phi_D\in\OO[[X,Y]]$ are formal power series
in the $2D$ variables $X,Y$. This also implies $[\Lat,\Lat]\subset\Lat$:
if the quadratic term of $\Phi(X,Y)$ is $\beta(X,Y)$, then
$[X,Y]=\beta(X,Y)-\beta(Y,X)$, and $\beta$ has integral coefficients.

For every \nchart{} $\kappa$ and lattice $\Lat$ there is $n_0\ge0$ such that $(\kappa,\pp^{n_0}\Lat)$ is standard.

Suppose that $(\kappa,\Lat)$ is standard. Then the sets
\begin{equation}\label{eq:tower}
K_n=\kappa^{-1}(\pp^n\Lat),\qquad n\ge1,
\end{equation}
are compact open subgroups of $G$ which are normal in $K_1$, and $\Ad(K_1)\Lat=\Lat$.
The groups $K_n$ are called the principal congruence subgroups with respect to $\Lat$.
For any $a,b\ge 1$, $x\in K_a$, $y\in K_b$, we have
\begin{subequations}
\begin{gather}
\kappa(xy)-\kappa(x)-\kappa(y)\in\pp^{a+b}\Lat,\qquad
\kappa(x^{-1})+\kappa(x)\in\pp^{2a}\Lat,\label{eq:law}\\
[K_a,K_b]\subset K_{a+b}.\label{eq:commutator}
\end{gather}
\end{subequations}
For $0\le s\le a$, $\kappa$ induces an isomorphism of Abelian groups
\[
K_a/K_{a+s}\rightarrow\pp^a\Lat/\pp^{a+s}\Lat.
\]
Moreover, the pushforward under $\kappa$ of a Haar measure on $K_1$ is a Haar measure on $\pp\Lat$ so that
$\vol(K_{a+s})=q^{-sD}\vol(K_a)$ for $a\ge1$, $s\ge0$.

For a fixed lattice $\Lat$, the principal congruence subgroups defined
by two standard pairs $(\kappa,\Lat)$ and $(\kappa',\Lat)$ coincide for all sufficiently large $n$.

For the rest of the paper we fix a standard pair $(\kappa,\Lat)$.
This determines the principal congruence subgroups $K_n$, the dual lattice
$\dLat=\Lat^\vee=\{\lf\in\gd:\lf(\Lat)\subset\OO\}$ and the cell characters below.

Auxiliary lattices introduced later do not change these data.

As a general rule, constants denoted $C$, $c$, \dots\ may depend on $F$, $\bG$, $(\kappa,\Lat)$ and
auxiliary choices but not on running variables and parameters $n,r,h$.


\subsection{Models}\label{sec:models}

The following result is standard  (cf.~\cite[Part~II, Chapter~IV, \S8]{MR2179691}; \cite[Definition~A.5.12, Lemma~A.5.13]{MR4520154}).

\begin{lemma}[Smooth models]\label{lem:model}
Let $\cG$ be a smooth affine group scheme of finite type over
$\OO$ with generic fiber $\bG$, and let $\Lat=\Lie(\cG)(\OO)$.
Then there is a \nchart{} $\kappa$ such that $(\kappa,\Lat)$ is standard
and $K_n=\ker(\cG(\OO)\to\cG(\OO/\pp^n))$ for all
$n\ge1$. Moreover $[\cG(\OO):K_1]=\abs{\cG(k)}$ and
$\vol(K_n)=\abs{\cG(k)}^{-1}q^{-(n-1)D}\vol(\cG(\OO))$.
\end{lemma}

For completeness we give some details.
Choose $t_1,\dots,t_D$ in the ideal $I$ of the identity
section whose classes form the basis of $I/I^2=\Lat^\vee$ dual to a basis
$X_1,\dots,X_D$ of $\Lat$, and put $\kappa=\sum_it_iX_i$. Since $\cG$ is
smooth, the completion of $\OO[\cG]$ along $I$ is $\OO[[t_1,\dots,t_D]]$, so
the group law becomes a formal group law over $\OO$, and Hensel's lemma
identifies $\ker(\cG(\OO)\to\cG(\OO/\pp^n))$ with $\kappa^{-1}(\pp^n\Lat)$
for $n\ge1$. Hence $(\kappa,\Lat)$ is standard. Hensel's lemma also
makes $\cG(\OO)\to\cG(k)$ surjective, giving the index formula, and the
volume formula follows from \cref{sec:standardgroups}.

A basic source of models is a realization of $\bG$ as an algebraic subgroup
of $\GL_d$ over $F$, that is, a homomorphism $\bG\to\GL_d$ over $F$ with trivial
scheme-theoretic kernel; such a homomorphism is a closed immersion, and we
write $\bG\subset\GL_d$. This gives rise to a lattice and principal congruence subgroups
\begin{equation}\label{eq:embeddingtower}
\Lat^\natural=\g\cap\Mat_d(\OO),\qquad
K^\natural_n=G\cap\bigl(1+\pp^n\Mat_d(\OO)\bigr)\quad(n\ge1).
\end{equation}
The groups $K^\natural_n$ are the kernels of reduction modulo $\pp^n$ on the
$\OO$-points of the schematic closure $\cG^\natural$ of $\bG$ in
$\GL_{d,\OO}$, since $\cG^\natural(\OO)=G\cap\GL_d(\OO)$.
Moreover, $\Lat^\natural=\Lie(\cG^\natural)(\OO)$. Indeed, if
$I\subset\OO[\GL_d]$ is the
ideal of $\cG^\natural$, then $X\in\Mat_d(\OO)$ lies in
$\Lie(\cG^\natural)(\OO)$ if and only if $df_1(X)=0$ for all $f\in I$, and
$I$ spans the ideal of $\bG$ over $F$. If $\cG^\natural$ is smooth,
\cref{lem:model} applies to it. In general $\cG^\natural$ need not be
smooth. Nevertheless, the results of this paper apply to $(K^\natural_n)_n$,
with the constant $\elld{(\Lat^\natural)^\vee}$, because of the following lemma.


\begin{lemma}[Towers from a linear embedding]\label{lem:embedding}
There are a \nchart{} $\kappa^\natural$ of $G$ and an integer $n_1$ such that
$K^\natural_n=(\kappa^\natural)^{-1}(\pp^n\Lat^\natural)$ for all
$n\ge n_1$. Hence, for $n_2\ge n_1$ such that
$(\kappa^\natural,\pp^{n_2}\Lat^\natural)$ is standard
(\cref{sec:standardgroups}), the tower of $\pp^{n_2}\Lat^\natural$ is
$(K^\natural_{n+n_2})_{n\ge1}$.
\end{lemma}

\begin{proof}
The lattice $\Lat^\natural$ is saturated in $\Mat_d(\OO)$, so there is
a direct summand $V_{\OO}$ with $\Mat_d(\OO)=\Lat^\natural\oplus V_{\OO}$;
let $V=V_{\OO}\otimes F$ and let $\pi:\Mat_d(F)\to\g$ be the projection
along $V$. The map $\iota(g)=g-1$ is an analytic immersion of $G$ near
$1$ with $d\iota_1$ the inclusion $\g\subset\Mat_d(F)$, so
$\kappa^\natural=\pi\circ\iota$ is a \nchart{} of $G$, and $\rho=(1-\pi)\circ\iota:G\to V$ satisfies
$\rho(1)=0$, $d\rho_1=0$. Hence there are $n_0$ and $c$ such that
$\kappa^\natural$ maps a neighborhood of $1$ bijectively onto
$\pp^{n_0}\Lat^\natural$ and
$\rho\circ(\kappa^\natural)^{-1}(X)\in\pp^{2\val(X)-c}V_{\OO}$
for $\val(X)\ge n_0$, with valuations taken with respect to
$\Lat^\natural$.

For $n\ge\max(n_0,c)$ and $g$ near $1$, $g\in K^\natural_n$ if and only if
$\iota(g)\in\pp^n(\Lat^\natural\oplus V_{\OO})$, i.e., if and only if
$\kappa^\natural(g)\in\pp^n\Lat^\natural$ and $\rho(g)\in\pp^nV_{\OO}$;
and the first condition implies the second, since then
$\rho(g)\in\pp^{2n-c}V_{\OO}$. It remains to note that $K^\natural_n$ is
contained in any given neighborhood of $1$ for $n$ large.
\end{proof}

\begin{remark}[Parahoric towers]\label{rem:badplaces}
Let $x\in\Bldg(\bG)$ and let $\cG_x$ be the parahoric group scheme of
Bruhat--Tits (with connected special fibre). Its congruence kernels
$\ker(\cG_x(\OO)\to\cG_x(\OO/\pp^n))$, $n\ge1$, are the Moy--Prasad groups
$G_{x,n}$ of integral depth $n$ (not $G_{x,n+}$), for the
valuation normalized by $\val(F^\times)=\Z$ \cite[Proposition~8.5.16]{MR4520154};
they form the tower of the smooth model $\cG_x$ (\cref{lem:model}). If $(K_n)$ and
$(K'_n)$ are two towers \eqref{eq:tower}, for instance this one and that of
\cref{lem:embedding}, then their lattices are commensurable and their charts
differ by $X\mapsto X+O(X^2)$; consequently $K'_{n+c}\subset K_n\subset
K'_{n-c}$ for a fixed $c$ and all sufficiently large $n$. These inclusions compare levels, while the limiting
densities retain their dependence on the lattices. When $x$ is hyperspecial,
\cref{thm:uniform} applies to the tower of $\cG_x$.
\end{remark}

\subsection{Unramified groups}\label{sec:integraldata}\footnote{This subsection is needed only for the explicit uniform bounds
(\cref{lem:integralatlas,thm:regularunram,thm:uniform}).}
Following Borel, we say that $\bG$ is \emph{unramified} if
it is quasi-split and splits over an unramified extension of $F$.
Equivalently, $\bG$ has a \emph{reductive model}, that is, a reductive group
scheme $\cG$ over $\OO$ (smooth and affine, with connected reductive geometric
fibres) with generic fibre $\bG$; the groups $\cG(\OO)$ are then exactly the
hyperspecial parahoric subgroups of $G$. For these facts, and for the
structure theory of reductive group schemes over $\OO$ recalled below, we
refer to \cite{MR3362641} and \cite{MR4520154}.


By the \emph{unramified case} we mean the following situation: $\bG$ is
unramified, $\cG$ is a reductive model of $\bG$, $\Lat=\Lie(\cG)(\OO)$, and
$\kappa=\sum_it_iX_i$ is the chart of \cref{lem:model}, so that
$(\kappa,\Lat)$ is standard and
\begin{equation}\label{eq:integralKn}
K_n=\kappa^{-1}(\pp^n\Lat)=\ker\bigl(\cG(\OO)\to\cG(\OO/\pp^n)\bigr),\qquad n\ge1.
\end{equation}
In the unramified case, all norms are those induced by $\Lat$, its
direct summands below, and their duals.
The model contains closed $\OO$-tori $\cS\subset\cT$ with $\cS$
split and maximal among split tori, and $\cT$ a maximal torus which splits
over the maximal unramified extension.

A Borel subgroup of $\bG$ containing $T=\cT_F$
extends uniquely to a Borel subgroup scheme $\cB\supset\cT$ of $\cG$, and
$B=\cB_F$ is a minimal parabolic subgroup. Let $\Delta$ be the set of
simple roots of $\bG_{\bar F}$ with respect to $(\cT,\cB)$ and $W$ the
absolute Weyl group; the Galois group of $F$ acts on $\Delta$ through
the Galois group of $k$, and the group $X_*(\cT)$ of cocharacters of the
$\OO$-torus $\cT$, i.e., of $T$ over $F$, is $X_*(T_{\bar F})^{\Gal(\bar F/F)}$.
Every parabolic subgroup $\bP$ of $\bG$ is $G$-conjugate to a standard
one, $\bP\supset B$; we take the representatives $\bP\in\Pc$ standard, i.e., $\bA_0=\cS_F$ and $\bP_0=B$ in \cref{sec:parabolics}.
The type of a standard maximal parabolic subgroup is a Galois orbit
$I\subset\Delta$, the simple roots not in its Levi subgroup. Let $e_I$ be the smallest positive integer for which there is a
Galois-invariant cocharacter $\eta_I\in X_*(T_{\bar F})$, determined up to a central cocharacter, satisfying
\begin{equation}\label{eq:cocharchoice}
\langle\alpha,\eta_I\rangle=
\begin{cases}e_I&\alpha\in I,\\0&\alpha\in\Delta\setminus I.\end{cases}
\end{equation}
Fix one such $\eta_I$.
It extends to a cocharacter of $\cT$ and satisfies $\bP=P_{\bG}(\eta_I)$.
The closed subgroup schemes
$\cP=P_{\cG}(\eta_I)$, $\cM=Z_{\cG}(\eta_I)$ and $\cU^\pm=U_{\cG}(\pm\eta_I)$
are smooth with generic fibres $\bP$, $\bM$, $\bU^\pm$, $\cP=\cM\ltimes\cU$,
and the multiplication map
\begin{equation}\label{eq:integralcell}
\cU\times\cM\times\cU^-\to\cG
\end{equation}
is an open immersion. Moreover the action of $\Gm$ on $\cU$ by
conjugation through $\eta_I$ extends to an action of the multiplicative
monoid $\Aff^1$, so that conjugation by $t=\eta_I(\varpi)$ is an endomorphism of the
$\OO$-scheme $\cU$; similarly for $t^{-1}$ on $\cU^-$.

The quotient
$\cG/\cP$ is a smooth projective $\OO$-scheme, $\Fl_P=\bFl_P(F)=(\cG/\cP)(\OO)$,
and the big cell $\cU^-\cP/\cP\simeq\cU^-$ is an open subscheme, and the Iwasawa decomposition $G=P\,\cG(\OO)$ holds.

Let $W_k=N_{\cG}(\cS)(k)/Z_{\cG}(\cS)(k)$ be the relative Weyl group of $\cG_k$ and $W_{k,P}$ that of $\cM_k$.
We may lift the Bruhat decomposition of $\cG_k$ to $\OO$.
More precisely, first, we may lift any $w\in W_k/W_{k,P}$ to $\dot w\in N_{\cG}(\cS)(\OO)$.
Put $\bP_w=\dot w\bP\dot w^{-1}$, and let $\cP_w$ and $\cU^\pm_w$ be the conjugates of $\cP$ and $\cU^\pm$ by $\dot w$. Then
\begin{equation}\label{eq:integralcover}
(\cG/\cP)(\OO)=\bigcup_{w\in W_k/W_{k,P}}\dot w\,\cU^-(\OO)\cP(\OO)/\cP(\OO);
\end{equation}
that is, every parabolic subgroup conjugate to $\bP$ is $u\bP_wu^{-1}$ for
some $w$ and some $u\in\cU^-_w(\OO)$. Moreover, the $\bP_w$,
$w\in W_k/W_{k,P}$, are pairwise distinct elements of $\Pcss$.

Finally, $\cU^\pm$ is isomorphic as an $\OO$-scheme to affine space.
We fix such isomorphisms $\xi_\pm:\Aff^{d_P}\to\cU^\pm$ such that $\xi_\pm(0)=1$ and $d\xi_\pm$ is the identity with
respect to $\Lat_\pm=\Lie(\cU^\pm)(\OO)$; the lattices $\Lat_\pm$
are direct summands of $\Lat=\Lie(\cG)(\OO)$ and are the direct sums
of their intersections with the weight spaces of $\eta_I$.

In the coordinates $\kappa$, the group law, the inverse, the conjugation by
any element of $\cG(\OO)$, the maps \eqref{eq:integralcell} and its inverse,
and the two contracting conjugations, by $t$ on $\cU$ and by $t^{-1}$ on
$\cU^-$, are given by power series with coefficients in $\OO$: they are
morphisms of $\OO$-schemes between neighborhoods of the identity, computed in
the formal completions along the identity sections.

Define
\begin{equation}\label{eq:mG}
m_{\cG}=\max_{\bP\in\Pc}\ \max_{\langle\alpha,\eta_I\rangle>0}\langle\alpha,\eta_I\rangle,
\end{equation}
the largest weight of $\eta_I$ on $\uu_P$, the inner maximum being over
the roots $\alpha$ of $\bG_{\bar F}$ and $I$ the type of $\bP$. If $\Pc=\emptyset$, then $\bG$, being quasi-split, is a torus, and we set $m_{\cG}=0$.

For instance, for $\bG=\GL_d$, $d\ge2$, one has $m_{\cG}=1$, using $\eta_j(t)=\diag(t\one_j,\one_{d-j})$,
while for $\bG=\SL_d$ one has $m_{\cG}=d$: the cocharacter attached to the Levi subgroup $\mathrm S(\GL_j\times\GL_{d-j})$
is $t\mapsto(t^{(d-j)/g}\one_j,t^{-j/g}\one_{d-j})$ with $g=\gcd(j,d)$, of weight $d/g$ on $\uu_P$.

The invariants $\Nss$ and $m_{\cG}$ depend only on the root datum of $\bG$ with its Galois action.
They control the constants in the estimates of
\cref{lem:integralatlas}, \cref{thm:regularunram}, and \cref{thm:uniform}.

\subsection{Cells}\label{sec:cells}
Let $n\ge1$. For $\lf\in \dLat$, let
\begin{equation}\label{eq:chi}
\chi_{n,\lf}(g)=\psi\bigl(\varpi^{-n}\lf(\kappa(g))\bigr),\qquad g\in K_{\lceil n/2\rceil}.\footnote{Replacing $\varpi$
by $u\varpi$ with $u\in\OO^\times$ replaces $\chi_{n,\lf}$ by $\chi_{n,u^{-n}\lf}$.}
\end{equation}
This is a character of $K_{\lceil n/2\rceil}$, trivial on $K_n$. For any $k\in K_1$ we have
\begin{equation}\label{eq:K1action}
\chi_{n,\Ad^*(k)\lf}(x)=\chi_{n,\lf}(k^{-1}xk),\qquad
x\in K_{\lceil n/2\rceil}.
\end{equation}
Let $r\ge0$ with $n\ge2r$. The restriction of $\chi_{n,\lf}$ to
$K_{n-r}\subset K_{\lceil n/2\rceil}$, which we still denote by
$\chi_{n,\lf}$, depends only on $\cell=\lf+\pp^r\dLat$, and $\cell\mapsto\chi_{n,\lf}$ is a bijection between
$\dLat/\pp^r\dLat$ and the character group of $K_{n-r}/K_n$.

For $\GL_d$, these are Howe's dual blobs \cite{MR0492088}: his
shallow characters of $K_v=1+\pp^v\Mat_d(\OO)$ are those trivial on
$K_{2v}$, and their dual blobs are cosets of $\pp^{-v}\Mat_d(\OO)$ in
$\pp^{-2v}\Mat_d(\OO)$ \cite[\S2]{MR0492088}; up to the scaling in
\eqref{eq:chi}, these are our cells. Moy and Prasad \cite{MR1253198}
developed the corresponding filtration and dual-coset framework for general
reductive groups. At positive depth their unrefined minimal $K$-types use
nondegenerate characters of the quotients $G_{x,r}/G_{x,r+}$. We use
characters of the larger abelian congruence quotients $K_{n-r}/K_n$ and
impose no minimality condition on individual cells
(cf.~\cref{rem:badplaces}).

Set
\[
e_{n,\cell}=e_{K_{n-r}}^{\chi_{n,\lf}}.
\]
The following formula shows that this definition is compatible with refinement of the partition into cells:
\begin{equation}\label{eq:cellformula}
\frac{e_{n,\cell}}{dg}(g)=\frac{\one_{K_1}(g)}{\vol(K_n)}\int_\cell
\psi\bigl(-\varpi^{-n}\lf(\kappa(g))\bigr)\,d\nu_{\dLat}(\lf),\qquad g\in G.
\end{equation}
More generally, for a union $\cell\subset\dLat$ of cosets of $\pp^r\dLat$, let
$e_{n,\cell}=\sum_{\cell'}e_{n,\cell'}$, the sum over the cosets $\cell'$ of $\pp^r\dLat$
contained in $\cell$. Thus \eqref{eq:cellformula} continues to hold and $e_{n,\dLat}=e_{K_n}$.

By Fourier inversion on the finite abelian group $K_{n-r}/K_n$,
transported to $\Hk(G)$ by \cite[\S4.4]{MR2234120},
\begin{equation}\label{eq:partition}
e_{K_n}=\sum_{\cell\in \dLat/\pp^r\dLat}e_{n,\cell},
\end{equation}
and $e_{n,\cell}$, $\cell\in\dLat/\pp^r\dLat$, are mutually orthogonal idempotents.
By the Plancherel formula, for every union $\cell$ of cosets of $\pp^r\dLat$
\begin{equation}\label{eq:cellmass}
\vol(K_n)\int_{\Gtemp}\tr\pi(e_{n,\cell})\,d\mu_G(\pi)=\nu_{\dLat}(\cell).
\end{equation}

For $n,r$ as above and a cell $\cell\in \dLat/\pp^r\dLat$ let
\begin{equation}\label{eq:f}
f_{n,r}(\cell)=q^{rD}\vol(K_n)\int_{\Gsc}\tr\pi(e_{n,\cell})\,d\mu_G(\pi)
\end{equation}
be the supercuspidal proportion of the Plancherel mass of $\cell$, whose
total mass is $q^{-rD}=\nu_{\dLat}(\cell)$ by \eqref{eq:cellmass}; we regard $f_{n,r}$
as a function on $\dLat$ constant on cells. Since $\tr\pi(e_{n,\cell})\ge0$,
\eqref{eq:cellmass} shows that $f_{n,r}$ takes values in $[0,1]$, and
\eqref{eq:partition} gives
\begin{equation}\label{eq:Pscdecomp}
\scd{K_n}=\sum_{\cell\in \dLat/\pp^r\dLat}\nu_{\dLat}(\cell)f_{n,r}(\cell)=\int_\dLat f_{n,r}\,d\nu_{\dLat}.
\end{equation}


\subsection{Elliptic cells}\label{sec:elliptic}

For $\lf\in \dLat\cap\Ell{\gd}$ let $\erad(\lf)$ be the minimal integer $R\ge0$ with
$\lf+\pp^R\dLat\subset\Ell{\gd}$, the \emph{ellipticity radius} of $\lf$; it is
finite since $\Ell{\gd}$ is open (\cref{sec:incidencemaps}), and
$\erad(\Ad^*(k)\lf)=\erad(\lf)$ for $k\in K_1$.

The following is a quantitative analogue of Moy--Prasad's criterion
for anisotropic unrefined minimal $K$-types \cite[Corollary~4.9]{MR1371680}.
It is closely related to \cite{MR3463423} and \cite[Appendix~D]{MR4325284}, which work in
characteristic zero and use the exponential map; for $\SL_2$ see also
\cite[\S6.1]{MR3753921}.

\begin{proposition}[Elliptic cells]\label{prop:elliptic}
There is a constant $C$ such that for every $\lf\in \dLat\cap\Ell{\gd}$, proper parabolic subgroup $P$ and integers $r\ge\erad(\lf)$ and $n\ge 2r+C$ we have 
\begin{equation}\label{eq:nontrivialonU}
\chi_{n,\lf}|_{K_{n-r}\cap U_P}\ne1.
\end{equation}
Hence, every irreducible representation of $G$ which contains the character $\chi_{n,\lf}$ of $K_{n-r}$ is
supercuspidal. In particular, $f_{n,r}(\lf)=1$.
In the unramified case (\cref{sec:integraldata}) one can take $C=0$.
\end{proposition}

\begin{proof}
There are only finitely many $K_1$-conjugacy classes of proper parabolic subgroups.
Therefore, by \eqref{eq:K1action} it is enough to show \eqref{eq:nontrivialonU} for a fixed $P$
and all $\lf\in \dLat\cap\Ell{\gd}$ with $r\ge\erad(\lf)$ and $n\ge 2r+C$, with $C$ depending on $P$.

Let $\uu=\uu_P$, $\Lat_{\uu}=\Lat\cap\uu$ and
let $\rho:\gd\to\uu^*$ be the restriction. Since $\dLat=\Lat^\vee$, we have
$\rho(\dLat)=\Lat_{\uu}^\vee$. Fix a \npar{} $\xi$ of $U_P$. Since $\kappa\circ\xi(Y)=Y+O(Y^2)$, there is $c$ such that
$\kappa(\xi(Y))-Y\in\pp^{2a-c}\Lat$ for $Y\in\pp^a\Lat_{\uu}$,
$a\ge c$. In particular $\xi(\pp^b\Lat_{\uu})\subset K_b\cap U_P$,
$b=n-r$, and if $2b-c\ge n$, then
\[
\chi_{n,\lf}(\xi(Y))=\psi\bigl(\varpi^{-n}\lf(Y)\bigr),\qquad Y\in\pp^b\Lat_{\uu}.
\]
This character of $\pp^b\Lat_{\uu}$ is trivial if and only if
$\varpi^{-r}\rho(\lf)\in\Lat_{\uu}^\vee=\rho(\dLat)$, i.e., if and only if
$\lf+\pp^r\dLat$ meets $\uu^\perp\subset\Par{\gd}$, which is impossible for
$r\ge\erad(\lf)$. This proves \eqref{eq:nontrivialonU}.

Let $\Pi=\cInd_{K_b}^G\chi_{n,\lf}$ and $\chi=\chi_{n,\lf}$. For a proper
parabolic subgroup $P=MU$, the Mackey formula gives
\[
\Pi|_P\simeq\bigoplus_{g\in K_b\backslash G/P}\cInd_{P\cap g^{-1}K_bg}^P\chi^g,
\qquad\chi^g(x)=\chi(gxg^{-1}).
\]
The $U$-coinvariants of $\cInd_{P\cap g^{-1}K_bg}^P\chi^g$ vanish as soon as
$\chi^g$ is nontrivial on $U\cap g^{-1}K_bg$, i.e., as soon as $\chi$ is
nontrivial on $K_b\cap gUg^{-1}$. By \eqref{eq:nontrivialonU}, applied to the
parabolic subgroup $gPg^{-1}$, this holds for every $g$, so $\Pi_U=0$.
Thus every proper Jacquet module of $\Pi$ is zero. On the other hand, by Frobenius
reciprocity, an irreducible representation $\pi$ containing $\chi_{n,\lf}$
is a quotient of $\Pi$. Therefore, $\pi$ is supercuspidal by the exactness of the Jacquet functor.
The assertion about $f_{n,r}$ follows from \eqref{eq:f} and \eqref{eq:cellmass}.

In the unramified case, put $K=\cG(\OO)$. By the Iwasawa decomposition, $K$
acts transitively on each conjugacy class of parabolic subgroups, and
conjugation by $K$ preserves $\Lat$, $\dLat$, every $K_b$ and, in the integral
coordinates, the cell characters when $2b\ge n$; so \eqref{eq:K1action} holds
for $k\in K$, and it suffices to treat a standard $\bP$. The lattice
$\Lat_\uu=\Lie(\cU)(\OO)$ is a direct summand of $\Lat$, and the integral
parametrization $\xi_+$ satisfies $\kappa(\xi_+(Y))-Y\in\pp^{2a}\Lat$ for
$Y\in\pp^a\Lat_\uu$, $a\ge1$. Thus the argument above works with $c=0$, and
$n\ge2r$ suffices.
\end{proof}

\begin{remark}[Existence of an irreducible representation containing a type]\label{rem:Ktype}
For a compact open subgroup $K$ of $G$ and an irreducible representation
$\rho$ of $K$ there is an irreducible representation $\pi$ of $G$
containing $\rho$. Indeed, $\cInd_K^G\rho$ is finitely generated -- it is generated by
the functions supported on $K$ -- so it has an irreducible quotient $\pi$, and $\Hom_K(\rho,\pi)=\Hom_G(\cInd_K^G\rho,\pi)\ne0$.
In particular, the representations in \cref{prop:elliptic} exist.
\end{remark}


\section{Parabolic incidence maps}\label{sec:incidence}

Using the quantitative Sard theorem of \cref{sec:sard} (\cref{thm:sard}), we
construct decreasing \bnets{} $\Bad_h\subset\dLat$, each a union of cosets of
$\pp^{2h}\dLat$, with $\nu_{\dLat}(\Bad_h)\ll q^{-\eexp h}$, such that every coset of
$\pp^{2h}\dLat$ outside $\Bad_h$ consists entirely of elliptic forms or
entirely of nonelliptic forms. In the latter case, the incidence map has a
lift with quantitatively controlled inverse differential. This is the
geometric input for the cellwise comparison of \cref{prop:density}.


\subsection{The pairing and the geometric alternatives}\label{sec:levi}
Fix a semistandard parabolic subgroup $\bP=\bM\bU$ and its opposite
semistandard $\bP^-=\bM\bU^-$ with common semistandard Levi subgroup $\bM$. Write
$\g=\uu\oplus\m\oplus\uu^-$, with $\uu^+=\uu$ and $\bU^+=\bU$,
and put $d_P=\dim\uu=\dim\uu^-$.
For $\lf\in\gd$ define
\begin{equation}\label{eq:levipairing}
\Bf_\lf(X,Y)=\lf([X,Y]),\qquad X\in\uu,\ Y\in\uu^-.
\end{equation}
For $\lf\in\uu^\perp$, let $Q_{\lf,P}:\uu^-\to\uu^*$ and $Q_{\lf,P}^*:\uu\to(\uu^-)^*$ be the corresponding maps.
Suppressing $P$ when it is understood from the context, we have
\begin{equation}\label{eq:pairing}
(Q_\lf Y)(X)=(Q_\lf^*X)(Y)=\Bf_\lf(X,Y).
\end{equation}
Identify $\m^*$ with the forms vanishing on $\uu\oplus\uu^-$, and let
$\Omega_M\subset\m^*$ be the open, possibly empty, set on which $Q_\lf$ is invertible
(it is Zariski open, hence also open for the topology of $F$).


For each $\bP\in\Pcss$, endow $\operatorname{Hom}_F(\uu^*,\uu^-)$ with a norm once and for all.
The choice affects only the constant $c$ in the following proposition.

\begin{proposition}[Geometric alternatives outside the \bnets{}]\label{prop:geometric}
There are a compact subset $\mathcal C\subset G$, compact sets
$\mathcal A_P\subset\uu_P^\perp$ for $\bP\in\Pcss$, an integer $c\ge0$,
and decreasing sets $\Bad_h\subset\dLat$, $h\ge1$ satisfying the following conditions.
\begin{enumerate}[label=\textup{(\arabic*)},ref=\arabic*]
\item\label{item:geo-cosets} Each $\Bad_h$ is a union of cosets of $\pp^{2h}\dLat$.
\item\label{item:geo-measure} $\bigcap_{h\ge1}\Bad_h=\Crv\cap\dLat$ and
$\nu_{\dLat}(\Bad_h)\ll q^{-\eexp h}$.
\item\label{item:geo-alt} For every $\lf\in\dLat\setminus\Bad_h$ either $\lf+\pp^{2h}\dLat\subset\Ell{\gd}$, or
$\lf+\pp^{2h}\dLat\subset\Par{\gd}$, in which case there are
$\bP\in\Pcss$, $g\in\mathcal C$ and $\nf\in\mathcal A_P$ with
\[
\lf=\Ad^*(g)\nf,\qquad \norm{Q_{\nf,P}^{-1}}\le q^{h+c}.
\]
\item\label{item:geo-tube} $\nu_{\dLat}((\partial\Ell{\gd}+\pp^r\dLat)\cap\dLat)\le\nu_{\dLat}((\Crv+\pp^r\dLat)\cap\dLat)\ll q^{-\eexp r/2}$
for $r\ge2$.
\end{enumerate}
\end{proposition}

In the unramified case we will make the statement more precise in \cref{prop:geometricunram}.


\subsection{The differential and the coordinate Jacobian}\label{sec:incidencejacobian}
For each $\bP\in\Pcss$, fix a normalized parametrization
$\xi_-:\uu^-\to U^-$ which is an isomorphism of $F$-varieties (\cref{sec:standardgroups}).
For a semistandard conjugate $\bP$, we identify $\Inc_P$ with the
corresponding component of $\Inc$. Let $\Incc{P}$ be the inverse image
in $\Inc_P$ of the big cell
$U^-P/P\subset\Fl_P$. The trivialization
\begin{equation}\label{eq:trivialization}
U^-\times\uu^\perp\longrightarrow\Incc{P},\qquad
(u,\nf)\longmapsto u\cdot(\bP,\nf)
\end{equation}
and $\xi_-$ give a polynomial parametrization
\begin{equation}\label{eq:bigcellchart}
\phi_P:\uu^-\times\uu^\perp\longrightarrow\Incc{P},\qquad
(y,\nf)\longmapsto\xi_-(y)\cdot(\bP,\nf).
\end{equation}
The incidence map in these coordinates is
\[
f_P=p\circ\phi_P:\uu^-\times\uu^\perp\longrightarrow\gd,\qquad
(y,\nf)\longmapsto\Ad^*(\xi_-(y))\nf.
\]
Fix bases of $\uu^-$, $\uu^\perp$ and $\uu^*$. Order the basis of
$\uu^-\oplus\uu^\perp$ with the basis of $\uu^\perp$ first, and use for $\gd$
the basis of $\uu^\perp$ followed by lifts of the basis of $\uu^*$ (via the
exact sequence $0\to\uu^\perp\to\gd\to\uu^*\to0$). A linear map
$\uu^-\oplus\uu^\perp\to\gd$ which is the identity on $\uu^\perp$ and induces
$Q:\uu^-\to\uu^*$ on the quotients then has determinant $\det Q$ in these
bases.\footnote{In characteristic zero, a
non-degenerate invariant symmetric bilinear form on $\g$ identifies $\gd$
with $\g$ and $\uu^*$ with $\uu^-$, so that $Q_\nf$ becomes an endomorphism
of $\uu^-$ and $\det Q_\nf$ is defined without further choices. In positive
characteristic such a form need not exist.} After fixing bases,
$\nf\mapsto\det Q_\nf$ is a polynomial, well defined up to a nonzero constant.

\begin{lemma}[Differential and Jacobian of the incidence map]\label{lem:levi}
In the trivialization \eqref{eq:trivialization},
\begin{equation}\label{eq:dp}
(dp_P)_{(\bP,\nf)}(Y,\mu)=\mu+\ad^*(Y)\nf,
\qquad Y\in\uu^-,\ \mu\in\uu^\perp.
\end{equation}
Moreover,
\begin{equation}\label{eq:jacobianfiber}
\det Q_\nf=\det Q_{\nf|_\m},
\end{equation}
and hence
\begin{equation}\label{eq:levi}
(dp_P)_{(\bP,\nf)}\text{ is invertible}
\iff Q_\nf\text{ is invertible}\iff\nf|_\m\in\Omega_M.
\end{equation}
Finally, in the bases above,
\begin{equation}\label{eq:detfield}
\det df_P(y,\nf)=\det Q_\nf\qquad\text{for all }(y,\nf)\in\uu^-\times\uu^\perp.
\end{equation}
\end{lemma}

\begin{proof}
Differentiating the coadjoint action gives \eqref{eq:dp}. The map is
the identity on $\uu^\perp$ and induces $Q_\nf$ on the quotients.
For \eqref{eq:jacobianfiber}, let $\g=\bigoplus_k\g_k$ be the weight
decomposition for a cocharacter $\eta\in X_*(\bA_0)$ defining $\bP$, so that
$\uu=\bigoplus_{k>0}\g_k$, $\m=\g_0$ and $\uu^-=\bigoplus_{k>0}\g_{-k}$.
If $X\in\g_i$, $Y\in\g_{-j}$ with $i>j>0$, then $[X,Y]\in\g_{i-j}\subset\uu$,
so $\Bf_\nf(X,Y)=0$. Thus $Q_\nf$ is block triangular with respect to the
weight decompositions, and its diagonal blocks $\g_{-i}\to\g_i^*$ only
involve $\nf$ on $[\g_i,\g_{-i}]\subset\m$. Hence
$\det Q_\nf=\prod_{i>0}\det(Q_\nf|_{\g_{-i}\to\g_i^*})=\det Q_{\nf|_\m}$,
proving \eqref{eq:jacobianfiber} and hence \eqref{eq:levi}. The same
argument shows that $Q_\nf$ is invertible if and only if each pairing
$\g_i\times\g_{-i}\to F$, $(X,Y)\mapsto\nf([X,Y])$, is non-degenerate.

For the last assertion fix $y$, put $u=\xi_-(y)$. Writing multiplication by $u$ in the $\uu^-$ coordinates, we obtain a polynomial automorphism
$\theta_y$ of $\uu^-$ such that $\xi_-(y')=u\xi_-(\theta_y(y'))$ and in particular $\theta_y(y)=0$. Thus
\[
f_P(y',\nf')=\Ad^*(u)f_P(\theta_y(y'),\nf').
\]
Keeping $y$ and $u$ fixed and differentiating at $(y,\nf)$ gives by the chain rule
\[
df_P(y,\nf)=\Ad^*(u)\circ df_P(0,\nf)\circ(T_y\oplus1)
\]
where $T_y=d(\theta_y)_y$.
Since $\theta_y$ is a polynomial automorphism for every geometric point $y$, $\det T_y$ is a polynomial in $y$ which has no zeros over $\bar F$
and it equals $1$ at $y=0$. Hence $\det T_y\equiv1$.
Since $\bU^-$ is unipotent, $\det\Ad^*(u)=1$. The quotient calculation
above shows that $\det df_P(0,\nf)=\det Q_\nf$ in the bases
preceding the lemma, since $d\xi_-$ is the identity at $0$. This proves \eqref{eq:detfield}.
\end{proof}

In particular, $\det Q_\nf$ is a homogeneous polynomial of degree
$d_P$ in the fiber coordinate. It is nonzero if and only if
$\Omega_M\ne\emptyset$, equivalently if $p_P$ is not everywhere critical.


\subsection{Construction of the \bnets{}}\label{sec:bad}
\begin{proof}[Proof of \cref{prop:geometric}]
If $\Pc=\emptyset$, take $\Bad_h=\emptyset$; all forms are elliptic.
Otherwise the big cells of the semistandard parabolic subgroups
cover $\Inc$. Indeed, let $\bP\in\Pc$. By the Bruhat decomposition
$G=\bigcup_{\dot w}P_0\dot wP$, together with
$\dot w^{-1}P_0\dot w\subset U_P^-P$, we have $G=\bigcup_{\dot w}\dot wU_P^-P$,
$\dot w\in N_G(A_0)$. Hence the translates $\dot wU_P^-P/P$ cover
$\Fl_P$, and $\dot wU_P^-P/P$ is the big cell of the semistandard
conjugate $\dot w\bP\dot w^{-1}$.

Let $x\in\Inc_\dLat=p^{-1}(\dLat)$, a compact open set. Choose a
semistandard $\bP_x$ whose big cell contains $x$, write
$x=\phi_{P_x}(y_x,\nf_x^0)$, and fix lattices $H^-\subset\uu_{P_x}^-$ and
$H^0\subset\uu_{P_x}^\perp$ with bases. The map
$(z,\mu)\mapsto f_{P_x}(y_x+z,\nf_x^0+\mu)$ is polynomial with constant term
$p(x)\in\dLat$. Replacing $(H^-,H^0)$ by $(\pp^sH^-,\pp^sH^0)$, with the
rescaled bases, multiplies the coefficient of each monomial of degree $k$
by $\varpi^{sk}$. As there are finitely many coefficients, for $s$ large
all coefficients of positive degree lie in $\dLat$. Then
$(z,\mu)\mapsto\phi_{P_x}(y_x+z,\nf_x^0+\mu)$ maps $\pp^sH^-\oplus\pp^sH^0$
into $\Inc_\dLat$, and its image is an open neighborhood of $x$, because
$\phi_{P_x}$ is an open immersion. By compactness, finitely many of these
charts cover $\Inc_\dLat$; we write them as
\[
\psi_i:\Lat_i\longrightarrow\Inc_\dLat,\qquad
\psi_i(z,\mu)=\phi_{P_i}(y_i+z,\nf_i^0+\mu),
\]
where $\Lat_i=H_i^-\oplus H_i^0$ is a product of lattices in
$\uu_{P_i}^-\oplus\uu_{P_i}^\perp$, and the polynomial maps
$f_i=p\circ\psi_i:\Lat_i\to\dLat$ have coefficients in $\dLat$.
Thus the $\psi_i$ form an \iatlas{} in the sense of \cref{sec:bnets}.
Computing $\det df_i$ in the bases of $\Lat_i$ and $\dLat$, and $\det Q$ in
fixed bases of $\uu_{P_i}^-$ and $\uu_{P_i}^*$, \eqref{eq:detfield} gives
constants $a_i\in F^\times$ such that
\begin{equation}\label{eq:atlasjacobian}
\det df_i(z,\mu)=a_i\det Q_{\nf_i^0+\mu,P_i}.
\end{equation}


Let $\Bad_h$ be the \bnets{} \eqref{eq:badf} for this atlas.
By \cref{thm:sard} they decrease to $\Crv\cap\dLat$ and are unions of
cosets of $\pp^{2h}\dLat$. Since $p(\Inc)=\Par{\gd}$, the same theorem gives
\begin{equation}\label{eq:dichotomy}
\lf+\pp^{2h}\dLat\not\subset\Bad_h\implies
\text{either }\lf+\pp^{2h}\dLat\subset\Ell{\gd}
\text{ or }\lf+\pp^{2h}\dLat\subset\Par{\gd}.
\end{equation}
The construction excludes neighborhoods of all critical values,
including those which also have noncritical lifts.

In the parabolic case choose any $x\in p^{-1}(\lf)$ and a chart
containing $x$, say $x=\psi_i(z,\mu)$. Put
$\nf=\nf_i^0+\mu$ and $g=\xi_-(y_i+z)$. Then
$\lf=\Ad^*(g)\nf$, with $g$ and $\nf$ in the fixed compact sets
\[
\mathcal C=\bigcup_i\xi_-(y_i+H_i^-),\qquad
\mathcal A_P=\bigcup_{i:P_i=P}(\nf_i^0+H_i^0).
\]
By \eqref{eq:sardchart} and \eqref{eq:atlasjacobian},
$\abs{\det Q_\nf}>\abs{a_i}^{-1}q^{-h}$.
The entries of $Q_\nf$ are bounded on these compact sets, so Cramer's
rule gives $\norm{Q_\nf^{-1}}\le q^{h+c}$, with $c$ depending on
the fixed reference norms and compact families.

It remains to estimate the measure.
By \eqref{eq:sardmeasure} $\nu_{\dLat}(\Bad_h)\le Nq^{-h}$.
If $\Omega_M\ne\emptyset$ for every $\bP\in\Pc$, the same holds for
their semistandard conjugates. Thus \eqref{eq:atlasjacobian} is a nonzero
polynomial of degree $d_{P_i}$ on each source lattice. The polynomial
sublevel estimate of \cref{lem:sublevel} gives
\[
\vol(S_h(f_i))\ll q^{-h/d_{P_i}}\le q^{-h/\dG{}},
\qquad
\nu_{\dLat}(\Bad_h)\le q^{-h}\sum_i\vol(S_h(f_i))
\ll q^{-(1+1/\dG{})h}.
\]
When $\bG$ is not \peculiar{}, every $\Omega_M$, $\bP\in\Pc$, is nonempty by
\eqref{eq:levi} and \cref{def:peculiar}, so the improved bound applies.
Otherwise the bound $O(q^{-h})$ suffices.
These are precisely the exponents in \eqref{eq:tau}.

Finally, \cref{item:geo-tube} follows from \cref{item:geo-measure} and
\eqref{eq:bntube}, since $\partial\Ell{\gd}\subset\Crv$ by
\eqref{eq:boundary}.
\end{proof}

\subsection{The unramified case}\label{sec:integralatlas}
In the unramified case (\cref{sec:integraldata}) the charts can be taken
integral, and their number is bounded in terms of the root datum. Let
$\bP\in\Pc$, let $\dot w$, $\bP_w$, $\cP_w$ and $\cU^\pm_w$,
$w\in W_k/W_{k,P}$, be as in \cref{sec:integraldata}, let $\xi^w_\pm$ be the
conjugates of $\xi_\pm$ by $\dot w$, and put
$\Lat^w_-=\Lie(\cU^-_w)(\OO)$. Consider the maps
\begin{equation}\label{eq:integralchart}
f_{P,w}:\Lat^w_-\times(\uu_{P_w}^\perp\cap \dLat)\to \dLat,\qquad
f_{P,w}(y,\nf)=\Ad^*(\xi^w_-(y))\nf,
\end{equation}
and the parametrizations $\psi_{P,w}(y,\nf)=\xi^w_-(y)\cdot(\bP_w,\nf)$ of
$\Inc_P$.

\begin{lemma}[An \iatlas{} with a uniform number of charts]\label{lem:integralatlas}
The maps $f_{P,w}$ are polynomial with coefficients
in $\dLat$, and the $\psi_{P,w}$, for $\bP\in\Pc$ and $w\in W_k/W_{k,P}$, form
an \iatlas{} of $p$ over $\dLat$ with at most $\Nss$ charts, with
$\Lat_i=\Lat^w_-\oplus(\uu_{P_w}^\perp\cap \dLat)$. Consequently, the
\bnets{} \eqref{eq:badf} of this atlas satisfy
\begin{equation}\label{eq:badexplicit}
\nu_{\dLat}(\Bad_h)\le \Nss\,q^{-h}\qquad(h\ge1).
\end{equation}
\end{lemma}

\begin{proof}
The map \eqref{eq:integralchart} is polynomial with integral
coefficients because $\cU^-_w$ acts on $\Lie(\cG)$ by an action of
$\OO$-group schemes and $\xi^w_-$ is an isomorphism of $\OO$-schemes
mapping $\Lat^w_-$ onto $\cU^-_w(\OO)$, and
$\psi_{P,w}$ is the trivialization \eqref{eq:trivialization} of $\Inc_P$ over the big cell of $\bP_w$, hence a parametrization of an open subset of $\Inc_P$. If
$(\bQ,\lf)\in p_P^{-1}(\dLat)$ then $\bQ=u\bP_wu^{-1}$ with $u\in\cU_w^-(\OO)$ by
\eqref{eq:integralcover}, and $\nf=\Ad^*(u^{-1})\lf\in \dLat$ vanishes on
$\uu_{P_w}$. The number of charts is at most $\Nss$, since the $\bP_w$ are
pairwise distinct elements of $\Pcss$ (\cref{sec:integraldata}), and those
attached to different $\bP\in\Pc$ are not conjugate. The last assertion is \cref{lem:tube} applied to the
$f_{P,w}$, which have coefficients in $\OO$.
\end{proof}

\begin{proposition}[Geometric alternatives in the unramified case]\label{prop:geometricunram}
In the unramified case (\cref{sec:integraldata}), the sets $\Bad_h$ of
\cref{prop:geometric} can be chosen so that moreover
$\nu_{\dLat}(\Bad_h)\le \Nss q^{-h}$ for all $h\ge1$, with $\mathcal C\subset\cG(\OO)$,
$\mathcal A_P=\uu_P^\perp\cap\dLat$ and $c=0$.
\end{proposition}

\begin{proof}
Let $\Bad_h$ be the \bnets{} of the \iatlas{} of \cref{lem:integralatlas}.
The proof of \cref{prop:geometric} applies to this atlas, with
$\mathcal C=\bigcup_w\xi_-^w(\Lat_-^w)\subset\cG(\OO)$ and
$\mathcal A_{P_w}=\uu_{P_w}^\perp\cap\dLat$, and \eqref{eq:badexplicit} gives
$\nu_{\dLat}(\Bad_h)\le \Nss q^{-h}$. It remains to see that $c=0$. If
$\lf=f_{P,w}(y,\nf)$ lies outside $\Bad_h$, then
$\abs{\det df_{P,w}(y,\nf)}>q^{-h}$ by \eqref{eq:sardchart}. In the
factorization of $df_{P,w}$ in the proof of \cref{lem:levi}, the coadjoint
action of $u=\xi_-^w(y)$ and the differential $T_y$ are lattice
automorphisms, so $\abs{\det Q_\nf}=\abs{\det df_{P,w}(y,\nf)}>q^{-h}$. As
$Q_\nf$ is integral, Cramer's rule gives $\norm{Q_\nf^{-1}}<q^h$.
\end{proof}




\section{Regular cells detect no supercuspidal representation}\label{sec:regular}

We prove the counterpart of \cref{prop:elliptic}: a cell admitting a
quantitatively regular parabolic lift cannot occur in a supercuspidal
representation. Throughout this section, fix $\bP=\bM\bU\in\Pcss$
and a cocharacter $\eta:\Gm\to Z(\bM)$ defining $\bP$ and $\bM$.
Write $\bP^-=\bM\bU^-$ for the opposite parabolic defined by $-\eta$.
Recall that a form $\nf\in\uu^\perp$ defines the map
\[
Q_\nf:\uu^-\longrightarrow\uu^*,\qquad
(Q_{\nf}Y)(X)=\nf([X,Y]).
\]
We fix a norm on $\operatorname{Hom}_F(\uu^*,\uu^-)$ as in \cref{sec:levi}.

\begin{theorem}[Regular cells]\label{thm:regular}
Fix compact sets $\mathcal C\subset G$ and $\mathcal K\subset\uu^\perp$.
There is a constant $C$ with the following property. Let $h\in\Z_{\ge0}$ and let
$\lf\in\dLat$ be of the form $\lf=\Ad^*(g)\nf$ with $g\in\mathcal C$,
$\nf\in\mathcal K$ and $\norm{Q_\nf^{-1}}\le q^h$. If $n$ and $r\ge1$ are integers with
\begin{equation}\label{eq:margins}
n-2r\ge C,\qquad 3r-n-2h\ge C,
\end{equation}
then no supercuspidal representation of $G$ contains the character
\[
\chi_{n,\lf}(k)=\psi\bigl(\varpi^{-n}\lf(\kappa(k))\bigr),
\qquad k\in K_{n-r}.
\]
In particular, $f_{n,r}(\lf)=0$.
\end{theorem}

In the unramified case the margins can be made explicit
(\cref{thm:regularunram}).

The constant $C$ may depend on the compact sets $\mathcal C,\mathcal K$
and the chosen norm. Once these data are fixed for each
$\bP\in\Pcss$, finiteness of $\Pcss$ gives a common $C$, independent
of $n,r,h,\lf$ and $g$.

We first work in product coordinates, writing $(J,\theta_J)$ for the
group and its cell character. We enlarge $J$ to a group $H$ and
normalize an extension $\theta$ of $\theta_J$ so that it is trivial
on $S=H\cap U$.
We then show that every nonzero vector in $V^{H,\theta}$ survives
passage to the Jacquet module $V_U$. The key point is a simple
averaging identity: after averaging over a larger subgroup of $U$,
projection back to $V^{H,\theta}$ recovers a nonzero multiple of
the vector. Finally, conjugating by $g$ and comparing charts to return to the original
cell (\cref{sec:returncell}) we obtain the result. The coordinate estimates and the
verification of the group constructions are collected in
\cref{sec:jacquetdetails}.



The proof is closely related to the mechanism of the theory of covers of Bushnell and Kutzko
\cite[\S\S6--8]{MR1643417}. The contracting-average identity behind \cref{cor:jacquet} is already in Deligne's
proof that, for $\bG$ semisimple, the character of a supercuspidal representation is supported on the compact elements \cite{MR0425033}.
See also \cite[Proposition~6.7]{MR1371680} and \cite[\S17.3]{MR2276772} for closely related statements.



\subsection{The statement in the product chart}\label{sec:localsetting}

Call a lattice in $\g$ \emph{\etacompat{}} if it is the direct sum
of its intersections with the $\eta$-weight spaces. Retain the normalized
parametrization $\xi_-$ from \cref{sec:incidencejacobian}, and choose
normalized parametrizations $\xi_+$ of $U$ and $\xi_M$ of $M$, with
$\xi_+$ an isomorphism of $F$-varieties. The product chart is
\begin{equation}\label{eq:productchart}
\kappa'\bigl(\xi_+(X)\xi_M(Z)\xi_-(Y)\bigr)=X+Z+Y.
\end{equation}
As justified in \cref{sec:estimates}, choose an \etacompat{} lattice
$\Lat'$ small enough that $(\kappa',\Lat')$ is standard, conjugation
by $t=\eta(\varpi)$ on $U$ and by $t^{-1}$ on $U^-$ has integral
coordinate coefficients, and $\mathcal K\subset(\Lat')^\vee$.
Write $\Lat'=\Lat_+\oplus\Lat_M\oplus\Lat_-$, where
$\Lat_\pm=\Lat'\cap\uu^\pm$ and $\Lat_M=\Lat'\cap\m$.

For $d\ge1$, write $U_d=\xi_+(\pp^d\Lat_+)$, $M_d=\xi_M(\pp^d\Lat_M)$ and
$U^-_d=\xi_-(\pp^d\Lat_-)$. We may prove \cref{thm:regular} using
the operator norm induced by these lattices and their duals.
Indeed, equivalence of norms gives a fixed integer $d\ge0$ such that
a bound $q^h$ in the original norm implies a bound $q^{h+d}$ in the
new norm. Replacing $h$ by $h+d$ and increasing $C$ by $2d$ in
\eqref{eq:margins} absorbs this change.
Henceforth in this section and \cref{sec:jacquetdetails}, all norms
are induced by $\Lat'$ and its summands and duals; we use the same
notation $\norm{\cdot}$. Write $\val(Z)=-\log_q\norm{Z}$, with
$\val(0)=+\infty$, and measure group depths in these coordinates.
These parametrizations and lattices remain fixed throughout
\cref{sec:localsetting,sec:groups,sec:jacquet}.

For each integer $h\ge0$, define
\begin{equation}\label{eq:Qbound}
\mathcal A_h=\{\lf\in(\Lat')^\vee\cap\uu^\perp:
Q_\lf\text{ is invertible and }\norm{Q_\lf^{-1}}\le q^h\}.
\end{equation}
Fix $h\ge0$ and $\lf\in\mathcal A_h$, and put $Q=Q_\lf$.
The same inverse bound holds for $Q^*$. In particular, for every
$X\in\uu$ there is $Y\in\Lat_-$ with
$\abs{\Bf_\lf(X,Y)}\ge q^{-h}\norm{X}$, and the analogous statement
holds with the two factors interchanged.

Let $b,r\ge1$ be integers: $b$ is the depth of the
subgroup and $r$ the precision of the cell. Put $n=b+r$. The assumptions are
\begin{equation}\label{eq:localmargins}
b\ge r+C,\qquad 2(r-h)\ge b+C.
\end{equation}
These are \eqref{eq:margins} with $n=b+r$ and the lattice-induced norm.
The first makes the cell character well defined; the second controls the commutators in the
enlarged groups used below. Per our standard convention, constants are independent of $h$ and $\lf\in\mathcal A_h$.
As we shall verify in \cref{lem:groups}, for $C$ sufficiently large,
\[
J=U_bM_bU^-_b=\kappa'^{-1}(\pp^b\Lat')
\]
is a compact open subgroup, and its cell character associated with $\lf$ is given by
\begin{equation}\label{eq:thetaJ}
\theta_J\bigl(\xi_+(X)\xi_M(Z)\xi_-(Y)\bigr)=\psi\bigl(\varpi^{-n}\lf(Z+Y)\bigr).
\end{equation}

We will deduce \cref{thm:regular} from the following proposition.

\begin{proposition}[Jacquet nonvanishing]\label{prop:local}
Let $\lf\in\mathcal A_h$ and assume \eqref{eq:localmargins}.
Every smooth representation $V$ of $G$ containing $\theta_J$ has $V_U\ne0$.
\end{proposition}

The proof occupies \cref{sec:groups,sec:jacquet}.

\subsection{Extending and normalizing the character}\label{sec:groups}
We enlarge $U_b$ to the subgroup determined by the annihilator lattice
\begin{equation}\label{eq:Slattice}
\sS=\{X\in\uu:\Bf_\lf(X,\pp^b\Lat_-)
\subset\pp^n\}
=\pp^r(Q^*)^{-1}\Lat_-^\vee,\qquad S=\xi_+(\sS).
\end{equation}
Its size is controlled by
\[
\pp^r\Lat_+\subset\sS\subset\pp^{r-h}\Lat_+.
\]
Set
\begin{equation}\label{eq:JH}
R=U^-_b,\qquad H=SM_bR\supset J.
\end{equation}
The groups $J$ and $R$ are independent of $\lf$, whereas $S$ and $H$
depend on $\lf$ through $Q_\lf$; we suppress this dependence in the notation.
By \cref{lem:groups}, $J$ is normal in $H$, the quotient
$H/J\simeq\sS/\pp^b\Lat_+$ is finite abelian, and the formula
\eqref{eq:thetaJ} extends $\theta_J$ to a character $\theta_\lf$ of
$H$ that is trivial on $S$.

Consequently $V^{J,\theta_J}$ decomposes into the spaces
$V^{H,\theta}$ as $\theta$ runs through the extensions of $\theta_J$.
Every extension other than $\theta_\lf$ is nontrivial on $S$.
The next lemma shows that any extension can be conjugated to one
trivial on $S$.

\begin{lemma}[Normalization of an extension]\label{lem:normalize}
Every character $\theta$ of $H$ extending $\theta_J$ admits an
element $y\in U^-$ normalizing $H$ such that
$\theta^y=\theta_{\lf'}$ is trivial on $S$, where
$\theta^y(z)=\theta(yzy^{-1})$ for $z\in H$.
The form $\lf'$ belongs to $\mathcal A_h$ and defines the same groups $S,H$.
\end{lemma}

The proof of \cref{lem:normalize}, including the choice of $y$ and
the bounds for $\lf'$, is given in \cref{sec:normalizationdetails}.
To prove \cref{prop:local}, suppose $V^{J,\theta_J}\ne0$. The
decomposition above and \cref{lem:normalize} then give
$V^{H,\theta_{\lf'}}\ne0$ for some $\lf'\in\mathcal A_h$.
Thus it remains to prove that $V^{H,\theta_\lf}\to V_U$ is injective
for every $\lf\in\mathcal A_h$, under the depth assumptions
\eqref{eq:localmargins}.

\subsection{Why averaging cannot kill the vector}\label{sec:jacquet}
For a compact subgroup $C$ let $\Avg_C$ denote averaging with Haar
measure of total mass $1$. If $\theta$ is a character of $C$, let
$\Avg_C^\theta$ be the projection onto $V^{C,\theta}$.
The following elementary observation is the core of the proof. 

\begin{lemma}[One-step averaging]\label{lem:criterion}
Let $H$ be a compact open subgroup of $G$, $\theta$ a character of
$H$, and $S\subset S^+$ compact subgroups with $S\subset H$ and
$\theta|_S=1$, where $[S^+:S]<\infty$. Suppose that for every
$x\in S^+\setminus S$ there is $u\in H$ such that
$xux^{-1}\in H$ and $\theta(xux^{-1})\ne\theta(u)$.
Then, for $E=\Avg_H^\theta$,
\begin{equation}\label{eq:leftinverse}
E\Avg_{S^+}v=[S^+:S]^{-1}v\qquad(v\in V^{H,\theta}).
\end{equation}
In particular, $\Avg_{S^+}$ is injective on $V^{H,\theta}$.
\end{lemma}

\begin{proof}
For $x\in S^+\setminus S$, choose $u$ as in the statement. Then
\[
\theta(u)E\pi(x)v=E\pi(x)\pi(u)v=E\pi(xux^{-1})\pi(x)v=\theta(xux^{-1})E\pi(x)v,
\]
so $E\pi(x)v=0$. On the other hand, since $v$ is $S$-invariant,
\[
\Avg_{S^+}v=[S^+:S]^{-1}\sum_{x\in S^+/S}\pi(x)v.
\]
Applying $E$ kills every coset except $S$, which contributes $v$.
\end{proof}

To iterate this observation, put $t=\eta(\varpi)$ and
$\lf_j=\Ad^*(t^j)\lf$. For every $j\ge0$, \cref{lem:chain} gives
$\lf_j\in\mathcal A_h$.
Applying the preceding construction to $\lf_j$, with $b,r,n,h$ fixed,
defines $S_j,H_j,\theta_j$, with $S_0=S$, $H_0=H$ and
$\theta_0=\theta_\lf$.
That lemma also gives
\[
tS_jt^{-1}\subset S_{j+1},\qquad
\theta_{j+1}(u)=\theta_j(t^{-1}ut)\quad(u\in M_bR),
\]
and shows that the increasing compact subgroups
$C_j=t^{-j}S_jt^j$ exhaust $U$. Finally, the hypothesis of
\cref{lem:criterion} holds for $H_j$, $\theta_j$ and
$S_j^+=t^{-1}S_{j+1}t$: the commutator pairing supplies the element
$u\in R$ that separates each new coset.

\begin{corollary}[Injectivity of the Jacquet map]\label{cor:jacquet}
Let $\lf\in\mathcal A_h$, assume \eqref{eq:localmargins}, and let $H$,
$\theta_\lf$ and $U=U_P$ be as in \cref{sec:groups}. Then
for every smooth representation $V$ of $G$, the canonical map
$V^{H,\theta_\lf}\longrightarrow V_U$ is injective.
\end{corollary}

\begin{proof}
Fix $v\in V^{H,\theta_\lf}$, and put
\[
v_j=\Avg_{S_j}\pi(t^j)v,\qquad E_j=\Avg_{H_j}^{\theta_j}.
\]
Since $\theta_\lf|_S=1$, we have $v_0=v$. The nesting
$tS_jt^{-1}\subset S_{j+1}$ gives
$v_{j+1}=\Avg_{S_{j+1}}\pi(t)v_j$.
Inductively, $v_j\in V^{H_j,\theta_j}$: the inclusion
$t^{-1}Rt\subset R$ and character compatibility on $M_bR$ make
$\pi(t)v_j$ transform there by $\theta_{j+1}$, and
\[
E_{j+1}=\Avg_{S_{j+1}}\Avg_{M_bR}^{\theta_{j+1}}.
\]
This factorization follows from the unique product decomposition
$H_{j+1}=S_{j+1}(M_bR)$ and the triviality of $\theta_{j+1}$ on
$S_{j+1}$, with all Haar measures normalized to have total mass one.
The one-step averaging identity, applied to $H_j$ and $S_j^+$, gives
\[
E_j\pi(t^{-1})v_{j+1}=[S_j^+:S_j]^{-1}v_j.
\]
Thus $v_j\ne0$ for every $j$ if $v\ne0$. On the other hand,
\begin{equation}\label{eq:telescope}
\pi(t^{-j})v_j=\Avg_{C_j}v.
\end{equation}
Since the $C_j$ exhaust $U$, the kernel of $V\to V_U$ is the union of
the kernels of $\Avg_{C_j}$ \cite[\S3.2]{CassNotes}. Hence $v$ has
nonzero image in $V_U$ whenever $v\ne0$.
\end{proof}

\begin{proof}[Proof of \cref{prop:local}]
If $V^{J,\theta_J}\ne0$, choose an extension $\theta$ with
$V^{H,\theta}\ne0$. By \cref{lem:normalize}, there are $y\in U^-$
normalizing $H$ and $\lf'\in\mathcal A_h$ such that
\[
\pi(y^{-1})V^{H,\theta}=V^{H,\theta^y}
=V^{H,\theta_{\lf'}}.
\]
The latter space is nonzero and injects into $V_U$ by
\cref{cor:jacquet}, so $V_U\ne0$.
\end{proof}

\subsection{Return to the original cell}\label{sec:returncell}
\begin{proof}[Proof of \cref{thm:regular}]
Use the lattice-induced norm as justified in \cref{sec:localsetting}.
Write $\lf=\Ad^*(g)\nf$ as in the hypothesis. Since
$\nf\in\mathcal K\subset(\Lat')^\vee$ and $\nf|_\uu=0$, the bound
$\norm{Q_\nf^{-1}}\le q^h$ gives $\nf\in\mathcal A_h$.
Set
\[
\kappa_g(x)=\Ad(g)\kappa'(g^{-1}xg),\qquad
\Lat_g=\Ad(g)\Lat'.
\]
Compactness gives a fixed $s\ge1$ with
$\pp^s\Lat_g\subset\Lat$ for every $g\in\mathcal C$.
Apply \cref{lem:compare} to this family of transported pairs, with
$\mathcal B=\dLat$. Its constants are uniform in $g$.
Put $b=n-r$ and $r'=r-s$. Then
\[
(b+s)-r'=n-2r+2s,\qquad
2(r'-h)-(b+s)=3r-n-2h-3s.
\]
Choose $C$ large enough for the chart comparison and to absorb the
fixed loss $3s$, and with $C\ge s$. Since \eqref{eq:margins} gives
$r\ge2h+2C$, this ensures $r'=r-s\ge C\ge1$. The margins \eqref{eq:margins} then imply the
hypotheses of \cref{prop:local} at depth $b+s$, with precision $r'$
and inverse bound $h$. No supercuspidal representation contains
the resulting character of $J'=U_{b+s}M_{b+s}U^-_{b+s}$.

After conjugation by $g$, this is the character
$x\mapsto\psi(\varpi^{-n}\lf(\kappa_g(x)))$ on
$J_g=gJ'g^{-1}$. By \cref{lem:compare}, $J_g\subset K_b$ and this
character is the restriction of $\chi_{n,\lf}$. A representation
containing $\chi_{n,\lf}$ would contain its restriction to $J_g$,
so cannot be supercuspidal.
\end{proof}

\subsection{The unramified case}\label{sec:regularunram}

\begin{theorem}[Regular cells in the unramified case]\label{thm:regularunram}
In the unramified case (\cref{sec:integraldata}), the conclusion of
\cref{thm:regular} holds with $\mathcal C=\cG(\OO)$ and
$\mathcal K=\uu^\perp\cap\dLat$ whenever
\begin{equation}\label{eq:explicitmargins}
n\ge2r,\qquad 3r-n-2h\ge2m_{\cG}+1,
\end{equation}
with $m_{\cG}$ as in \eqref{eq:mG}.
\end{theorem}

\begin{proof}
We may assume that $\bP=\dot w\bP_I\dot w^{-1}$ with $\bP_I\in\Pc$ standard and
$\dot w\in N_{\cG}(\cS)(\OO)$, and conjugate the data of
\cref{sec:integraldata} by $\dot w$. Take $\Lat'=\Lat$, and for $\kappa'$ the
product chart \eqref{eq:productchart} built from $\xi_\pm$ and from the chart
of \cref{lem:model} for $\cM$. It has integral coefficients and, by
\eqref{eq:integralcell}, identifies $K_b$ with $U_bM_bU^-_b$ for $b\ge1$, where
$U^\pm_b=\cU^\pm(\OO)\cap K_b$ and $M_b=\cM(\OO)\cap K_b$. Its transition map
to $\kappa$ is $X+O(X^2)$ with integral coefficients, so \cref{lem:compare}
holds with $b_0=1$, $s=0$ and $C=0$ for $(\kappa',\Lat)$ and its conjugates by
elements of $\cG(\OO)$. Every $\nf\in\uu^\perp\cap\dLat$ is integral on
$\Lat$, and the conjugations by $t$ on $U$ and by $t^{-1}$ on $U^-$ have
integral coefficients. Thus the requirements of \cref{sec:localsetting} hold
without rescaling or change of norm, and
\cref{lem:estimates,lem:contraction} apply directly, with $w_+\le m_{\cG}$.
With $a=r-h$ and $b=n-r$, the depth inequalities \eqref{eq:used} become
\[
2a\ge b,\qquad b\ge r,\qquad r\ge2h+2m_{\cG}+1.
\]
They follow from \eqref{eq:explicitmargins}: the first margin gives $b\ge r$,
and the second gives
\[
2a\ge b+2m_{\cG}+1,\qquad
r-2h\ge2m_{\cG}+1+(n-2r)\ge2m_{\cG}+1,
\]
whence $a\ge h+2m_{\cG}+1$ and $b\ge r\ge2h+2m_{\cG}+1$. Finally, $g\in\cG(\OO)$
and $s=0$, so the return to the original cell (\cref{sec:returncell})
introduces no shift, and the proof concludes as before.
\end{proof}

\section{Statements and proofs of the main results}\label{sec:proofs}

We can now state and prove our main results on the relation between supercuspidal and elliptic densities
for principal congruence subgroups.

\subsection{Statements}\label{sec:statements}
We use the notation of \cref{sec:notation}: $(\kappa,\Lat)$ is
the fixed standard pair of \cref{sec:standardgroups}, $(K_n)$ its
principal congruence subgroups, $\dLat=\Lat^\vee$, $\scd{K}$ the
supercuspidal density (\cref{sec:representations}),
$\elld{\dLat}$ the elliptic density (\cref{sec:ellipticforms}), and
$\eexp$ the exponent \eqref{eq:tau}.

\begin{theorem}[Supercuspidal Plancherel mass]\label{thm:main}
We have $0<\elld{\dLat}\le1$ and
\begin{equation}\label{eq:main}
\scd{K_n}=\elld{\dLat}+O(q^{-\eexp n/4})\qquad(n\to\infty).
\end{equation}
In particular, if $Z(\bG)(F)$ is compact then
\begin{equation}\label{eq:formal}
\sum_{\pi\in\Gsc}\fdeg{\pi}(K_n)\dim\pi^{K_n}=\elld{\dLat}+O(q^{-\eexp n/4}).
\end{equation}
\end{theorem}

The implied constants depend on $F$, $\bG$ and $\Lat$; the exponent
$\eexp/4\ge1/4$ depends only on $\bG$ (see \cref{rem:sharpness}).

For a smooth model $\cG$ as in \cref{lem:model}, with
$\Lat=\Lie(\cG)(\OO)$ and $\vol(\cG(\OO))=1$, substituting
$\vol(K_n)^{-1}=\abs{\cG(k)}q^{(n-1)D}$ into \cref{thm:main} gives
\[
\int_{\Gsc}\dim\pi^{K_n}\,d\mu_G(\pi)=\elld{\dLat}\abs{\cG(k)}q^{(n-1)D}+O\bigl(q^{(D-\eexp/4)n}\bigr).
\]

The elliptic density $\elld{\dLat}$ depends on the lattice and not
only on $\bG$; it is studied in \cref{sec:positivity}, where we also
determine when it equals $1$ (\cref{thm:endpoint}).

When $\bG$ is unramified, that is, quasi-split and split over an unramified
extension of $F$, the constants can be made explicit. Let $\Nss$ be the number of maximal semistandard parabolic subgroups of
$\bG$ (\cref{sec:parabolics}; for $\GL_d$ we have $\Nss=2^d-2$), and let
$m_{\cG}$ be the invariant \eqref{eq:mG} of the root datum of $\bG$ with its
Galois action, the largest weight on the unipotent radicals of the
cocharacters defining the maximal parabolic subgroups.

\begin{theorem}[Explicit uniform bound]\label{thm:uniform}
Suppose that $\bG$ is unramified, let $\cG$ be a reductive model of $\bG$ over
$\OO$, $\Lat=\Lie(\cG)(\OO)$ and $K_n=\ker(\cG(\OO)\to\cG(\OO/\pp^n))$. Then
for every $n\ge4m_{\cG}+9$,
\begin{equation}\label{eq:explicitbound}
\bigl|\scd{K_n}-\elld{\dLat}\bigr|\le \Nss\,q^{\,m_{\cG}+2-n/4}.
\end{equation}
\end{theorem}

\begin{corollary}[Chevalley groups]\label{cor:chevalley}
Let $\cG$ be a Chevalley group scheme over $\Z$ with generic fibre $\bG$. For
every non-archimedean local field $F$, of any characteristic, with
$K_n=\ker(\cG(\OO)\to\cG(\OO/\pp^n))$ and $\Lat=\Lie(\cG)(\OO)$,
\[
\bigl|\scd{K_n}-\elld[\bG_F]{\dLat}\bigr|\le \Nss\,q^{\,m_{\cG}+2-n/4}\qquad(n\ge4m_{\cG}+9),
\]
where $\Nss$ and $m_{\cG}$ depend only on the root datum.
\end{corollary}

For instance, for $\bG=\GL_d$, $d\ge2$, the estimate becomes
\[
\abs{\scd{K_n}-\elld[\GL_d]{\Mat_d(\OO)}}\le(2^d-2)q^{3-n/4},\ \ n\ge13.
\]

For a reductive group over a global field with a fixed linear embedding, the
towers defined by the embedding come from reductive models at almost all
places. This gives a bound uniform over all places.

\begin{theorem}[Uniformity over a global field]\label{thm:global}
Let $k$ be a global field, $\bG$ a connected reductive group over $k$ and
$\bG\subset\GL_d$ a closed embedding over $k$. For a finite place $v$ of $k$, with ring of integers $\OO_v$, maximal ideal $\pp_v$
and residue field of cardinality $q_v$, let $\Lat_v=\g\cap\Mat_d(\OO_v)$,
$K_{v,0}=\bG(k_v)\cap\GL_d(\OO_v)$ and
$K_{v,n}=\bG(k_v)\cap(1+\pp_v^n\Mat_d(\OO_v))$ for $n\ge1$, and let $\scd{K_{v,n}}$ be
the supercuspidal density of $K_{v,n}$ in $\bG(k_v)$. There are a finite set
$S$ of places of $k$ and an integer $m_{\bG}\ge0$ depending only on $\bG$
(\cref{sec:global}) with the following properties. For $v\notin S$, the
schematic closure $\cG_v$ of $\bG$ in $\GL_{d,\OO_v}$ is a reductive group
scheme, $\Lat_v=\Lie(\cG_v)(\OO_v)$, $K_{v,0}=\cG_v(\OO_v)$ is a hyperspecial
parahoric subgroup, $K_{v,n}=\ker(\cG_v(\OO_v)\to\cG_v(\OO_v/\pp_v^n))$ for
$n\ge1$, and
\begin{equation}\label{eq:globalgood}
\bigl|\scd{K_{v,n}}-\elld[\bG_{k_v}]{\Lat_v^\vee}\bigr|\le N_{\bG_{k_v}}\,q_v^{\,m_{\bG}+2-n/4}
\qquad(n\ge4m_{\bG}+9).
\end{equation}
Moreover, there are constants $n_0$ and $C$, depending on $\bG$ and the embedding, such that
\begin{equation}\label{eq:globalall}
\bigl|\scd{K_{v,n}}-\elld[\bG_{k_v}]{\Lat_v^\vee}\bigr|\le C\,q_v^{\,m_{\bG}+2-n/4}
\end{equation}
for all finite places $v$ and all $n\ge n_0$.
\end{theorem}


\subsection{Proof of \Cref{thm:main}}\label{sec:proofsgeneral}

Recall that $f_{n,r}$, defined in \eqref{eq:f}, is constant on each
cell and records the supercuspidal proportion of its Plancherel mass.
In particular, $0\le f_{n,r}\le1$ and
\[
\scd{K_n}=\int_\dLat f_{n,r}\,d\nu_{\dLat}.
\]

We first compare the cellwise masses; the main estimates follow by integration.

\begin{remark}[Choice of parameters]\label{rem:parameters}
For each $\bP\in\Pcss$, apply \cref{thm:regular} with
$\mathcal K=\mathcal A_P$ and the compact set $\mathcal C$ supplied
by \cref{prop:geometric}.
Take the maximum of the constants in \cref{thm:regular} over the
finite set $\Pcss$. Fix $C\ge1$ larger than this maximum and the
constant in \cref{prop:elliptic}, with an additional allowance of $2c$.

If integers $n,r,h$ with $h\ge1$ satisfy \eqref{eq:margins}, then $n\ge2r$ and
$r\ge2h+2C$. In particular, each cell in $\dLat/\pp^r\dLat$ is
either contained in $\Bad_h$ or disjoint from it.
For suitable fixed nonnegative integers $C_1,C_2$ (for instance any $C_1\ge C/2$ and $C_2\ge(C+3C_1+2)/2$) the parameters
\[
r_n=\lfloor n/2\rfloor-C_1,\qquad
h_n=\lfloor n/4\rfloor-C_2
\]
satisfy \eqref{eq:margins} and $h_n\ge1$ for all sufficiently large $n$,
and by \cref{prop:geometric} $\nu_{\dLat}(\Bad_{h_n})\ll q^{-\eexp n/4}$.
\end{remark}

\begin{proposition}[Cellwise comparison]\label{prop:density}
Suppose that integers $n,r,h$ with $h\ge1$ satisfy \eqref{eq:margins}.
Let $\cell=\lf+\pp^r\dLat$ be a cell in $\dLat$ disjoint from
$\Bad_h$, and let $\chi_{n,\lf}$ be its associated character of
$K_{n-r}/K_n$, independent of the representative $\lf$.
Then $\cell$ is contained in either $\Ell{\gd}$ or $\Par{\gd}$.
In the elliptic case, every irreducible representation containing
$\chi_{n,\lf}$ is supercuspidal; in the parabolic case, no
supercuspidal representation contains $\chi_{n,\lf}$.
Consequently, on $\dLat$,
\begin{equation}\label{eq:cellcomparison}
\abs{f_{n,r}-\one_{\Ell{\gd}}}\le\one_{\Bad_h},
\end{equation}
and
\begin{equation}\label{eq:masscomparison}
-\nu_{\dLat}(\Bad_h\cap\Ell{\gd})
\le\scd{K_n}-\elld{\dLat}
\le\nu_{\dLat}(\Bad_h\cap\Par{\gd}).
\end{equation}
\end{proposition}

\begin{proof}
By \cref{rem:parameters}, $r\ge2h+C$. For
$\lf\in\dLat\setminus\Bad_h$, \cref{prop:geometric} gives an elliptic
coset $\lf+\pp^{2h}\dLat$, or a parabolic lift with
$\norm{Q_{\nf,\bP}^{-1}}\le q^{h+c}$. In the parabolic case, apply
\cref{thm:regular} with $h+c$ in place of $h$; the extra $2c$ in
the choice of $C$ ensures its depth assumptions.
In the elliptic case, $\erad(\lf)\le2h$, so \cref{prop:elliptic}
applies. This proves both
assertions about representations and hence
$f_{n,r}=\one_{\Ell{\gd}}$ outside $\Bad_h$.
Since $0\le f_{n,r}\le1$, \eqref{eq:cellcomparison} follows.
The function $f_{n,r}-\one_{\Ell{\gd}}$ is nonpositive on
$\Ell{\gd}$ and nonnegative on $\Par{\gd}$, so integration using
\eqref{eq:Pscdecomp} gives \eqref{eq:masscomparison}.
\end{proof}

\begin{proof}[Proof of \cref{thm:main}]
Apply \eqref{eq:masscomparison} with $r=r_n$ and $h=h_n$.
By \cref{rem:parameters}, we obtain
\[
\abs{\scd{K_n}-\elld{\dLat}}
\le\nu_{\dLat}(\Bad_{h_n})\ll q^{-\eexp n/4}.
\]
This proves the asymptotic \eqref{eq:main}.
The remaining assertion $\elld{\dLat}>0$ is supplied by
\cref{thm:endpoint} in \cref{sec:positivity}; the upper bound
$\elld{\dLat}\le1$ follows from the definition.
\end{proof}

\begin{remark}[The lower bound needs only elliptic cells]\label{rem:lower}
The first inequality in \eqref{eq:masscomparison}, and hence
the lower half of \eqref{eq:main}, does not require \cref{thm:regular}.
In particular, if $\elld{\dLat}=1$, then
$0\le1-\scd{K_n}\ll q^{-n/4}$.
\end{remark}

\begin{proof}[Proof of \eqref{eq:formal}]
If $Z(\bG)(F)$ is compact, the supercuspidal representations are
square-integrable, hence atoms of $\mu_G$ whose masses are their formal
degrees \cite[Theorem~VIII.1.1]{MR1989693}. Viewing $\fdeg{\pi}$ as a Haar
measure on $G$, we have $\fdeg{\pi}(K_n)=\vol(K_n)\mu_G(\{\pi\})$, so the
left-hand side of \eqref{eq:formal} is $\scd{K_n}$, and \eqref{eq:formal} is
\eqref{eq:main}.
\end{proof}


For a fixed compact subset of $\dLat\setminus\Crv$, the \bnets{} are
eventually disjoint from it, so the cellwise comparison is eventually
exact there. Globally, the examples of $\GL_2$ and $\PGL_2$ in
characteristic $2$ show that the error need not vanish. The localized
form of the estimate is as follows.

\begin{corollary}[Localization]\label{cor:localized}
Let $\cell\subset \dLat$ be a union of cosets of $\pp^m\dLat$ and, for $n$ with
$r_n\ge m$, let $e_{n,\cell}$ be the sum of the idempotents $e_{n,\cell'}$ over
the cosets $\cell'$ of $\pp^{r_n}\dLat$ contained in $\cell$; it does not depend
on $r_n$ (\cref{sec:cells}). Then
\begin{equation}\label{eq:localized}
\vol(K_n)\int_{\Gsc}\tr\pi(e_{n,\cell})\,d\mu_G(\pi)
=\nu_{\dLat}(\cell\cap\Ell{\gd})+O(q^{-\eexp n/4}),
\end{equation}
and if $\cell\cap\Crv=\emptyset$ then the left-hand side equals
$\nu_{\dLat}(\cell\cap\Ell{\gd})$ for $n\gg1$.
\end{corollary}

\begin{proof}[Proof of \cref{cor:localized}]
We have
$\vol(K_n)\int_{\Gsc}\tr\pi(e_{n,\cell})\,d\mu_G=\int_\cell f_{n,r_n}\,d\nu_{\dLat}$ by
\eqref{eq:f}, and \eqref{eq:localized} follows from \eqref{eq:cellcomparison}
with $r=r_n$ and $h=h_n$ and from \cref{rem:parameters}. If
$\cell\cap\Crv=\emptyset$, then $\cell\cap\Bad_h=\emptyset$ for $h\gg1$ by
\eqref{eq:bncompact}, so that $f_{n,r_n}=\one_{\Ell{\gd}}$ on $\cell$
for $n\gg1$ by \cref{prop:density}.
\end{proof}

\subsection{The unramified case}\label{sec:uniform}
We now make the constants explicit in the unramified case
(\cref{sec:integraldata}). In general, the constants in the proof of
\cref{thm:main} come from three sources: a finite \iatlas{} of the incidence
map $p$, the shifts needed to compare the original tower with the auxiliary
product charts and norms, and depth thresholds governed by the weights of
the defining cocharacters. In the unramified case integral coordinates
eliminate the comparison shifts, and the number of charts and the depth
thresholds are bounded in terms of the root datum, by
\cref{lem:integralatlas}, the unramified case of \cref{prop:elliptic}, and
\cref{prop:geometricunram,thm:regularunram}.

\begin{proposition}[Explicit cellwise comparison]\label{prop:explicitcell}
In the unramified case, let $(\Bad_h)$ be the \bnets{} of
\cref{lem:integralatlas}, so that $\nu_{\dLat}(\Bad_h)\le \Nss q^{-h}$. If integers
$n,r,h$ with $h\ge1$ satisfy \eqref{eq:explicitmargins}, then
$\abs{f_{n,r}-\one_{\Ell{\gd}}}\le\one_{\Bad_h}$ on $\dLat$.
\end{proposition}

\begin{proof}
If $\Pc=\emptyset$, then $\Ell{\gd}=\gd$, every irreducible representation
is supercuspidal and $f_{n,r}=1$. Otherwise \eqref{eq:explicitmargins} gives
$r\ge2h+2m_{\cG}+1\ge2h$, so each cell $\lf+\pp^r\dLat$ is either contained
in $\Bad_h$ or disjoint from it. Let it be disjoint. By
\cref{prop:geometricunram}, either $\lf+\pp^{2h}\dLat\subset\Ell{\gd}$, and then
$\erad(\lf)\le2h\le r$ and \cref{prop:elliptic} with $C=0$ gives
$f_{n,r}(\lf)=1$; or $\lf=\Ad^*(g)\nf$ with $g\in\cG(\OO)$,
$\nf\in\uu_P^\perp\cap\dLat$ and $\norm{Q_\nf^{-1}}\le q^h$, and then
\cref{thm:regularunram} gives $f_{n,r}(\lf)=0$. Since
$0\le f_{n,r}\le1$, the claim follows.
\end{proof}

\begin{proof}[Proof of \cref{thm:uniform}]
Take $r=\lfloor n/2\rfloor$ and $h=\lfloor(n-5)/4\rfloor-m_{\cG}$, which is
$\ge1$ for $n\ge4m_{\cG}+9$. Then $n\ge2r$, and $3r-n\ge(n-3)/2$, so
$3r-n-2h\ge(n-3)/2-(n-5)/2+2m_{\cG}=2m_{\cG}+1$. By \cref{prop:explicitcell}
and \eqref{eq:badexplicit},
\[
\int_\dLat\abs{f_{n,r}-\one_{\Ell{\gd}}}\,d\nu_{\dLat}\le\nu_{\dLat}(\Bad_h)\le \Nss q^{-h}
\le \Nss q^{\,m_{\cG}-(n-8)/4},
\]
since $\lfloor(n-5)/4\rfloor\ge(n-8)/4$. The bound for $\scd{K_n}$ follows as
in the proof of \cref{thm:main}.
\end{proof}

\begin{remark}[Uniform localized bound]\label{rem:uniformlocalized}
In the setting of \cref{thm:uniform}, let $n\ge4m_{\cG}+9$ and let
$\cell\subset\dLat$ be a union of cosets of $\pp^m\dLat$ with
$m\le\lfloor n/2\rfloor$. Then the error term in \eqref{eq:localized} is at
most $\Nss\,q^{\,m_{\cG}+2-n/4}$. Indeed, as in the proof of
\cref{cor:localized}, it equals
$\int_\cell(f_{n,r}-\one_{\Ell{\gd}})\,d\nu_{\dLat}$ with
$r=\lfloor n/2\rfloor$, which is bounded by the integral in the proof of
\cref{thm:uniform}.
\end{remark}

\begin{proof}[Proof of \cref{cor:chevalley}]
By base change, a Chevalley group scheme over $\Z$ gives a reductive model of
$\bG_F$ over $\OO$ for every $F$, and $\Nss$ and $m_{\cG}$ depend only on its
root datum. For $\bG=\GL_d$ we have $\Nss=2^d-2$ and $m_{\cG}=1$
(\cref{sec:integraldata}), so that $4m_{\cG}+9=13$ and $m_{\cG}+2=3$.
\end{proof}

\subsection{Uniformity over a global field}\label{sec:global}
Recall the setting of \cref{thm:global}: $k$ is a global field, $\bG$ is a
connected reductive group over $k$ with a closed embedding $\bG\subset\GL_d$,
and for a finite place $v$ of $k$ we have $\Lat_v=\g\cap\Mat_d(\OO_v)$,
$K_{v,0}=\bG(k_v)\cap\GL_d(\OO_v)$ and
$K_{v,n}=\bG(k_v)\cap(1+\pp_v^n\Mat_d(\OO_v))$ for $n\ge1$. (The second
formula is not meant for $n=0$: for $\bG=\GL_1$ it would give
$\OO_v\setminus\{0\}$.) For $n\ge1$ these are the lattice
$\Lat^\natural$ and the tower $(K^\natural_n)$ of \cref{sec:models} for the
embedding $\bG_{k_v}\subset\GL_{d,k_v}$ over $F=k_v$.
If the schematic closure $\cG^\natural$ of \cref{sec:models} is a
reductive group scheme, then $\Lat^\natural$ and $K_n^\natural$ are the
lattice and the tower of \cref{thm:uniform}, for all $n\ge1$. Without this reductivity assumption, the constants of
\cref{thm:main} for $(K^\natural_n)$ depend on the embedding and on $F$.
Over a global field, only finitely many places require these
field-dependent constants; this gives \cref{thm:global}, with $m_{\bG}$ defined as follows. Let $\Gamma$ be
the finite image of the Galois group of $k$ on the based absolute root
datum of $\bG$, acting through the $*$-action (so that no Borel subgroup
over $k$ is needed when $\bG$ is not quasi-split). For every
subgroup $\Gamma'\subset\Gamma$, compute the invariant in \eqref{eq:mG} using
the $\Gamma'$-orbits on $\Delta$ and the $\Gamma'$-invariant cocharacters of
\eqref{eq:cocharchoice}; let $m_{\bG}$ be the maximum of these finitely many
invariants.

\begin{proof}[Proof of \cref{thm:global}]
Let $\cG$ be the schematic closure of $\bG$ in $\GL_d$ over the ring of
integers $\OO_k$ (or over a Dedekind ring of $S_0$-integers if $k$ is a
function field). It is a flat affine group scheme of finite type with
generic fibre $\bG$, hence smooth over a dense open subset of
$\Spec\OO_k$ by generic smoothness, with reductive fibres
over a dense open subset by \cite[Proposition~3.1.9]{MR3362641}; the
fibres there are connected by \cite[Proposition~3.1.12]{MR3362641} or
\cite[Proposition~A.8.2]{MR4520154}. Let $S$ be the complement,
together with the places omitted in the function field case. (Compare
\cite[Proposition~18.1.8]{MR4520154}, which proves in the same way that
$\bG(k_v)\cap\SL_d(\OO_v)$ is hyperspecial for almost all $v$ for an
embedding into $\SL_d$.)

For $v\notin S$, $\cG_v=\cG\otimes\OO_v$ is a reductive group scheme
and is the schematic closure of $\bG$ in $\GL_{d,\OO_v}$, since
schematic closure commutes with the flat base change $\OO_k\to\OO_v$.
The identifications of $\Lat_v$ and $K_{v,n}$ were noted before
\cref{lem:embedding}. Moreover $K_{v,0}=\bG(k_v)\cap\GL_d(\OO_v)=\cG_v(\OO_v)$,
since an $\OO_v$-point of $\GL_d$ whose generic fibre lies in $\bG$ factors
through the schematic closure $\cG_v$, and $\cG_v(\OO_v)$ is hyperspecial by
\cite[Theorem~9.9.3]{MR4520154}. Thus \cref{thm:uniform} applies to $\cG_v$.
Moreover,
the number of maximal semistandard parabolic subgroups of $\bG_{k_v}$ is at
most $N_{\bG_{\bar k}}$ by \cref{sec:parabolics}, since $\bG_{k_v}$ is quasi-split, while the types of the
standard maximal parabolic subgroups of $\bG_{k_v}$ are the orbits
of the image of its decomposition group in $\Gamma$. This is one of
the subgroups used to define $m_{\bG}$, so
$m_{\cG_v}\le m_{\bG}$. This gives
\eqref{eq:globalgood}.

For $v\in S$, \cref{lem:embedding} and
\cref{thm:main} give constants $C_v$, $n_v$ with
$\abs{\scd{K_{v,n}}-\elld[\bG_{k_v}]{\Lat_v^\vee}}\le C_vq_v^{-n/4}$ for
$n\ge n_v$, and \eqref{eq:globalall} holds with
$n_0=\max(\{4m_{\bG}+9\}\cup\{n_v:v\in S\})$ and
$C=\max(\{N_{\bG_{\bar k}}\}\cup\{C_vq_v^{-m_{\bG}-2}:v\in S\})$; if $S$ is
empty, these are $4m_{\bG}+9$ and $N_{\bG_{\bar k}}$.
\end{proof}

\begin{remark}[Sharpness of the exponent]\label{rem:sharpness}
The exponent $\eexp/4$ in \cref{thm:main} need not be optimal:
for $\GL_2$ and $\Lat=\Mat_2(\OO)$, where $\eexp=2$ and
$\elld{\dLat}=\frac12$ (\cref{prop:gl2}), $\frac12-\scd{K_n}$ is of exact order $q^{-2n}$ (\cref{prop:gl2exact}), while for $\PGL_2$
in characteristic $2$, where $\eexp=1$ and $\elld{\dLat}=1$, $1-\scd{K_n}$ is of exact
order $q^{-n/2}$ (\cref{prop:pgl2}). We do not know the optimal exponent in general.

The explicit bound \eqref{eq:explicitbound} is weaker still: for $\GL_2$ it
gives $2q^{3-n/4}$. We have kept the exponent $1/4$ there rather than
$\eexp/4$ because the implied constant in \cref{prop:geometric} depends on the
polynomials $\det \Bf_\lf$ through \cref{lem:sublevel} and we have not
made it explicit.
For $\GL_2$ with the trace pairing, the \bnets{} $\Bad_h$ of \cref{lem:integralatlas} are
contained in the set of $\lf\in\Mat_2(\OO)$ with $\val(\disc\lf)\ge2h$, where
$\disc\lf=(a-d)^2+4bc$; for odd residual characteristic this set has measure
$\ll q^{-2h}$. This is consistent with $\eexp=2$ in \cref{prop:geometric};
the loss in \eqref{eq:explicitbound} comes from using only the bound
$\Nss q^{-h}$. A further loss comes from the margins
\eqref{eq:margins}, which force $h\le n/4+O(1)$; we do not know how far they
can be relaxed.
\end{remark}

\subsection{The supercuspidal part of the Hecke algebra}\label{sec:squares}
The cellwise analysis of \cref{sec:proofsgeneral} also gives a lower bound
for the sum of the squares of the dimensions of the spaces of $K_n$-fixed
vectors in supercuspidal representations, i.e., the dimension of the supercuspidal part
of the Hecke algebra of $K_n$-biinvariant functions.

\begin{theorem}[Supercuspidal part of the Hecke algebra]\label{thm:squares}
Suppose that $Z(\bG)(F)$ is compact, and let $K$ be a compact open
subgroup of $G$ such that $K_n$ is a normal subgroup of $K$ for all
$n\gg1$, e.g.\ $K=K_1$. Then
\begin{equation}\label{eq:squares}
\sum_{\pi\in\Gsc}(\dim\pi^{K_n})^2\ge\bigl(\elld{\dLat}-O(q^{-\eexp n/4})\bigr)[K:K_n].
\end{equation}
\end{theorem}

For a smooth model as in \cref{lem:model}, one may take $K=\cG(\OO)$; then $[K:K_n]=\abs{\cG(k)}q^{(n-1)D}$.

\begin{proof}
Put $r=r_n$, $h=h_n$, and take $n$ large enough that
$\Gamma=K/K_n$ and $H=K_{n-r}/K_n$ are defined. The group $H$ is
abelian of order $q^{rD}$, with characters indexed by the cells in
$\dLat/\pp^r\dLat$. Let $\mathcal W$ be the cells contained in
$\Ell{\gd}\setminus\Bad_h$. By \cref{prop:density}, every irreducible
representation of $G$ containing one of their characters is
supercuspidal, and
\[
\frac{|\mathcal W|}{|H|}
\ge\elld{\dLat}-\nu_{\dLat}(\Bad_h).
\]

In $\C[H]$ let $e$ be the sum of the character idempotents for
$\mathcal W$, and call $\rho\in\Irr(\Gamma)$ good if $\rho(e)\ne0$.
Since $\C[\Gamma]$ is free of rank $[\Gamma:H]$ as a right
$\C[H]$-module,
\[
\dim\C[\Gamma]e=|\mathcal W|[\Gamma:H].
\]
This left ideal is contained in the sum of the matrix blocks of
$\C[\Gamma]$ corresponding to the good representations. Hence
\[
\sum_{\rho\text{ good}}(\dim\rho)^2
\ge|\mathcal W|[\Gamma:H]
\ge\bigl(\elld{\dLat}-\nu_{\dLat}(\Bad_h)\bigr)[K:K_n].
\]

Every good $\rho$, inflated to $K$, occurs in an irreducible
representation $\pi$ of $G$ by \cref{rem:Ktype}, and this $\pi$ is
supercuspidal because $\rho(e)\ne0$. Assign each good $\rho$ to one
such $\pi$. The distinct representations assigned to a given $\pi$
occur as a direct sum in $\pi^{K_n}$, so
\[
\sum_{\pi\in\Gsc}(\dim\pi^{K_n})^2
\ge\sum_{\rho\text{ good}}(\dim\rho)^2.
\]
The preceding bound and \cref{rem:parameters} give \eqref{eq:squares}.
\end{proof}

\section{The elliptic density}\label{sec:positivity}

In this section we prove strict positivity of the elliptic density,
completing the proof of \cref{thm:main}, and characterize when it equals
$1$ (\cref{thm:endpoint}).
We also establish the noncriticality statements used in \cref{sec:incidence} for the improved error exponent.
For \ordinary{} groups (\cref{def:ordinary}), in
particular whenever $\chr F\ne2$, regular torus-invariant forms
on elliptic maximal tori and on maximal tori of Levi subgroups supply
elliptic forms and noncritical points (\cref{lem:coroots}). In
characteristic $2$ the adjoint odd orthogonal factors, which are not
\ordinary{}, are treated separately.

The elliptic density $\elld{\dLat}$ is unchanged if $\dLat$ is scaled or conjugated
(\cref{rem:ball}), but it is otherwise a genuine invariant of the
lattice and not of $\bG$ alone; see \cref{prop:gl2,prop:iwahori} for $\GL_2$.

\subsection{Elementary properties}

\begin{remark}[Scaling, conjugation and positivity]\label{rem:ball}
$\Ell{\gd}$ is stable under multiplication by $F^\times$, so
$\elld{a\dLat}=\elld{\dLat}$ for $a\in F^\times$. Likewise, $\Ell{\gd}$ is
$\Ad^*(G)$-stable, so $\elld{\Ad^*(g)\dLat}=\elld{\dLat}$ for $g\in G$.
Moreover, $\elld{\dLat}>0$ if and only if $\Ell{\gd}\ne\emptyset$.
Indeed, scaling takes any elliptic form into $\dLat$, and openness of
$\Ell{\gd}$ (\cref{sec:incidencemaps}) then gives a nonempty open
subset of $\dLat\cap\Ell{\gd}$.
\end{remark}


\begin{lemma}[Full-density criterion]\label{lem:one}
$\elld{\dLat}=1$ if and only if $\Omega_M=\emptyset$ for every
$\bP\in\Pc$ (with the notation of \cref{sec:levi}), i.e., if and only if
every incidence map $p_P$ is everywhere critical. In particular this
condition does not depend on $\Lat$.
\end{lemma}

\begin{proof}
If $\Omega_M=\emptyset$ for all $\bP\in\Pc$ then every point of every
$\Inc_P$ is critical by \eqref{eq:levi}, so $\Par{\gd}=\Crv$ has measure zero
by \cref{thm:sard}, and $\elld{\dLat}=1$. Conversely, if
$\lf\in\Omega_M$ then $dp_P$ is invertible at $(\bP,\lf)$, so
$\Par{\gd}\supset p_P(\Inc_P)$ contains a nonempty open set, which after
scaling (\cref{rem:ball}) meets $\dLat$ in a set of positive measure.
\end{proof}

\subsection{Elliptic tori}
A maximal torus $\bT$ of $\bG$ (defined over
$F$) is \emph{elliptic} if its maximal split subtorus is the maximal
split central torus of $\bG$. Equivalently, $\bT$ is contained in no
proper parabolic subgroup of $\bG$.
Elliptic maximal tori exist: this is proved in \cite[Theorem~6.21]{MR1278263}
in characteristic zero and in \cite[\S2.4]{MR2214792} (an argument of
Prasad) in general; cf.~\cite[Theorem~5.3]{MR4617942}.

For $\lf\in\gd$ let $\g_\lf=\{X\in\g:\lf([X,Y])=0\ \forall Y\in\g\}$ be its infinitesimal stabilizer. 

\begin{lemma}[Elliptic forms from elliptic tori]\label{lem:toral}
Let $\bT$ be an elliptic maximal torus with Lie algebra $\liet$ and let
$\lf\in\gd$ be $\bT$-invariant with $\g_\lf=\liet$. Then $\lf\in\Ell{\gd}$.
\end{lemma}

\begin{proof}
Suppose on the contrary that $\lf|_{\uu_P}=0$ for a proper parabolic subgroup $\bP$.
Consider the incidence variety $\bInc_P$ of \cref{sec:incidencemaps},
and work after base change to $\bar F$
at the point $(\bP,\lf)$. By \eqref{eq:dp}, the
image of the differential of $p_P:\bInc_P\to\Lie(\bG)^*$ at this point is
$\uu_P^\perp+\ad^*(\uu_P^-)\lf$. Its annihilator in $\g$ is
$\{X\in\uu_P:\lf([X,\uu_P^-])=0\}=\uu_P\cap\g_\lf$, since
$\lf([\uu_P,\Lie(\bP)])\subset\lf(\uu_P)=0$, and
$\uu_P\cap\g_\lf=\uu_P\cap\liet=0$. Thus the differential is bijective, and hence $(\bP,\lf)$ is an
isolated point of the fiber over $\lf$. The connected group
$\bT_{\bar F}$ fixes $\lf$, hence acts on this fiber and fixes its
isolated points. Therefore $\bT$ normalizes $\bP$, so $\bT\subset\bP$
since parabolic subgroups are self-normalizing. This contradicts the ellipticity of $\bT$.
\end{proof}

For a maximal torus $\bT$ and an absolute root $\alpha$ write
$H_\alpha=d\alpha^\vee(1)\in\liet_{\bar F}$.

\begin{definition}[\Ordinary{} groups]\label{def:ordinary}
We call $\bG$ \emph{\ordinary{}} if
\begin{equation}\label{eq:corootcondition}
H_\alpha\ne0\qquad\text{for every absolute root }\alpha.
\end{equation}
The condition is imposed over $\bar F$ and is independent of the
choice of $\bT$.
\end{definition}

This is the condition called coroot-unramified in \cite[Definition~2.4]{Chi26}.
In characteristic $p>0$, this is equivalent to requiring
$\alpha^\vee\notin pX_*(\bT_{\bar F})$ for every $\alpha$.
Geometrically, it says that the regular semisimple locus in $\gd$ is
nonempty. Indeed, over $\bar F$, identify $\liet^*$ with the forms
vanishing on all root spaces. As observed by Goresky--Kottwitz--MacPherson
\cite[\S6.4]{MR2209851} (see \cite[Lemma~2.2(1) and Corollary~2.3]{Chi26}
for a proof), for $\lf\in\liet^*$ one has
\begin{equation}\label{eq:toralstabilizer}
\g_\lf=\liet\oplus
\bigoplus_{\lf(H_\alpha)=0}\g_\alpha.
\end{equation}
Thus $\g_\lf=\liet$ exactly when every $\lf(H_\alpha)$ is nonzero;
such forms exist exactly under \eqref{eq:corootcondition}.
This is the coadjoint counterpart of the usual root criterion for
regular semisimple elements of $\g$; for the classical adjoint
discussion, see Springer--Steinberg \cite[\S II.3.17]{MR0268192}. For Jordan
decompositions in $\gd$ in arbitrary characteristic, see \cite{ST26}; by
\cite[Theorem~1.1(iv)]{ST26}, the closure of the set of semisimple elements of
$\gd$ is all of $\gd$ unless $\chr F=2$ and $\bG_{\bar F}$ has a direct factor
$\SO_{2d+1}$.

Over $F$ this condition has the following consequence.

\begin{lemma}[Regular torus-invariant forms]\label{lem:coroots}
Assume that $\bG$ is \ordinary{}. Then for every maximal torus $\bT$
of $\bG$ there is a nonempty open subset of the space of $\bT$-invariant
$\lf\in\gd$ with $\g_\lf=\liet$. Consequently
$\Ell{\gd}\ne\emptyset$, and $\Omega_M\ne\emptyset$ for every proper
parabolic subgroup $\bP=\bM\bU$ of $\bG$.
\end{lemma}

\begin{proof}
The projection to the zero weight space identifies $\liet^*$ with the
space of $\bT$-invariant elements of $\gd$. Since each $H_\alpha\ne0$,
$\{\lf\in\liet^*(F):\lf(H_\alpha)=0\}$ is a proper $F$-subspace
and the condition $\lf(H_\alpha)\ne0$ for all $\alpha$ defines a
nonempty open subset of $\liet^*(F)$ (a finite union of proper subspaces
does not cover a vector space over an infinite field).
By \eqref{eq:toralstabilizer}, these forms satisfy $\g_\lf=\liet$.
Taking $\bT$ elliptic (whose existence was recalled at the start of
this subsection), \cref{lem:toral} gives $\lf\in\Ell{\gd}$. Taking
$\bT\subset\bM$, the pairing $\Bf_{\lf|_\m}$ of
\eqref{eq:levipairing} is, over the splitting field, diagonal with
respect to root vectors with entries $\pm\lf(H_\alpha)\ne0$ for the
roots $\alpha$ of $\uu$, hence non-degenerate. Indeed, $\lf$ is
$\bT$-invariant, so it vanishes on the root spaces; hence
$\lf([X_\alpha,X_{-\beta}])=0$ for $\beta\ne\alpha$, while
$[X_\alpha,X_{-\alpha}]=\pm H_\alpha$ for root vectors from a Chevalley
basis, an identity over $\Z$ valid in every characteristic.
\end{proof}

\subsection{Characteristic two}

Chi \cite[Remark~2.5]{Chi26} notes that the only possible geometric
obstructions to being \ordinary{} are odd orthogonal factors
in characteristic $2$. For our application we need the corresponding
direct-product decomposition over $F$, which we establish below.

\begin{lemma}[\Oddorth{} and \ordinary{} parts]\label{lem:split}
If $\chr F\ne2$ then $\bG$ is \ordinary{}. If
$\chr F=2$ there is a unique decomposition $\bG=\Godd\times\Gord$ over $F$
such that $\Gord$ is \ordinary{} and $(\Godd)_{\bar F}$ is a
product of adjoint groups of type $\mathrm B_d$, $d\ge1$ (where
$\mathrm B_1=\mathrm A_1$); either factor may be trivial. If $\chr F\ne2$
we put $\Godd=1$ and $\Gord=\bG$. We call $\Godd$ and $\Gord$ the
\emph{\oddorth{} part} and the \emph{\ordinary{} part} of $\bG$. In
particular, $\bG$ is \ordinary{} if and only if $\Godd=1$.
\end{lemma}

\begin{proof}
If $\alpha^\vee\in pX_*$ then $p$ divides $\langle\alpha,\alpha^\vee\rangle=2$,
so $p=2$. Let $p=2$. Since $X_*$ is stable under the Weyl group and every root is
conjugate to a simple root, it suffices to treat simple roots: let
$\alpha$ be a simple root (for some ordering)
of an irreducible component $\Phi_i$ of the absolute root system with
$\alpha^\vee\in2X_*$. Then all the entries $\langle\beta,\alpha^\vee\rangle$
of the column of $\alpha^\vee$ in the Cartan matrix are even, and
inspection of the Cartan matrices shows that this happens only for the
short simple root of $\mathrm B_d$, $d\ge1$. (Dually, simply connected
$\mathrm C_d$ is the only reduced irreducible root datum with a root
divisible in the character lattice; see \cite[proof of Theorem~2.3.10 and
Remark~5.3.11]{MR3362817}. The same reduction to adjoint factors of type
$\mathrm B_d$ appears in \cite[\S4]{ST26}.)

We show that this component splits off as an adjoint factor.
Conceptually, the divisibility of $\alpha^\vee$ forces the cocharacter
lattice of the component to be the full coweight lattice, i.e., the
corresponding normal subgroup is adjoint, and an adjoint factor always
splits off, because the projection of $X_*$ to the span of its coroots
lies in the coweight lattice. In detail:
Let $V_i\subset X_*\otimes\Q$ be the span of the coroots of $\Phi_i$, and let $P_i^\vee$ be the
coweight lattice of $\Phi_i$. In the standard
coordinates, $P_i^\vee=\Z^d$ and $\tfrac12\alpha^\vee=\pm e_j$ for some
$j$; since $X_*$ is stable under the Weyl group, $X_*\cap V_i\supset P_i^\vee$,
hence $X_*\cap V_i=P_i^\vee$. If $V_i'$ is the common kernel of the roots
of $\Phi_i$, then $X_*\otimes\Q=V_i\oplus V_i'$, and the projection of
$X_*$ to $V_i$ pairs integrally with the roots of $\Phi_i$, hence lies
in $P_i^\vee$. Therefore $X_*=P_i^\vee\oplus(X_*\cap V_i')$, and
dually the character lattice and the roots decompose accordingly. Thus
the root datum of $\bG_{\bar F}$ is the direct sum of the adjoint root
datum of type $\mathrm B_d$ and a complementary root datum, and
$\bG_{\bar F}=\bH_i\times\bG_i'$ accordingly.

It remains to descend the product of these factors to $F$.
Let $\Godd$ be the product of the normal subgroups $\bH_i$ over all components $\Phi_i$ containing
a coroot divisible by $2$. Since $\bG$ splits over the separable closure $F^{\mathrm{sep}}$ of $F$, the $\bH_i$ are defined over $F^{\mathrm{sep}}$, and this set of components is stable under $\Gal(F^{\mathrm{sep}}/F)$, so $\Godd$ is defined over $F$, and so is
$\Gord=Z_\bG(\Godd)$, which satisfies $\bG=\Godd\times\Gord$ over $\bar F$
because $\Godd$ has trivial center. Hence $\bG=\Godd\times\Gord$ over $F$,
and no coroot of $\Gord$ is divisible by $2$.

For uniqueness, let $\bG=\bH'\times\bG'$ be another such decomposition,
and work over $\bar F$ with a maximal torus $\bT=\bT_{\bH'}\times\bT_{\bG'}$,
so that $X_*=X_*(\bT_{\bH'})\oplus X_*(\bT_{\bG'})$. Each component
$\Phi_i$ is a component of the root system of $\bH'$ or of $\bG'$. Every
component of $\bH'$ has a coroot divisible by $2$ in $X_*(\bT_{\bH'})$,
hence in $X_*$, since the coroots corresponding to the short roots of an adjoint group of type
$\mathrm B_d$ are divisible by $2$ in its cocharacter lattice. Conversely,
if a coroot $\alpha^\vee$ of $\bG'$ satisfies $\alpha^\vee=2\mu$ with
$\mu\in X_*$, then the projection of $\mu$ to $X_*(\bT_{\bG'})$ shows that
$\alpha^\vee\in2X_*(\bT_{\bG'})$, which is impossible since $\bG'$ is
\ordinary{}. Hence the components of $\bH'$ are exactly those
containing a coroot divisible by $2$, and $\bH'=\Godd$, both being
generated by the root groups of these components. Finally
$\bG'=Z_\bG(\bH')=\Gord$, since $\bH'$ has trivial center.
\end{proof}

\begin{lemma}[Criticality for adjoint odd orthogonal factors]\label{lem:oddorthogonal}
Let $\chr F=2$ and let $\bH$ be an $F$-group such that
$\bH_{\bar F}$ is a product of adjoint groups of type $\mathrm B_d$, $d\ge1$.
Then every incidence map $p_P$ of $\bH$ is everywhere critical, so that
$\Par{\lieh^*}$ has measure zero, where $\lieh=\Lie(\bH)(F)$.
\end{lemma}

\begin{proof}
Criticality can be checked over $\bar F$, so we may assume that
$\bH$ is a product of split adjoint groups of type $\mathrm B_d$.
For a product, a proper parabolic is proper in at least one factor,
and $\Bf_\lf$ is the direct sum of the pairings for the factors.
Moreover, a maximal proper parabolic subgroup over $F$ need not remain
maximal over $\bar F$. If $\bP\subset\bQ$ with $\bQ$ maximal proper, then
$\uu_Q\subset\uu_P$ gives a map $\Inc_P\to\Inc_Q$ through which
$p_P$ factors, so $p_P$ is everywhere critical if $p_Q$ is. Thus it suffices to treat maximal proper parabolics
in a single split factor. We will exhibit a degenerate weight block
of $\Bf_\lf$.

We first describe the Lie algebra. Consider a split factor
$\SO(V)$, where $V$ is a vector space of dimension $2d+1$ over $\bar F$
with a non-degenerate quadratic form $Q$, whose polar form $\beta$ has
a one-dimensional radical on which $Q\ne0$. Here $\SO(V)$ denotes the
special orthogonal group of $(V,Q)$ in the sense of
\cite[Definition~C.2.10]{MR3362641}; in characteristic $2$ and odd
dimension it is the smooth connected adjoint group of type $\mathrm B_d$,
and $\mathrm O(Q)=\SO(V)\times\mu_2$ (\cite[Theorem~C.2.11]{MR3362641} and the remark after \cite[Definition~C.2.10]{MR3362641}). The map
\[
\bigwedge^2V\to\mathfrak{so}(V),\qquad
u\wedge v\mapsto\bigl(z\mapsto\beta(v,z)u+\beta(u,z)v\bigr)
\]
is an $\SO(V)$-equivariant isomorphism (cf.~\cite[\S\S3--4]{MR2558787}).
Indeed, every endomorphism $X$ in the image satisfies $\beta(z,Xz)=0$
for all $z$ and $\tr X=0$, so lies in $\mathfrak{so}(V)$.
To see injectivity, decompose $V=\rad\beta\oplus W$ with $\beta|_W$
non-degenerate. Then $\bigwedge^2W$ injects into $\End W$ and
$\rad\beta\wedge W$ into $\Hom(W,\rad\beta)$. Finally, both sides
have dimension $d(2d+1)$.

We now compute the relevant weight block. Let $\bP_i$ be the stabilizer of an isotropic
subspace $E$ of dimension $i$, $1\le i\le d$. Write
$V=E\oplus V_0\oplus E'$ with $E'$ isotropic, $\beta(e_a,f_b)=\delta_{ab}$
for bases $(e_a)$ of $E$ and $(f_b)$ of $E'$, and
$V_0=(E\oplus E')^\perp$ of dimension $2(d-i)+1$. For the cocharacter
acting by $t,1,t^{-1}$ on the three summands, the weight spaces of
$\uu^\pm$ are
\[
\begin{aligned}
\uu_1&=E\wedge V_0, & \uu_2&=\bigwedge^2E,\\
\uu_{-1}&=E'\wedge V_0, & \uu_{-2}&=\bigwedge^2E'.
\end{aligned}
\]
Computing commutators of the endomorphisms above gives, for $v,w\in V_0$,
\[
[e_a\wedge v,f_b\wedge w]=\delta_{ab}\,v\wedge w+\beta_0(v,w)\,e_a\wedge f_b,
\qquad\beta_0=\beta|_{V_0}.
\]
Let $\lf\in\m^*$ and put $a_0(v,w)=\lf(v\wedge w)$,
$C_{ab}=\lf(e_a\wedge f_b)$. Then on $\uu_1\times\uu_{-1}$,
\[
\Bf_\lf(e_a\wedge v,f_b\wedge w)=\delta_{ab}a_0(v,w)+C_{ab}\beta_0(v,w).
\]
Choose a basis of $V_0$, and let $A$ and $B$ be the matrices of
$a_0$ and $\beta_0$. The displayed block has matrix
$I_i\otimes A+C\otimes B$. Over $\bar F$, choose an eigenvector
$x\ne0$ of $C$, with eigenvalue $c$. Since $A+cB$ is alternating
of odd size, there is a nonzero vector $v$ with $(A+cB)v=0$.
Hence
\[
(I_i\otimes A+C\otimes B)(x\otimes v)
=x\otimes(A+cB)v=0.
\]
Thus this block is degenerate for every $\lf$. Since $\Bf_\lf$ is
block diagonal with respect to the weights (as in the proof of
\cref{lem:levi}), $\Omega_M=\emptyset$.
By \eqref{eq:levi}, $p_{P_i}$ is everywhere critical.
The measure-zero assertion follows from \cref{thm:sard}.
\end{proof}

Note that if $\bG=\bG_1\times\bG_2$, with $\g_i=\Lie(\bG_i)(F)$, so that
$\gd=\g_1^*\oplus\g_2^*$, then the proper parabolic subgroups of $\bG$ are
the products $\bP_1\times\bP_2$ of parabolic subgroups $\bP_i$ of $\bG_i$, not
both equal to $\bG_i$, and the maximal ones are $\bP_1\times\bG_2$ and
$\bG_1\times\bP_2$ with $\bP_i$ maximal proper. Since the Lie algebra
of the unipotent radical of
$\bP_1\times\bP_2$ is $\uu_{P_1}\oplus\uu_{P_2}$, it follows that
\[
\Ell{\gd}=\Ell{(\g_1^*)}\times\Ell{(\g_2^*)}.
\]
Moreover, if $\bP=\bP_1\times\bG_2$, then $\uu_P=\uu_{P_1}$ and
$\uu_P^-=\uu_{P_1}^-$ lie in $\g_1$, so that $Q_{\lf,P}=Q_{\lf|_{\g_1},P_1}$.
Hence $\Omega_M=\Omega_{M_1}\times\g_2^*$ for $\bM=\bM_1\times\bG_2$, and by
\eqref{eq:levi} $p_P$ is everywhere critical if and only if $p_{P_1}$ is.

\begin{proposition}[Peculiar groups]\label{prop:peculiar}
The group $\bG$ is \peculiar{} (\cref{def:peculiar}) if and only if
$\chr F=2$ and $\bG$ has a direct factor over $F$ which is isotropic over $F$
and whose base change to $\bar F$ is a product of adjoint groups of type
$\mathrm B_d$, $d\ge1$; in other words, if and only if its \oddorth{} part
$\Godd$ (\cref{lem:split}) is isotropic over $F$.
In this case, for $\bP\in\Pc$ the map $p_P$ is everywhere critical if and only if
$\bP=\bP_{\mathrm{odd}}\times\Gord$ with $\bP_{\mathrm{odd}}$ a maximal proper parabolic subgroup of
$\Godd$; in particular, if $\bG$ is \peculiar{} and $F$-simple, then every $p_P$
is everywhere critical.
\end{proposition}

\begin{proof}
By \eqref{eq:levi}, $p_P$ is everywhere critical if and only if
$\Omega_M=\emptyset$. If $\bG$ is \ordinary{}, in particular if
$\chr F\ne2$ (\cref{lem:split}), then $\Omega_M\ne\emptyset$ for every
$\bP\in\Pc$ by \cref{lem:coroots}. Let $\chr F=2$ and write
$\bG=\Godd\times\Gord$ as in \cref{lem:split}. By the note above, every
maximal proper parabolic subgroup of $\bG$ is of the form
$\bP_{\mathrm{odd}}\times\Gord$ or $\Godd\times\bP_{\mathrm{ord}}$ with
$\bP_{\mathrm{odd}}$, $\bP_{\mathrm{ord}}$ maximal proper. In the first case
$p_P$ is everywhere critical by \cref{lem:oddorthogonal} and the note. In the
second, $\Omega_M=\godd^*\times\Omega_{M_{\mathrm{ord}}}$ is nonempty by
\cref{lem:coroots}, since $\Gord$ is \ordinary{}. This proves the second
assertion, and shows that $\bG$ is \peculiar{} if and only if $\Godd$ has a
proper parabolic subgroup, i.e., is isotropic. Finally, a direct factor of $\bG$ as in
the first formulation is a direct factor of $\Godd$, since the coroots of the
short roots of an adjoint group of type $\mathrm B_d$ are divisible by $2$ in
its cocharacter lattice; and $\Godd$ itself
is such a factor when it is isotropic.
\end{proof}

\begin{remark}[Coroots and incidence maps]\label{rem:corootincidence}
Being \ordinary{} implies that $\bG$ is not \peculiar{}, but the converse is not true.
For example, let $\chr F=2$. Then $\bH=\PGL_1(D)$ for a quaternion division
algebra $D$ over $F$ is anisotropic but not \ordinary{}; for $\bG=\PGL_1(D)\times\GL_2$ one
has $\dG{}=1$ and $\eexp=2$.
\end{remark}

\subsection{Positivity and full density}

\begin{theorem}[Positivity and full density]\label{thm:endpoint}
\begin{enumerate}[label=\textup{(\arabic*)},ref=\arabic*]
\item\label{item:end-pos} $\Ell{\gd}\ne\emptyset$, hence $\elld{\dLat}>0$ for every lattice
$\dLat\subset\gd$.
\item\label{item:end-factor} Let $\bG=\Godd\times\Gord$ be as in \cref{lem:split}. Put $\gord=\Lie(\Gord)(F)$, and let
$\dLat_{\mathrm{ord}}$ be the projection of $\dLat$ to $\gord^*$, with probability Haar measure
$\nu_{\dLat_{\mathrm{ord}}}$. Then
\begin{equation}\label{eq:factor}
\elld{\dLat}=\nu_{\dLat_{\mathrm{ord}}}(\dLat_{\mathrm{ord}}\cap\Ell{(\gord^*)}).
\end{equation}
\item\label{item:end-full} The following are equivalent: $\elld{\dLat}=1$ for some (or every) lattice
$\dLat\subset\gd$; $\Par{\gd}$ is a null set; $\Gord$ is anisotropic modulo its center.
In particular, if $\chr F\ne2$, they hold if and only if $\bG$ is anisotropic modulo its center.
\end{enumerate}
\end{theorem}

\begin{proof}
By the note before \cref{prop:peculiar},
$\Ell{\gd}=\Ell{(\godd^*)}\times\Ell{(\gord^*)}$, where
$\godd=\Lie(\Godd)(F)$. Here $\Ell{(\godd^*)}$ is the
complement of a null set by \cref{lem:oddorthogonal}, in particular
nonempty, and $\Ell{(\gord^*)}\ne\emptyset$ by \cref{lem:coroots}. Hence
$\Ell{\gd}\ne\emptyset$ and $\elld{\dLat}>0$ by \cref{rem:ball}. Since
$\Par{(\godd^*)}\times\gord^*$ is null, $\one_{\Ell{\gd}}$ coincides almost everywhere
with the pull-back of $\one_{\Ell{(\gord^*)}}$, and the projection $\dLat\to \dLat_{\mathrm{ord}}$
maps $\nu_{\dLat}$ to $\nu_{\dLat_{\mathrm{ord}}}$; this gives \eqref{eq:factor}.

For \cref{item:end-full}, note that $\elld{\dLat}=1$ if and only if
$\dLat\cap\Par{\gd}$ is a null set. Since $\Par{\gd}$ is stable under
$F^\times$ and $\gd=\bigcup_{k\ge0}\varpi^{-k}\dLat$, this holds for one
lattice if and only if $\Par{\gd}$ is a null set, and then it holds for
every lattice. If $\Gord$ has no proper parabolic subgroup then
$\Ell{(\gord^*)}=\gord^*$ and $\elld{\dLat}=1$ by \eqref{eq:factor}.
Otherwise \cref{lem:coroots} gives $\Omega_{M_{\mathrm{ord}}}\ne\emptyset$ for a
maximal proper parabolic subgroup of $\Gord$ with Levi subgroup
$\bM_{\mathrm{ord}}$, so that
$\nu_{\dLat_{\mathrm{ord}}}(\dLat_{\mathrm{ord}}\cap\Ell{(\gord^*)})<1$ by
\cref{lem:one} applied to $\Gord$, and $\elld{\dLat}<1$ by
\eqref{eq:factor}. Hence these conditions hold if and only if $\Gord$ has no
proper parabolic subgroup, i.e., $\Gord$ is anisotropic modulo its
center. If $\chr F\ne2$ then $\bG=\Gord$.
\end{proof}

This completes the proof of \cref{thm:main}.

\appendix

\section{A quantitative Sard theorem}\label{sec:sard}

We prove a quantitative version of Sard's theorem on the image of critical points as well as a quantitative
inverse function theorem for analytic maps between manifolds of the same dimension, in every characteristic.
The proof provides \bnets{} around the critical values.
Finally, we give an elementary polynomial sublevel estimate.

Sard's theorem in this analytic setting was proved in
\cite[Theorem~2.5]{MR3803728} (in the equal-dimensional case considered
here), and an ultrametric inverse function theorem
in \cite[Theorem~2.4]{MR3057121}. The elementary estimates below also
apply when the differential is identically singular. Our quantitative
version is in the spirit of Oesterl\'e's bounds \cite{MR656627}.

\subsection{The local image estimate and inverse theorem}
Endow $F^D$ with the maximum norm and $\OO^D$ with probability Haar
measure. Let $f:\OO^D\to\OO^D$ be given by power series with integral
coefficients converging on $\OO^D$: their coefficients tend to zero
as the total degree tends to infinity. Expanding the monomials at
$x\in\OO^D$ gives, in every characteristic,
\begin{equation}\label{eq:taylor}
f(x+z)=f(x)+df_x(z)+R_x(z),\qquad
R_x(z)=\sum_{\abs\beta\ge2}b_\beta(x)z^\beta,
\end{equation}
with $df_x\in M_D(\OO)$ and $b_\beta(x)\in\OO^D$.
The expansions converge uniformly. Thus
\begin{equation}\label{eq:remainder}
\norm{R_x(z)}\le\norm z^2,\qquad
\norm{R_x(z)-R_x(z')}\le
\max(\norm z,\norm{z'})\norm{z-z'}.
\end{equation}
Indeed, each difference of monomials of degree at least two contains
a factor $z_i-z_i'$ and at least one further coordinate factor.
Also $J(x)=\det df_x$ is a convergent integral power series, hence
$1$-Lipschitz. Consequently, for $h\ge1$,
\begin{equation}\label{eq:Sh}
S_h(f)=\{x\in\OO^D:\abs{\det df_x}\le q^{-h}\}
\end{equation}
is a union of cosets of $\pp^h\OO^D$.

The second assertion of the next lemma is a quantitative form of the
classical multivariable Newton lemma (cf.\ \cite[\S3, Lemma~2]{MR0207700} and
\cite[Theorem~2.2.1 and Corollary~2.2.1]{MR1743467}).

\begin{lemma}[Image estimate and effective inverse theorem]\label{lem:tube}
For $h\ge1$,
\begin{equation}\label{eq:tube}
\vol\bigl(f(S_h(f))+\pp^{2h}\OO^D\bigr)
\le q^{-h}\vol(S_h(f))\le q^{-h}.
\end{equation}
For every $x\notin S_h(f)$,
\begin{equation}\label{eq:inverseball}
f(x+\pp^h\OO^D)\supset f(x)+\pp^{2h}\OO^D.
\end{equation}
\end{lemma}

\begin{proof}
Fix $x\in S_h(f)$. The remainder estimate \eqref{eq:remainder} gives
\[
f(x+\pp^h\OO^D)+\pp^{2h}\OO^D
\subset f(x)+df_x(\pp^h\OO^D)+\pp^{2h}\OO^D.
\]
If the elementary divisors of $df_x$ have valuations
$s_1,\dots,s_D\in\Z_{\ge0}\cup\{\infty\}$, the set on the right has volume
\[
q^{-Dh-\sum_j\min(s_j,h)}\le q^{-(D+1)h},
\]
since $\sum_js_j\ge h$ implies $\sum_j\min(s_j,h)\ge h$.
There are $q^{Dh}\vol(S_h(f))$ cosets of $\pp^h\OO^D$ in $S_h(f)$;
summing proves \eqref{eq:tube}.

Now let $x\notin S_h(f)$. Cramer's rule gives
$\norm{(df_x)^{-1}}\le\abs{\det df_x}^{-1}\le q^{h-1}$.
For $\delta\in\pp^{2h}\OO^D$, the map
\[
z\longmapsto(df_x)^{-1}(\delta-R_x(z))
\]
maps $\pp^h\OO^D$ to itself and, by \eqref{eq:remainder}, has
Lipschitz constant at most $q^{-1}$. Its fixed point satisfies
$f(x+z)=f(x)+\delta$, proving \eqref{eq:inverseball}.
\end{proof}

The Frobenius map $f(x)=x^p$ in characteristic $p$ has identically
zero differential, so $S_h(f)=\OO$ for every $h$. Its image satisfies
\[
\vol(f(\OO)+\pp^{2h})=q^{\lceil2h/p\rceil-2h}\le q^{-h},
\]
because $x^p\equiv y^p\pmod{\pp^{2h}}$ if and only if
$x\equiv y\pmod{\pp^{\lceil2h/p\rceil}}$. Equality holds for $p=2$.
Thus \cref{lem:tube} applies even when the differential vanishes
identically, and its bound is attained by a polynomial map in
characteristic $2$. Letting $h\to\infty$ also shows that the image
itself is null, in agreement with \cref{thm:sard}.

The lemma holds equally for a convergent integral power series map
between lattices in $D$-dimensional vector spaces, with probability
Haar measures and determinants taken in lattice bases. The conditions
and conclusions are independent of the chosen bases.

\subsection{Finite integral atlases and \bnets{}}\label{sec:bnets}
Let $X$ be an analytic manifold of dimension $D$, let $V$ be an
$F$-vector space of dimension $D$, and let $f:X\to V$ be proper and
analytic. Write $\Vf=f(X)$ and $\Crv=f(\Crit f)$; both are closed.
The inverse function theorem gives $\partial\Vf\subset\Crv$.
Fix a lattice $\dLat\subset V$, let $\nu_{\dLat}$ be the probability Haar measure on $\dLat$, and put
$\XL=f^{-1}(\dLat)$, a compact open subset of $X$.

An \emph{\iatlas{}} of $f$ over $\dLat$ is a finite family of analytic
isomorphisms $\psi_i:\Lat_i\to U_i$ onto open subsets of $\XL$
covering $\XL$, where each $\Lat_i$ is a lattice in a
$D$-dimensional vector space and
$f_i=f\circ\psi_i:\Lat_i\to\dLat$ is a convergent power series with
coefficients in $\dLat$. We use probability Haar measure on each
$\Lat_i$. The empty atlas is allowed when $\XL=\emptyset$.

Such atlases exist: around each $x\in\XL$ choose local coordinates
centered at $x$. The constant term of $f$ lies in $\dLat$; restricting
to a sufficiently small lattice makes all the other coefficients
integral and the series convergent on that lattice
(cf.\ \cite[Lemma~7.4.3(i)]{MR1743467}). Compactness then gives a finite
cover. For polynomial coordinate maps this is just translation and
rescaling of finitely many coefficients.

\begin{theorem}[Quantitative Sard theorem]\label{thm:sard}
Fix an \iatlas{} $(\psi_i:\Lat_i\to U_i)_{1\le i\le N}$, put
$S_{i,h}=S_h(f_i)$, and define the \emph{\bnets{}}
\begin{equation}\label{eq:badf}
\Bad_h=\bigcup_{i=1}^N\bigl(f_i(S_{i,h})+\pp^{2h}\dLat\bigr)
\subset\dLat,\qquad h\ge1.
\end{equation}
These sets decrease, are unions of cosets of $\pp^{2h}\dLat$, and
every such coset outside $\Bad_h$ is either contained in $\Vf$ or
disjoint from $\Vf$. Moreover,
\begin{equation}\label{eq:sardchart}
\abs{\det df_i(y)}>q^{-h}
\quad\text{if }f_i(y)\in\dLat\setminus\Bad_h,
\end{equation}
\begin{equation}\label{eq:sardmeasure}
\nu_{\dLat}(\Bad_h)\le q^{-h}\sum_{i=1}^N\vol(S_{i,h})\le Nq^{-h},
\end{equation}
and $\bigcap_{h\ge1}\Bad_h=\Crv\cap\dLat$. In particular,
$\Crv$ is a null set.
\end{theorem}

\begin{proof}
The nesting and coset assertions follow from the definition, as does
\eqref{eq:sardchart}. If a coset outside $\Bad_h$ meets $\Vf$, write
a point of intersection as $f_i(y)$. Then $y\notin S_{i,h}$, and
\eqref{eq:inverseball} shows that the entire coset lies in
$f_i(\Lat_i)\subset\Vf$. Summing \eqref{eq:tube} over the charts
gives \eqref{eq:sardmeasure}.

Every critical point in $\XL$ lies in a chart and has zero chart
Jacobian, so $\Crv\cap\dLat\subset\Bad_h$ for all $h$.
Conversely, if $\lf\in\bigcap_h\Bad_h$, finiteness of the atlas
gives an index $i$ and arbitrarily large $h$ with
$\lf=f_i(y_h)+z_h$, where $y_h\in S_{i,h}$ and
$z_h\in\pp^{2h}\dLat$. Along a convergent subsequence in the compact
lattice $\Lat_i$, the limit $y$ satisfies
$\det df_i(y)=0$ and $f_i(y)=\lf$. Hence $\lf\in\Crv$.
The measure bound makes $\Crv\cap\dLat$ null. Applying this to
$\pp^{-k}\dLat$, $k\ge0$, shows that $\Crv$ is a null set.
\end{proof}

Two useful consequences are
\begin{equation}\label{eq:bntube}
(\Crv+\pp^r\dLat)\cap\dLat
\subset\Bad_{\lfloor r/2\rfloor}\qquad(r\ge2),
\end{equation}
since the nets contain $\Crv\cap\dLat$ and are unions of the
corresponding cosets, and, for every compact $\cell\subset\dLat$,
\begin{equation}\label{eq:bncompact}
\cell\cap\Crv=\emptyset\implies
\cell\cap\Bad_h=\emptyset\quad\text{for }h\gg1.
\end{equation}
The latter follows by compactness from the nesting and intersection
assertions. The estimate \eqref{eq:sardmeasure} retains the sizes of
the source sublevel sets; together with the next lemma, this yields
the improved rate in \cref{sec:incidence}.

\subsection{Polynomial sublevel sets}\label{sec:sublevel}

The following estimate is standard.
For $F=\Q_\ell$ it is stated without proof in \cite[\S3.2, (57)]{MR644559}.
In positive characteristic, see \cite[Theorem~1.4]{GH26} for a slightly weaker bound
in a much more general setting.

\begin{lemma}[Polynomial sublevel estimate, cf.\ {\cite[\S3.2, (57)]{MR644559}}]\label{lem:sublevel}
Let $f\in F[x_1,\dots,x_m]$ be a nonzero polynomial of degree $d\ge1$.
Then there is a constant $C_f$ such that
\[
\vol\{x\in\OO^m:\abs{f(x)}\le\epsilon\}\le C_f\epsilon^{1/d}\qquad(\epsilon>0).
\]
\end{lemma}

\begin{proof}
For completeness we provide the details.
Let $f_d$ be the homogeneous part of degree $d$ of $f$ and choose a
primitive vector $v\in\OO^m$ with $c=f_d(v)\ne0$. Write
$\OO^m=\OO v\oplus W$ for a direct summand $W$, with product Haar
measure. For $w\in W$, the polynomial $t\mapsto f(w+tv)$ has degree
$d$ and leading coefficient $c$, so it factors as
$c\prod_{j=1}^d(t-\alpha_j)$ over an algebraic closure, with the
extended absolute value. If $\abs{f(w+tv)}\le\epsilon$, then
$\abs{t-\alpha_j}\le\rho=(\epsilon/\abs c)^{1/d}$ for some $j$.
Each such set of $t\in F$ is empty or a ball of measure at most
$\rho$. Thus the set of $t\in\OO$ under consideration has measure
at most $d\rho$, uniformly in $w$. Fubini's theorem gives
$C_f=d\abs c^{-1/d}$.
\end{proof}


\section{Technical details for the Jacquet argument}\label{sec:jacquetdetails}

This appendix is part of the proof of \cref{thm:regular,thm:regularunram}, and hence of
\cref{thm:main}; it supplies the technical verifications used in
\cref{sec:regular} to prove that a representation containing the cell
character $\theta_J$ has a nonzero Jacquet module $V_U$.
We verify that the enlarged group $H$ and its characters are well defined,
that an extension of $\theta_J$ can be conjugated to one trivial on $S$,
and that the successive averaging maps used in the proof are injective.
The final subsection compares characters in different charts,
completing the passage from the product chart to the original cell.

We retain the auxiliary pair and notation of \cref{sec:localsetting};
all depths and norms in the group constructions refer to its lattices.
Thus $b$ is the
depth of $J$ and $r$ is the precision of the cell. Put $n=b+r$;
character values are computed using the additive character
$x\mapsto\psi(\varpi^{-n}x)$. Throughout, $\lf\in\mathcal A_h$:
it is integral on $\Lat'$, vanishes on $\uu$, and satisfies
$\norm{Q_\lf^{-1}}\le q^h$. We put $a=r-h$, so that $S\subset U_a$,
and assume \eqref{eq:localmargins}.

We first justify the initial choice of integral auxiliary coordinates.
With these coordinates fixed, the orders of the remainders give the
required depth estimates.
Coordinate errors of depth at least $b$ vanish in additive quotients
such as $\sS/\pp^b\Lat_+$, and only the class modulo $\pp^n$ of the
\emph{phase} $\Phi$ of a character value $\psi(\varpi^{-n}\Phi)$ matters.

\subsection{Integral coordinates and depth estimates}\label{sec:estimates}
We justify the initial auxiliary choice in \cref{sec:localsetting}.
Keep the product parametrizations fixed. Starting with an
\etacompat{} lattice, shrink it until $(\kappa',\Lat')$ is standard
(\cref{sec:standardgroups}). This gives bracket closure and integral
coefficients for the group law, inverse, and the maps obtained by
multiplying, reordering, and conjugating elements of the three factors.
A further fixed shrinking makes every form in $\mathcal K$ integral
on $\Lat'$, while preserving standardness.

We also require integral coefficients for conjugation by
$t=\eta(\varpi)$ on $U$ and by $t^{-1}$ on $U^-$.
Their linear terms preserve the respective lattices.
Replacing $\Lat'$ by $\pp^k\Lat'$ multiplies each coefficient of
degree $d\ge2$ by $\varpi^{k(d-1)}$; since these coordinate maps
are polynomial, a fixed $k$ makes all their coefficients integral.

Let $X,Y$ have depths at least $u,v\ge1$, respectively.
A mixed monomial---one containing at least one coordinate from each
input---has depth at least $u+v$. If its total degree is at least
three, it contains at least one additional coordinate, contributing
depth at least $\min(u,v)$. Its depth is therefore at least
$u+v+\min(u,v)$. These bounds also hold for convergent sums of such
monomials with integral coefficients.

We also need bounds for differences. Suppose $X$ has depth at least
$u$, $Y,Y'$ have depth at least $v$, and $Y-Y'$ has depth at least
$k\ge v$. In the difference of a mixed monomial evaluated at
$(X,Y)$ and $(X,Y')$, each term contains a factor from $Y-Y'$.
Its depth is therefore at least $u+k$, or at least
$u+k+\min(u,v)$ if the monomial has total degree at least three.
We use these elementary observations throughout.

Write $\mu_\pm,\mu_M$ for the group laws in $\xi_\pm,\xi_M$, and
$\mu$ for the group law in $\kappa'$. The other maps needed are
\[
\begin{aligned}
\rho(X,Y)&=\kappa'\bigl(\xi_-(Y)\xi_+(X)\bigr),\\
\gamma_-(X,Y)&=\kappa'\bigl(\xi_-(Y)\xi_+(X)\xi_-(Y)^{-1}\bigr),\\
\gamma_+(X,Y)&=\kappa'\bigl(\xi_+(X)\xi_-(Y)\xi_+(X)^{-1}\bigr).
\end{aligned}
\]
In the product chart $\kappa'$, write the group law as
\[
\mu(X,Y)=X+Y+\beta(X,Y)+O(\deg\ge3),
\]
where $\beta$ is the quadratic term.
Then $\beta(X,Y)-\beta(Y,X)=[X,Y]$, and
\begin{equation}\label{eq:serre}
\kappa'\bigl(\kappa'^{-1}(X)\kappa'^{-1}(Y)
\kappa'^{-1}(X)^{-1}\bigr)
=Y+[X,Y]+O(\deg\ge3)
\end{equation}
(\cite[Part~II, Chapter~IV, \S7, and Chapter~V, \S1]{MR2179691}).
In the product chart $\xi_+(X)\xi_-(Y)$ has coordinate $X+Y$.
Consequently
\begin{equation}\label{eq:quadratic}
\begin{aligned}
\rho(X,Y)&=X+Y+[Y,X]+O(\deg\ge3),\\
\gamma_-(X,Y)&=X+[Y,X]+\mathscr R(X,Y),\\
\gamma_+(X,Y)&=Y+[X,Y]+\mathscr R'(X,Y).
\end{aligned}
\end{equation}
The remainders are mixed and of total degree at least three.
Set $\mathscr R_\lf=\lf\circ\mathscr R$. We retain the full pairing
$\Bf_\lf$: the form $\lf$ need not vanish on the
$\uu^-$-component of $[Y,X]$.

\begin{lemma}[Product and phase estimates]\label{lem:estimates}
Let $u,v\ge1$, and let $\lf\in\gd$ be integral on $\Lat'$ and
vanish on $\uu$. In
\cref{item:est-mixed,item:est-cubic,item:est-diff}, let
$X\in\uu$ and $Y\in\uu^-$ have depths at least $u,v$, respectively.
\begin{enumerate}[label=\textup{(\arabic*)},ref=\arabic*]
\item\label{item:est-group} \emph{Group laws.}
Let $\mu_\bullet$ be any of $\mu_\pm,\mu_M,\mu$.
For inputs $Z,W$ in its coordinate space of depths at least $u,v$,
respectively, $\mu_\bullet(Z,W)-Z-W$ has depth at least $u+v$.
If $Z$ has depth at least $u$, the coordinate of its group inverse
plus $Z$ has depth at least $2u$.
\item\label{item:est-mixed} \emph{Reordering and conjugation.}
The errors
$\rho(X,Y)-X-Y$, $\gamma_-(X,Y)-X$, and $\gamma_+(X,Y)-Y$
have depth at least $u+v$. More generally, conjugating an element
of depth at least $u$ in one factor by an element of depth at least
$v$ in another changes its coordinates by an element of depth at
least $u+v$.
\item\label{item:est-cubic} \emph{The commutator phase.}
\begin{equation}\label{eq:mixedphase}
\begin{aligned}
\lf(\gamma_-(X,Y))&=-\Bf_\lf(X,Y)+\mathscr R_\lf(X,Y),\\
\val(\mathscr R_\lf(X,Y))&\ge u+v+\min(u,v).
\end{aligned}
\end{equation}
The phase $\lf(\rho(X,Y))-\lf(Y)$ has the same leading term and
error bound; for $\lf(\gamma_+(X,Y))-\lf(Y)$ the leading term is
$\Bf_\lf(X,Y)$, with the same bound.
\item\label{item:est-diff} \emph{Differences of remainders.}
If $Y'\in\uu^-$ satisfies $\val(Y')\ge v$ and
$\val(Y-Y')\ge k\ge v$, then
\begin{equation}\label{eq:difference}
\val\bigl(\mathscr R_\lf(X,Y)-\mathscr R_\lf(X,Y')\bigr)
\ge u+k+\min(u,v).
\end{equation}
\end{enumerate}
\end{lemma}

\begin{proof}
The group-law and conjugation errors are mixed, and the inversion
errors have degree at least two. This gives
\cref{item:est-group,item:est-mixed}.
Applying $\lf$ to \eqref{eq:quadratic} gives the leading phases in
\cref{item:est-cubic}, since $\lf|_\uu=0$.
Every remaining monomial has multidegree at least $(2,1)$ or $(1,2)$,
giving $\min(2u+v,u+2v)=u+v+\min(u,v)$.
For differences, expand $Y^\beta-Y'^\beta$ with one factor from
$Y-Y'$; the same argument gives \cref{item:est-diff}.
\end{proof}

Let $w_+$ be the largest weight of $\eta$ on $\uu$.

\begin{lemma}[Contraction estimates]\label{lem:contraction}
For $d\ge1$,
\[
tU_dt^{-1}\subset U_{d+1},\qquad
t^{-1}U^-_dt\subset U^-_{d+1},\qquad
\bigcup_{j\ge0}t^{-j}U_dt^j=U.
\]
For an element of $U$ of depth $u\ge2w_++1$, conjugation by
$t^{-1}$ gives an element of depth at least $u-w_+$.
\end{lemma}

\begin{proof}
In $\xi_+$ coordinates, conjugation by $t$ is
$\phi(X)=A_+X+\varepsilon(X)$, where
$A_+=\Ad(t)|_\uu$ has weights $\varpi^s$, $1\le s\le w_+$, and
$\varepsilon$ is integral of degree at least two. Thus an input of
depth $d\ge1$ has image of depth at least $d+1$. The argument for
$t^{-1}$ on $U^-$ is identical. The exhaustion follows from the
contracting action of $\eta$ on $U$.

On $\pp^{w_++1}\Lat_+$ the map
$X\mapsto X+A_+^{-1}\varepsilon(X)$ is a bijection: its nonlinear
part has Lipschitz constant at most $q^{-1}$, so the equation for
each prescribed image is solved by contraction. Hence
\[
\phi(\pp^{w_++1}\Lat_+)
=A_+(\pp^{w_++1}\Lat_+)\supset\pp^{2w_++1}\Lat_+.
\]
For such $X$,
$\val(\varepsilon(X))\ge2\val(X)>\val(A_+X)$, so
$\val(\phi(X))=\val(A_+X)\le\val(X)+w_+$.
Injectivity of conjugation now gives the inverse-depth bound for
every element of depth at least $2w_++1$.
\end{proof}

For the remaining proofs we use only the following depth inequalities:
\begin{equation}\label{eq:used}
2a\ge b,\qquad b\ge r,\qquad r\ge2h+2w_++1.
\end{equation}
They follow from $b-r\ge C$, $2a-b\ge C$ with $C\ge w_++1$, since
\[
r-2h=(b-r)+(2a-b)\ge2C.
\]
They imply $b\ge r\ge h+w_+$ and $a=r-h\ge h+2w_++1\ge2w_++1$.
All inputs below therefore have positive depth. In particular,
quadratic errors of depth $2a$ vanish modulo depth $b$, and errors
of depth $2b\ge n$ do not change the character values.

\subsection{The groups and the extension}\label{sec:groupdetails}
Recall $\sS$ from \eqref{eq:Slattice}. We will also need its
opposite counterpart:
\begin{equation}\label{eq:ST}
\begin{gathered}
\sS=\pp^r(Q^*)^{-1}\Lat_-^\vee,\qquad
\sT=\{Y\in\uu^-:\Bf_\lf(\pp^b\Lat_+,Y)
\subset\pp^n\}=\pp^rQ^{-1}\Lat_+^\vee,\\
S=\xi_+(\sS),\qquad T=\xi_-(\sT),\qquad R=U^-_b.
\end{gathered}
\end{equation}
Since $Q$ is integral and $\norm{Q^{-1}}\le q^h$, both lattices
lie between the indicated depth lattices:
\[
\begin{aligned}
\pp^b\Lat_+&\subset\pp^r\Lat_+\subset\sS\subset\pp^a\Lat_+,\\
\pp^b\Lat_-&\subset\pp^r\Lat_-\subset\sT\subset\pp^a\Lat_-.
\end{aligned}
\]
Set $J=U_bM_bR$ and $H=SM_bR$. The product factorizations are unique
by the big cell. On $H$ define $\theta_\lf$ by \eqref{eq:thetaJ}.

\begin{lemma}[Groups and character extension]\label{lem:groups}
The sets $S,T,R,J,H$ are compact subgroups; $J,H$ are open, and
$J$ is normal in $H$. The coordinate maps give
\[
H/J\simeq\sS/\pp^b\Lat_+,\qquad
T/R\simeq\sT/\pp^b\Lat_-.
\]
Moreover, $T$ normalizes $H,J,R$, and $\theta_\lf$ is a character
of $H$, trivial on $S$.
\end{lemma}

\begin{proof}
\emph{Subgroups.}
In $S$ and $T$, multiplication and inversion agree with addition and
negation modulo depth $2a\ge b$. Since the defining lattices
contain $\pp^b\Lat_\pm$, they are closed under these operations.
The same argument applies to $U_b^\pm$ and $M_b$. In the product
chart, multiplication and inversion preserve $\pp^b\Lat'$,
so $J=\kappa'^{-1}(\pp^b\Lat')$ is a compact open subgroup.
The error in its group law has depth $2b\ge n$; therefore
\[
\theta_J(x)=\psi(\varpi^{-n}\lf(\kappa'(x)))
\]
is a character of $J$, as in \cref{sec:cells}.

\smallskip\noindent\emph{Character invariance.}
For $s=\xi_+(X)\in S$ and $j\in J$, the mixed estimate shows that
conjugation changes the coordinates of $j$ only at depth
$a+b\ge b$. Applying this also to $s^{-1}$ shows that $S$
normalizes $J$. To prove invariance of $\theta_J$, it suffices to
consider the three factors of $J$.
On $U_b$ both character values are $1$. For $j\in M_b$,
$sjs^{-1}=s'j$ with $s'\in U_b$: indeed $s'\in U$ by the Levi
action, and $s'\in J$ by normalization. Thus the values again agree.
For $j=\xi_-(Y)\in R$, the change of phase is
\[
\Bf_\lf(X,Y)+\mathscr R'_\lf(X,Y).
\]
The first term has depth at least $n$ by the definition of $\sS$.
Since $a\le b$, the remainder has depth at least
$2a+b\ge n$; here $2a\ge b\ge r$. Thus the change of phase lies in $\pp^n$
and does not change the character value.

\smallskip\noindent\emph{Extension.}
It follows that $H=SJ$ is a compact open subgroup and $J$ is normal
in $H$. Since $S\cap J=U_b$, the coordinate map identifies
$H/J$ with $\sS/\pp^b\Lat_+$. The formula
$\theta_\lf(sj)=\theta_J(j)$ is well defined because $\theta_J$
is trivial on $U_b$. It is a character because $S$ fixes $\theta_J$:
\[
\begin{aligned}
\theta_\lf(s_1j_1s_2j_2)
&=\theta_J(s_2^{-1}j_1s_2)\theta_J(j_2)\\
&=\theta_J(j_1)\theta_J(j_2).
\end{aligned}
\]
This agrees with \eqref{eq:thetaJ} and is trivial on $S$.

\smallskip\noindent\emph{Normalization by $T$.}
For $y\in T$ and $s=\xi_+(X)\in S$, conjugation changes $X$ only
at depth $2a\ge b$ and introduces $\m$- and $\uu^-$-coordinates
of depth at least $b$. Hence $ysy^{-1}\in H$. Similarly,
$yJy^{-1}\subset J$, since the coordinate change has depth at least
$a+b\ge b$. Applying the same argument to $y^{-1}$ proves that
$T$ normalizes $H$ and $J$. Also $R=J\cap U^-$, so $T\subset U^-$
normalizes $R$. Finally, the group law of $T$ is addition modulo
$\pp^b\Lat_-$, giving the asserted description of $T/R$.
\end{proof}

\subsection{Normalizing an extension}\label{sec:normalizationdetails}
An extension of $\theta_J$ differs from $\theta_\lf$ by a character
of $S/U_b$. Write its parameter as $\delta$. Conjugation by
$y=\xi_-(Y)\in T$ changes the leading parameter from $\delta$ to
$\delta-QY$. Thus non-degeneracy of $Q$ suggests a unique class of
$Y$ for which the character becomes trivial on $S$. The proof below
shows that the nonlinear terms preserve this conclusion, then checks
that the new linear form satisfies the same bounds.

\begin{lemma}[Quantitative normalization]\label{lem:normalizationbounds}
Let $\theta$ extend $\theta_J$ to $H$. There is a unique class
$yR\in T/R$ such that $\theta^y|_S=1$, where
$\theta^y(z)=\theta(yzy^{-1})$. For this $y$ there is (a not necessarily unique)
$\lf'\in\mathcal A_h$ such that $\theta^y=\theta_{\lf'}$ and
\begin{equation}\label{eq:Aprime}
\begin{gathered}
\val(\lf'-\lf)\ge a,\qquad
\norm{Q_{\lf'}^{-1}}\le q^h,\\
(Q_{\lf'}^*)^{-1}\Lat_-^\vee=(Q^*)^{-1}\Lat_-^\vee.
\end{gathered}
\end{equation}
In particular $\lf'$ defines the same groups $S,H$.
\end{lemma}


\begin{proof}
\emph{Parameterizing the extensions.}
With respect to the pairing
$(\delta,X)\mapsto\psi(\varpi^{-n}\delta(X))$, the annihilators of
$\pp^b\Lat_+$ and $\sS$ in $\uu^*$ are, respectively,
$\pp^r\Lat_+^\vee$ and $\pp^bQ\Lat_-$. Consequently
\[
\widehat{S/U_b}\simeq
\mathcal E:=\pp^r\Lat_+^\vee/\pp^bQ\Lat_-.
\]
Every extension is therefore of the form
\[
\theta_\delta\bigl(\xi_+(X)\xi_M(Z)\xi_-(Y)\bigr)
=\psi\bigl(\varpi^{-n}(\lf(Z+Y)+\delta(X))\bigr),
\]
where we choose a representative
$\delta\in\pp^r\Lat_+^\vee$ of its class in $\mathcal E$.

For $y=\xi_-(Y)\in T$ and $X\in\sS$, the phase of
$\theta_\delta^y(\xi_+(X))$ is, modulo $\pp^n$,
\begin{equation}\label{eq:normalphase}
\delta(X)-\Bf_\lf(X,Y)+\mathscr R_\lf(X,Y).
\end{equation}
Indeed, let $\pi_+:\g\to\uu$ be the projection along $\m\oplus\uu^-$.
The phase differs from \eqref{eq:normalphase} by
$\delta(\pi_+\gamma_-(X,Y)-X)$. Since $X$ and $Y$ have depth at least $a$,
$\pi_+\gamma_-(X,Y)-X$ has depth at least $2a$ by
\cref{item:est-mixed} of \cref{lem:estimates}, and since $\delta$ has depth at
least $r$, this term has depth at least $r+2a\ge r+b=n$ by \eqref{eq:used}.

\smallskip\noindent\emph{Finding the normalizing element.}
First take $X\in\pp^b\Lat_+$, $Y\in\sT$. The two leading
terms in \eqref{eq:normalphase} have depth at least $n$: this follows
from $r+b=n$ for $\delta(X)$ and from the definition of $\sT$ for
$\Bf_\lf(X,Y)$. The remainder has depth at least
$b+2a\ge n$ by \eqref{eq:mixedphase} and \eqref{eq:used}.
Thus $\theta_\delta^y|_S$ is trivial on $U_b$.
It depends only on $yR$, because conjugation by $R\subset H$ fixes
every character of $H$. We obtain a map of finite sets
\[
\Upsilon:\sT/\pp^b\Lat_-\longrightarrow\widehat{S/U_b},
\qquad [Y]\longmapsto\theta_\delta^{\xi_-(Y)}|_S.
\]
Conceptually, $T$ acts trivially on $H/J$, because conjugation by an
element of $T$ changes an element of $S$ only by an element of $J$ (see the
proof of \cref{lem:groups}); hence it fixes each character of $H/J$, such as
$\theta_\delta/\theta_\lf$. Thus $\delta$ merely translates the target of $\Upsilon$.
The linear map $Q$ identifies its source with $\mathcal E$, so source
and target have the same cardinality. We prove injectivity by finding
an element of $S$ on which two distinct images differ.

Suppose $Y\not\equiv Y'\pmod{\pp^b\Lat_-}$, and put
$k=\val(Y-Y')<b$. Choose $\mu\in\Lat_-^\vee$ with
$\val(\mu(Y-Y'))=k$. For $z\in\OO$ set
\[
X_z=z\varpi^{n-k-1}(Q^*)^{-1}\mu.
\]
For this choice, the leading difference in \eqref{eq:normalphase} is
\[
-\Bf_\lf(X_z,Y-Y')=-z\varpi^{n-k-1}\mu(Y-Y').
\]
As $z$ varies, these values run through $\pp^{n-1}$.
Also $X_z\in\sS$, since $n-k-1\ge r$, and
\[
\val(X_z)\ge u:=n-k-h-1\ge a,\qquad k\ge a.
\]
Apply \eqref{eq:difference} with $v=a$. Its lower bound is
\[
u+k+a=n+r-2h-1,
\]
which is at least $n$ by \eqref{eq:used}. Hence the nonlinear
difference does not affect the character. Since $\psi$ is nontrivial
on $\pp^{-1}$, some $z$ distinguishes the two images of
$\Upsilon$. The map is therefore injective, hence bijective. The
trivial character has a unique preimage, proving existence and
uniqueness of $yR$.

\smallskip\noindent\emph{Controlling the new parameter.}
Choose such a $y=\xi_-(Y)$ and put $\theta'=\theta_\delta^y$.
Regard $\delta$ as a form on $\g$ vanishing on $\m\oplus\uu^-$,
and set $\tilde\lf=\Ad^*(y^{-1})(\lf+\delta)$.
On $M_bR$, conjugation by $y$ has linear term $\Ad(y)$ and error
of depth at least $2b\ge n$. Therefore
\[
\theta'(z)=\psi(\varpi^{-n}\tilde\lf(\kappa'(z)))
\qquad(z\in M_bR).
\]
Define $\lf'$ to agree with $\tilde\lf$ on $\m\oplus\uu^-$ and
to vanish on $\uu$. The integral conjugation series gives
$\val(\Ad(y^{\pm1})-1)\ge\val(Y)\ge a$. Since $\val(\delta)\ge r$,
we have $\val(\lf'-\lf)\ge a$. In particular, $\lf'$ is integral
on $\Lat'$ and vanishes on $\uu$.
Write $Q_{\lf'}=Q+\Delta$. Then
\[
\val(\Delta)\ge a,\qquad
\val(Q^{-1}\Delta)\ge r-2h\ge1.
\]
Thus $A=1+Q^{-1}\Delta$ is a lattice automorphism of $\Lat_-$
and $Q_{\lf'}=QA$. It follows that
\[
\begin{aligned}
(Q_{\lf'}^*)^{-1}\Lat_-^\vee
&=(Q^*)^{-1}(A^*)^{-1}\Lat_-^\vee
=(Q^*)^{-1}\Lat_-^\vee,\\
\norm{Q_{\lf'}^{-1}}&=\norm{A^{-1}Q^{-1}}\le q^h.
\end{aligned}
\]
This proves \eqref{eq:Aprime}. The first equality preserves $S$,
and hence $H$. By \cref{lem:groups}, $\theta_{\lf'}$ is a character
of this same $H$. It agrees with $\theta'$ on $M_bR$, and both are
trivial on $S$, so $\theta'=\theta_{\lf'}$.
\end{proof}

\subsection{Verification of the moving chain}\label{sec:chaindetails}
We verify that the groups and characters constructed at each stage
satisfy the hypotheses of \cref{lem:criterion}, so that successive
averaging remains injective.
Set $t=\eta(\varpi)$, $\lf_j=\Ad^*(t^j)\lf$ and
$Q_j=Q_{\lf_j}$. Define $\sS_j,S_j,H_j=S_jM_bR$ and
$\theta_j=\theta_{\lf_j}$ as before, with $\lf_j$ in place of
$\lf$. Thus $S_0=S$, $H_0=H$ and $\theta_0=\theta_\lf$.
We first show that the inverse bound for $Q_j$ is uniform in $j$.
The rest of the proof verifies the four properties used in
\cref{sec:jacquet}.

\begin{lemma}[The moving chain]\label{lem:chain}
For every $j\ge0$, the form $\lf_j$ lies in $\mathcal A_h$ and agrees
with $\lf$ on $\m$.
Thus \cref{lem:groups} applies with the same constants and
$\pp^r\Lat_+\subset\sS_j\subset\pp^a\Lat_+$.
Moreover:
\begin{enumerate}[label=\textup{(\alph*)},ref=\alph*]
\item\label{item:crit-incl} \emph{Nesting:}
$tS_jt^{-1}\subset S_{j+1}$.
\item\label{item:crit-cov} \emph{Character compatibility:}
$\theta_{j+1}(u)=\theta_j(t^{-1}ut)$ for $u\in M_bR$.
\item\label{item:crit-nondeg} \emph{Separation of new cosets:}
for $x\in t^{-1}S_{j+1}t\setminus S_j$ there is $u\in R$ with
$xux^{-1}\in H_j$ and $\theta_j(xux^{-1})\ne\theta_j(u)$.
Equivalently,
\[
S_j=\{x\in t^{-1}S_{j+1}t:\theta_j([x,u])=1
\ \ \forall u\in R\cap x^{-1}H_jx\}.
\]
\item\label{item:crit-exhaust} \emph{Exhaustion:}
$\bigcup_{j\ge0}t^{-j}S_jt^j=U$.
\end{enumerate}
Also $M_b$ normalizes $R$ and $t^{-1}Rt\subset R$.
\end{lemma}

\begin{proof}
\emph{Uniform control of the pairing.}
Since $t$ centralizes $M$ and $\Ad(t^{-j})$ contracts $\uu^-$,
the form $\lf_j$ agrees with $\lf$ on $\m$, still vanishes on
$\uu$, and remains integral on $\Lat'$. Put
$A_+=\Ad(t)|_{\uu}$ and $A_-=\Ad(t^{-1})|_{\uu^-}$; on the
weight spaces $\g_s$ and $\g_{-s}$ these are multiplication by
$\varpi^s$, $s>0$. Equivariance of the pairing gives
\begin{equation}\label{eq:Qj}
Q_j=(A_+^*)^{-j}Q_0A_-^j,\qquad
Q_j^{-1}=A_-^{-j}Q_0^{-1}(A_+^*)^j.
\end{equation}
The block of $Q_0$ from $\g_{-s'}$ to $\g_s^*$ vanishes when
$s>s'$, because $[\g_s,\g_{-s'}]\subset\uu$ is killed by $\lf$.
Thus $Q_0$ is block triangular (as in the proof of \cref{lem:levi}), and
so is its inverse. In \eqref{eq:Qj}, every possibly nonzero block is
multiplied by a nonnegative power of $\varpi$: by $\varpi^{j(s'-s)}$ for
$Q_0$, and by $\varpi^{j(s-s')}$ for the inverse block from
$\g_s^*$ to $\g_{-s'}$. Since the lattices split into weight spaces,
\[
Q_j\Lat_-\subset\Lat_+^\vee,\qquad
\norm{Q_j^{-1}}\le q^h.
\]
This proves the asserted uniformity and the bounds on $\sS_j$.
Also $M_b$ normalizes $R=J\cap U^-$, and
$t^{-1}Rt\subset R$ by \cref{lem:contraction}.

\smallskip\noindent\emph{Nesting.}
From \eqref{eq:Qj},
\[
(Q_{j+1}^*)^{-1}
=A_+(Q_j^*)^{-1}(A_-^*)^{-1}.
\]
Since $(A_-^*)^{-1}\Lat_-^\vee\supset\pp^{-1}\Lat_-^\vee$,
we get $\Ad(t)\sS_j\subset\pp\sS_{j+1}$.
The error of conjugation has depth $2a\ge r$, so it too lies in
$\sS_{j+1}$. Hence $tS_jt^{-1}\subset S_{j+1}$.

\smallskip\noindent\emph{Character compatibility.}
For $u=\xi_-(Y)\in R$, the error of conjugation by $t^{-1}$ has
depth $2b\ge n$. It does not affect the character, and therefore
\[
\theta_j(t^{-1}ut)
=\psi(\varpi^{-n}\lf_j(\Ad(t^{-1})Y))
=\theta_{j+1}(u).
\]
The equality on $M_b$ follows because $t$ centralizes $M$ and the
forms agree there. This proves compatibility on $M_bR$.

\smallskip\noindent\emph{Separation of new cosets.}
Take $x=\xi_+(X)\in t^{-1}S_{j+1}t\setminus S_j$.
By \cref{lem:contraction},
\[
\val(X)\ge a':=a-w_+.
\]
Since $X\notin\sS_j$, there is $Y_0\in\pp^b\Lat_-$ such that
$\Bf_{\lf_j}(X,Y_0)\notin\pp^n$. The ideal
$\varpi^{-n}\Bf_{\lf_j}(X,Y_0)\OO$ therefore contains $\pp^{-1}$.
Since $\psi$ is nontrivial on $\pp^{-1}$, some $z\in\OO$ gives
$\psi(\varpi^{-n}\Bf_{\lf_j}(X,zY_0))\ne1$.
Put $Y=zY_0\in\pp^b\Lat_-$ and $u=\xi_-(Y)$.
The mixed coordinate error in $xux^{-1}$ has depth at least
$a'+b\ge b$. Thus its $\m$- and $\uu^-$-coordinates have depth
at least $b$, and its $\uu$-coordinate lies in
$\pp^r\Lat_+\subset\sS_j$. Hence $xux^{-1}\in H_j$.

It remains to check that the nonlinear remainder lies in $\pp^n$
and hence does not change the character value. Since
$a'\le b$, its depth is at least $2a'+b$. The excess over the
required depth $n$ is
\[
2a'+b-n=r-2h-2w_+\ge1.
\]
This follows from \eqref{eq:used}.
Consequently
\[
\theta_j(xux^{-1})\theta_j(u)^{-1}
=\psi(\varpi^{-n}\Bf_{\lf_j}(X,Y))\ne1,
\]
which proves separation.

\smallskip\noindent\emph{Exhaustion.}
The inclusion $S_j\supset U_r$ holds for every $j$, while
$\bigcup_jt^{-j}U_rt^j=U$ by \cref{lem:contraction}.
\end{proof}


\subsection{Changing the chart}\label{sec:chartchange}
The following lemma compares cell characters in arbitrary normalized
charts; here $(\kappa',\Lat')$ need not be the auxiliary product pair.
It is used in \cref{sec:returncell} to pass from the statement in the
product chart to the original cell. Its constants are made explicit
in the unramified case in the proof of \cref{thm:regularunram}.
\begin{lemma}[Change of \nchart{}]\label{lem:compare}
Let $\kappa'$ be another \nchart{}, $\Lat'$ another
lattice in $\g$, and $\mathcal B\subset\gd$ a bounded set. Fix an integer $s\ge0$
with $\pp^s\Lat'\subset\Lat$. There are constants $b_0,C$ such
that for $\lf\in\mathcal B$ and integers $b,n$ with $b\ge b_0$, $2b-n\ge C$
(in particular $b_0$ is chosen so that $\pp^{b_0}\Lat$ and
$\pp^{b_0+s}\Lat'$ lie in the images of $\kappa$ and $\kappa'$), the formulas
\[
\chi_{n,\lf}(g)=\psi(\varpi^{-n}\lf(\kappa(g))),\qquad
\chi'(g)=\psi(\varpi^{-n}\lf(\kappa'(g)))
\]
define characters of $K_b=\kappa^{-1}(\pp^b\Lat)$ and
$J_b=\kappa'^{-1}(\pp^{b+s}\Lat')$, respectively.
Moreover, $J_b\subset K_b$ and $\chi'=\chi_{n,\lf}|_{J_b}$.
Consequently
\begin{equation}\label{eq:compare}
e_{J_b}^{\chi'}*e_{K_b}^{\chi_{n,\lf}}=e_{K_b}^{\chi_{n,\lf}}=e_{K_b}^{\chi_{n,\lf}}*e_{J_b}^{\chi'}.
\end{equation}
The constants can be taken uniform when $\kappa',\Lat'$ are
transported by conjugation by elements of a fixed compact subset of $G$,
with $s$ chosen so that the lattice inclusion holds throughout the family.
\end{lemma}

\begin{proof}
For $b$ sufficiently large, $J_b$ is a subgroup by
\cref{sec:standardgroups}. Since both charts are normalized, their
transition map has the form
\[
\kappa\circ\kappa'^{-1}(X)=X+R(X),\qquad R(X)=O(X^2).
\]
Thus, for $X\in\pp^{b+s}\Lat'\subset\pp^b\Lat$,
\[
R(X)\in\pp^{2b-C_0}\Lat
\]
with a fixed constant $C_0$. For large $b$ this lies in $\pp^b\Lat$,
so $J_b\subset K_b$.

Boundedness of $\mathcal B$ gives an integer $c\ge0$ such that
$\lf(\Lat)\subset\pp^{-c}$ for every $\lf\in\mathcal B$. Hence
\[
\val(\lf(R(X)))\ge2b-C_0-c.
\]
If $2b-n\ge C_0+c$, the transition error does not change the
character value, and $\chi'=\chi_{n,\lf}|_{J_b}$.
The standard group-law estimate \eqref{eq:law} also shows that
$\chi_{n,\lf}$ is a character when $2b-n\ge c$.
Thus $\chi'$ is a character by restriction, and the convolution
identities \eqref{eq:compare} follow from the subgroup inclusion
and agreement of the characters.

For uniformity, replace $\kappa'$ and $\Lat'$ by
$x\mapsto\Ad(g)\kappa'(g^{-1}xg)$ and $\Ad(g)\Lat'$, with $g$
in a fixed compact set. Compactness permits a common choice of $s$.
The coordinate maps depend analytically on $g$; a finite cover of
the compact set gives a common neighborhood and uniform bounds
for the group laws and transition remainders. The same proof
therefore gives uniform $b_0,C$.
\end{proof}


\section{Examples}\label{sec:examples}

We explicate the elliptic set in a few cases and compute the elliptic
densities for three lattices in $\Mat_2(F)$.
We also calculate the Plancherel masses for $\GL_2$ and for $\PGL_2$
in characteristic $2$, obtaining exact formulas that also measure the
sharpness in the general error estimate.

\subsection{The elliptic set}

\begin{example}[$\GL_d$]\label{ex:gld}
Let $\bG=\GL_d$ and identify $\gd$ with $\Mat_d(F)$ by the trace
pairing. For the maximal parabolic subgroup stabilizing a nonzero
proper subspace $W$, the Lie algebra $\uu_P=\Hom(F^d/W,W)$ has
annihilator equal to the set of endomorphisms preserving $W$.
Hence $\Ell{\gd}$ consists of the matrices with no nonzero proper
invariant subspace, or equivalently with characteristic polynomial
irreducible over $F$ (in any characteristic, including inseparable
polynomials). For $\Lat=\Mat_d(\OO)$ we have $\dLat=\Mat_d(\OO)$, and
reduction modulo $\pp$ gives the lower bound
\[
\elld[\GL_d]{\Mat_d(\OO)}\ge\frac{I_d(q)\abs{\GL_d(k)}}{q^{d^2}(q^d-1)},
\]
where $I_d(q)=\frac1d\sum_{e\mid d}\mu(e)q^{d/e}$ is the number of
monic irreducible polynomials of degree $d$ over $k$: each of them is
the characteristic polynomial of $\abs{\GL_d(k)}/(q^d-1)$ matrices over
$k$, and every lift of such a matrix has irreducible characteristic
polynomial over $F$.
\end{example}

The following elementary lemma computes the elliptic densities of
$\GL_2$ for three lattices, in every characteristic.
\begin{lemma}[Random monic quadratics]\label{lem:quadratics}
Let $\rho_{\OO}$, resp.\ $\rho_{\pp}$, be the probability that
$T^2+bT+c$ is irreducible over $F$ when $(b,c)$ is uniformly distributed
in $\OO^2$, resp.\ in $\pp\times\pp$. Then
\[
\rho_{\OO}=\frac{q+2}{2(q+1)},\qquad\rho_{\pp}=\frac{2q+1}{2(q+1)}.
\]
Moreover, let $x\in\OO$, $y\in\OO^\times$, and
$X=\begin{psmallmatrix}x&y\\z&w\end{psmallmatrix}$. If $(z,w)$ is uniformly
distributed in $\OO^2$, then the characteristic polynomial of $X$ is
irreducible with probability $\rho_{\OO}$; if $x\in\pp$ and $(z,w)$ is
uniformly distributed in $\pp\times\pp$, it is irreducible with probability
$\rho_{\pp}$. The same holds for the transpose of $X$.
\end{lemma}

\begin{proof}
A monic quadratic polynomial over $F$ is reducible if and only if it has
a root in $F$, and a root of a monic polynomial with coefficients in $\OO$
lies in $\OO$.

Let $(b,c)$ be uniform in $\pp\times\pp$. With probability $(q-1)/q$ we
have $\val(c)=1$, and $T^2+bT+c$ is an Eisenstein polynomial, hence
irreducible. With probability $1/q$ we have $c\in\pp^2$; then
$T^2+bT+c=\varpi^2\bigl(S^2+b'S+c'\bigr)$ with $S=T/\varpi$, $b'=b/\varpi$,
$c'=c/\varpi^2$, and conditionally $(b',c')$ is uniform in $\OO^2$. Hence
$\rho_{\pp}=\frac{q-1}q+\frac1q\rho_{\OO}$.

Let $(b,c)$ be uniform in $\OO^2$, and reduce $f=T^2+bT+c$ modulo $\pp$:
$\bar f$ is uniformly distributed among the $q^2$ monic quadratic polynomials
over $k$. If $\bar f$ is irreducible, which happens for $(q^2-q)/2$ of them,
then $f$ is irreducible. If $\bar f$ has two distinct roots in $k$, which
happens for $(q^2-q)/2$ of them, then $f$ has a root in $\OO$ by Hensel's
lemma. In the remaining $q$ cases $\bar f=(T-\bar a)^2$ for some
$\bar a\in k$. Fix a lift $a\in\OO$ of $\bar a$. Then
$f(T+a)=T^2+(b+2a)T+(c+ab+a^2)$, and $(b,c)\mapsto(b+2a,c+ab+a^2)$ is a
measure preserving bijection from the residue class of $(b,c)$ onto
$\pp\times\pp$: it is affine-triangular with unit diagonal coefficients
in $b$ and $c$, and its image consists of the coefficients of a lift of
$\bar f(T+\bar a)=T^2$. Hence
$\rho_{\OO}=\frac{q-1}{2q}+\frac1q\rho_{\pp}$. The two linear relations
give the stated values.

For the last assertion, the characteristic polynomial of $X$ (and of its
transpose) is $T^2-(x+w)T+(xw-yz)$, and
$(z,w)\mapsto(-(x+w),xw-yz)$ is a measure preserving bijection of $\OO^2$
onto itself, since $w\mapsto-(x+w)$ is one of $\OO$ and, for fixed $w$,
$z\mapsto xw-yz$ is one of $\OO$, as $y$ is a unit. If $x\in\pp$ it
restricts to a measure preserving bijection of $\pp\times\pp$ onto itself.
\end{proof}

\begin{proposition}[The hyperspecial $\GL_2$ lattice]\label{prop:gl2}
For every non-archimedean local field $F$, $\elld[\GL_2]{\Mat_2(\OO)}=\frac12$.
\end{proposition}

\begin{proof}
By \cref{ex:gld}, $\elld[\GL_2]{\Mat_2(\OO)}$ is the probability
$P$ that the characteristic polynomial of a uniformly distributed
$X\in\Mat_2(\OO)$ is irreducible over $F$. Conjugation by $\GL_2(\OO)$
preserves Haar measure on $\Mat_2(\OO)$ and this property, and
$\GL_2(\OO)\to\GL_2(k)$ is surjective. We stratify according to the
conjugacy class of $\bar X\in\Mat_2(k)$ \cite[5.3~Proposition,
6.1]{MR2234120}; conditionally on $\bar X$, the matrix $X$ is uniformly
distributed in a coset of $\pp\Mat_2(\OO)$.
\begin{enumerate}[label=\textup{(\alph*)},ref=\alph*]
\item $\bar X$ has irreducible characteristic polynomial: there are
$(q^2-q)/2$ such polynomials, each the characteristic polynomial of
$q(q-1)$ matrices, so this has probability $\frac{(q-1)^2}{2q^2}$, and
then the characteristic polynomial of $X$ is irreducible.
\item $\bar X$ has two distinct eigenvalues in $k$: probability
$\frac{q(q-1)}2\cdot\frac{q(q+1)}{q^4}=\frac{q^2-1}{2q^2}$, and then the
characteristic polynomial of $X$ has a root in $F$ by Hensel's lemma.
\item $\bar X=\bar a+\bar N$ with $\bar a\in k$ and $\bar N\ne0$
nilpotent: probability $\frac{q(q^2-1)}{q^4}=\frac{q^2-1}{q^3}$. After
conjugation by $\GL_2(\OO)$ we may assume
$\bar X=\begin{psmallmatrix}\bar a&1\\0&\bar a\end{psmallmatrix}$. Fix a
lift $a\in\OO$ of $\bar a$; then
$X-a=\begin{psmallmatrix}\varpi s&1+\varpi t\\\varpi u&\varpi v\end{psmallmatrix}$
with $s,t,u,v$ uniform in $\OO$, and the characteristic polynomial of $X$
is irreducible if and only if that of $X-a$ is. By the last assertion
of \cref{lem:quadratics} (with $x=\varpi s$, $y=1+\varpi t$), this has
conditional probability $\rho_{\pp}$.
\item $\bar X=\bar a$ is scalar: probability $q^{-3}$. Then
$X=a+\varpi Y$ with $Y$ uniform in $\Mat_2(\OO)$, and the characteristic
polynomial of $X$ is irreducible if and only if that of $Y$ is; the
conditional probability is $P$.
\end{enumerate}
Since $\frac{q^2-1}{q^3}\rho_{\pp}=\frac{(q-1)(2q+1)}{2q^3}$, we obtain
\[
P=\frac{(q-1)^2}{2q^2}+\frac{(q-1)(2q+1)}{2q^3}+q^{-3}P
=\frac{1-q^{-3}}2+q^{-3}P,
\]
whence $P=\frac12$.
\end{proof}


For $F=\Q_p$, including $p=2$, the same value follows from the joint
eigenvalue distribution of Haar random matrices in $\Mat_2(\Z_p)$
computed in \cite{SVP26}: by \cite[Theorem~1.7]{SVP26} with $n=m=2$
and $K_1=K_2=\Q_p$, the density of ordered pairs of distinct eigenvalues
in $\Z_p^2$ is $(1+p^{-1})\abs{x-y}$, so the probability that both eigenvalues lie in
$\Z_p$ is $\frac{1+p^{-1}}2\int_{\Z_p^2}\abs{x-y}\,dx\,dy=\frac12$. The
exact computation of $\scd{K_n}$ in \cref{prop:gl2exact} below gives
$\lim_n\scd{K_n}=\frac12$ for every $F$, in accordance with
\cref{thm:main} and \cref{prop:gl2}.

To see the dependence on the lattice, consider
\begin{equation}\label{eq:iwahorilattices}
\Lat_1=\begin{pmatrix}\OO&\OO\\\pp&\OO\end{pmatrix},\qquad
\Lat_2=\begin{pmatrix}\OO&\pp^{-1}\\\OO&\OO\end{pmatrix}.
\end{equation}
$\Lat_1$ is the Lie algebra of the Iwahori group scheme of $\GL_2$,
whose $\OO$-points form the Iwahori subgroup $\Iw$ of \cref{sec:gl2exact},
and the tower of $\Lat_1$ consists of its congruence kernels
$\{g\in\Iw:g\equiv1\ (\pp^n),\ g_{21}\equiv0\ (\pp^{n+1})\}$
(\cref{lem:model}).
By contrast, $\Lat_2$ is not closed under the bracket; its tower is
defined for $n$ large, as the tower of a rescaled standard pair
$(\kappa,\pp^{n_0}\Lat_2)$ (\cref{sec:standardgroups}).

\begin{proposition}[Iwahori-type lattices]\label{prop:iwahori}
For every non-archimedean local field $F$,
\[
\elld{\Lat_1^\vee}=\frac q{q+1},\qquad \elld{\Lat_2^\vee}=\frac1{q+1}.
\]
\end{proposition}

\begin{proof}
For $i=1,2$ let $\phi_i$ be the probability that the characteristic
polynomial of a uniformly distributed element of $\Lat_i$ is irreducible
over $F$. Under the trace pairing $\Lat_1^\vee=\Lat_2$ and
$\Lat_2^\vee=\Lat_1$, so $\elld{\Lat_1^\vee}=\phi_2$ and
$\elld{\Lat_2^\vee}=\phi_1$ by \cref{ex:gld}.

We first compute $\phi_1$. Let
$Y=\begin{psmallmatrix}a&b\\c&d\end{psmallmatrix}$ be uniformly distributed in
$\Lat_1$, so that $a,b,d$ are uniform in $\OO$ and $c$ is uniform in $\pp$. The
reduction of $Y$ is upper triangular with diagonal entries $\bar a,\bar d$.
If $\bar a\ne\bar d$, which has probability $(q-1)/q$, the characteristic
polynomial of $Y$ has a root in $F$ by Hensel's lemma. If $\bar a=\bar d$, fix
a lift $a_0\in\OO$; then
$Y-a_0=\begin{psmallmatrix}\varpi s&b\\\varpi c'&\varpi v\end{psmallmatrix}$
with $s,b,c',v$ uniform in $\OO$. If $b\in\OO^\times$, which has
probability $(q-1)/q$, the conditional probability of irreducibility is
$\rho_{\pp}$ by \cref{lem:quadratics} with $x=\varpi s$ and $y=b$. If
$b\in\pp$, then $Y-a_0=\varpi Z$ with $Z$ uniform in $\Mat_2(\OO)$, and the
conditional probability is $\frac12$ by \cref{prop:gl2}. Hence
\[
\phi_1=\frac1q\Bigl(\frac{q-1}q\,\rho_{\pp}+\frac1{2q}\Bigr)
=\frac1q\cdot\frac{(q-1)(2q+1)+(q+1)}{2q(q+1)}=\frac1{q+1}.
\]

For $\phi_2$ let $Y=\begin{psmallmatrix}a&b\\c&d\end{psmallmatrix}$ be
uniformly distributed in $\Lat_2$, and write $b=\varpi^{-1}\beta$, so that
$a,\beta,c,d$ are uniform in $\OO$. If
$\beta,c\in\OO^\times$, which has probability $(1-q^{-1})^2$, then
$\det Y=ad-\varpi^{-1}\beta c$ has valuation $-1$ while $\tr Y\in\OO$. A
root $t\in F$ of $T^2-(\tr Y)T+\det Y$ would satisfy
$\val(t^2-(\tr Y)t)=-1$, which is impossible: this valuation is $\ge0$ if
$\val(t)\ge0$, and equals $2\val(t)$ if $\val(t)<0$. So $Y$ is elliptic.
If $c\in\OO^\times$ and $\beta\in\pp$, then $(\beta/\varpi,d)$ is uniform in
$\OO^2$, and the conditional probability is $\rho_{\OO}$ by
\cref{lem:quadratics} applied to the transpose, with $x=a$ and $y=c$. If
$\beta\in\OO^\times$ and $c\in\pp$, then $Y$ has the same characteristic
polynomial as $\begin{psmallmatrix}a&\beta\\c/\varpi&d\end{psmallmatrix}$,
with $(c/\varpi,d)$ uniform in $\OO^2$, and the conditional probability is
again $\rho_{\OO}$. Each of these two cases has probability
$q^{-1}(1-q^{-1})$. Finally, if $\beta,c\in\pp$, then $Y$ is uniformly
distributed in $\Lat_1$, and the conditional probability is $\phi_1$. Hence
\[
\begin{aligned}
\phi_2&=\Bigl(1-\frac1q\Bigr)^2+\frac{2(q-1)}{q^2}\cdot\frac{q+2}{2(q+1)}+\frac1{q^2(q+1)}\\
&=\frac{(q-1)^2(q+1)+(q-1)(q+2)+1}{q^2(q+1)}=\frac q{q+1}.
\end{aligned}\qedhere
\]
\end{proof}


The three lattice calculations are
summarized below. By \cref{thm:main}, the same values are the limiting
supercuspidal densities of the corresponding congruence towers.
\begin{center}
\begin{tabular}{lcc}
\toprule
Lattice $\Lat$ & $\elld{\Lat^\vee}$ & Calculation \\
\midrule
$\Mat_2(\OO)$ & $1/2$ & \Cref{prop:gl2} \\
$\Lat_1$ (Iwahori) & $q/(q+1)$ & \Cref{prop:iwahori} \\
$\Lat_2=\Lat_1^\vee$ & $1/(q+1)$ & \Cref{prop:iwahori} \\
\bottomrule
\end{tabular}
\end{center}

\subsection{The Plancherel mass of \texorpdfstring{$\GL_2$}{GL2}}\label{sec:gl2exact}

We compute the nonsupercuspidal mass directly from the Plancherel
formula and subtract it from the total mass. (For $\GL_n$, the number of
discrete series representations with central character trivial on $\varpi$
and with nonzero $K_{j+1}$-fixed vectors is computed in
\cite[Theorem~1]{MR2303528}; that count by itself does not give their
contribution to the Plancherel integral of $\dim\pi^{K_{j+1}}$.) The calculation applies
to every non-archimedean local field, including those of residual
characteristic $2$.

Let $G=\GL_2(F)$, $K=\GL_2(\OO)$, let $\Iw\subset K$ be the Iwahori
subgroup of matrices which are upper triangular modulo $\pp$, and
let $K_n=1+\pp^n\Mat_2(\OO)$, $n\ge1$; this is the tower
\eqref{eq:tower} for $\kappa(g)=g-1$ and $\Lat=\Mat_2(\OO)$. We
normalize $dg$ by $\vol(K)=1$, so that
\begin{equation}\label{eq:volumes}
\vol(\Iw)=\frac1{q+1},\qquad
\vol(K_n)=\frac1{\abs{\GL_2(k)}\,q^{4(n-1)}}
=\frac1{q(q-1)^2(q+1)q^{4(n-1)}}.
\end{equation}
Let $B=TU$ be the Borel subgroup of upper triangular matrices and
$\delta_B$ its modulus character. For a unitary character $\chi$ of
$F^\times$ let $\cond(\chi)$ be the smallest $m\ge0$ such that $\chi$ is
trivial on $1+\pp^m$ (with $1+\pp^0=\OO^\times$). For
unitary characters $\chi_1,\chi_2$ let
$\pi(\chi_1,\chi_2)=\Ind_B^G(\delta_B^{1/2}(\chi_1\otimes\chi_2))$ be
the unitary principal series; here and below $\Ind$ denotes
unnormalized smooth induction. Let $\St$ be the Steinberg representation,
the infinite-dimensional irreducible subquotient of $\Ind_B^G\one$
\cite[(9.10.3)]{MR2234120}, and
$\St\otimes\chi=\St\otimes(\chi\circ\det)$. By the classification of
the irreducible representations of $G$ \cite[9.11~Lemma and Classification Theorem]{MR2234120}, valid
for every $F$, the representations $\pi(\chi_1,\chi_2)$ are irreducible,
$\pi(\chi_1,\chi_2)\simeq\pi(\chi_2,\chi_1)$, and every irreducible
non-supercuspidal representation of $G$ is an irreducible principal
series $\Ind_B^G(\chi_1\otimes\chi_2)$ for suitable characters $\chi_i$,
a character $\chi\circ\det$, or a twist $\St\otimes\chi$. The latter two
arise as the subquotients of reducible principal series. Hence
the discrete series of $G$, i.e., the irreducible representations with
unitary central character which are square-integrable modulo the center
of $G$, consist of $\Gsc$ and of the twists
$\St\otimes\chi$ with $\chi$ unitary, and by
\cite[Proposition~III.4.1]{MR1989693} $\Gtemp$ is the disjoint union of
$\Gsc$, of $\{\St\otimes\chi\}$ and of $\{\pi(\chi_1,\chi_2)\}$, the
latter modulo the symmetry $(\chi_1,\chi_2)\mapsto(\chi_2,\chi_1)$.

Let $\Xc$ be the group of unitary characters of $F^\times$ and
$\Xc_m=\{\chi\in\Xc:\cond(\chi)\le m\}$, $m\ge0$. These are open
subgroups; $\Xc_0$ is the circle of unramified characters, and
$\Xc_m/\Xc_0$ is the dual of $\OO^\times/(1+\pp^m)$. We fix the
Haar measure $d\chi$ on $\Xc$ with $\int_{\Xc_0}d\chi=1$, so that
\begin{equation}\label{eq:massXm}
\int_{\Xc_m}d\chi=[\OO^\times:1+\pp^m]=(q-1)q^{m-1}\quad(m\ge1).
\end{equation}

\begin{lemma}[Fixed vectors in the nonsupercuspidal tempered spectrum]\label{lem:fixed}
Let $\chi,\chi_1,\chi_2$ be unitary characters of $F^\times$ and
$n\ge1$.
\begin{enumerate}[label=\textup{(\arabic*)},ref=\arabic*]
\item\label{item:fixed-Kn} $\dim\pi(\chi_1,\chi_2)^{K_n}=q^{n-1}(q+1)$
if $\cond(\chi_1),\cond(\chi_2)\le n$ and $0$ otherwise;
$\dim(\St\otimes\chi)^{K_n}=q^{n-1}(q+1)-1$ if $\cond(\chi)\le n$ and
$0$ otherwise.
\item\label{item:fixed-I} $\dim\pi(\chi_1,\chi_2)^{\Iw}=2$ and
$\dim(\St\otimes\chi)^{\Iw}=1$ if the characters are unramified, and
these dimensions are $0$ otherwise; the same holds with $\Iw$ replaced
by $K$ and the values $2$, $1$ replaced by $1$, $0$.
\item\label{item:fixed-sc} $V^K=V^{\Iw}=0$ for every supercuspidal representation $(\pi,V)$.
\end{enumerate}
\end{lemma}

\begin{proof}
Let $\sigma$ be a character of $B$ and $J\subset K$ a compact open
subgroup. Since $G=BK$ \cite[(7.2.1)]{MR2234120}, restriction to $K$ identifies
$(\Ind_B^G\sigma)^J$ with the functions $f$ on $K$ such that
$f(bkj)=\sigma(b)f(k)$ for $b\in B\cap K$, $j\in J$, whence
\[
\dim(\Ind_B^G\sigma)^J=\#\{k\in(B\cap K)\backslash K/J:
\sigma|_{B\cap kJk^{-1}}=1\}.
\]
If $J=K_n$, which is normal in $K$, the double coset space is
$B(\OO/\pp^n)\backslash\GL_2(\OO/\pp^n)$, the projective line over
$\OO/\pp^n$, with $q^{n-1}(q+1)$ points, and the condition is that
$\sigma$ be trivial on $B\cap K_n$, whose diagonal entries range over
$1+\pp^n$. If $J=\Iw$ the double coset space is
$B(k)\backslash\GL_2(k)/B(k)=\{1,w\}$ with $w$ the nontrivial Weyl
element, and $B\cap\Iw$ and $B\cap w\Iw w^{-1}$ both have diagonal
$\OO^\times\times\OO^\times$; if $J=K$ there is one double coset and
$B\cap K$ has the same diagonal. The characters
$\sigma=\delta_B^{1/2}(\chi_1\otimes\chi_2)$ and $\sigma=\chi\circ\det$
are trivial on $U$, and $\delta_B$ is trivial on diagonal matrices with
entries in $\OO^\times$. Thus $\sigma$ is trivial on $B\cap K_n$ if and
only if $\cond(\chi_1),\cond(\chi_2)\le n$, resp.\ $\cond(\chi)\le n$,
and on $B\cap\Iw$, $B\cap w\Iw w^{-1}$, $B\cap K$ if and only if the
characters are unramified. This gives the dimensions of the $J$-fixed
vectors in $\pi(\chi_1,\chi_2)$ and in $\Ind_B^G(\chi\circ\det)$. The
latter contains $\chi\circ\det$ with quotient $\St\otimes\chi$, and
taking $J$-fixed vectors is exact \cite[2.3~Proposition]{MR2234120}; since $\det K_n=1+\pp^n$ and
$\det\Iw=\det K=\OO^\times$, the fixed vectors of $\chi\circ\det$ have
dimension $1$ if $\chi$ is trivial on $\det J$ and $0$ otherwise. This
proves \cref{item:fixed-Kn,item:fixed-I} (for $\dim\St^{\Iw}=1$ see also
\cite[14.4~Example]{MR2234120}).

For \cref{item:fixed-sc}, $V^K\subset V^{\Iw}$, and $V^{\Iw}=0$ for
supercuspidal $V$ by \cite[14.3~Proposition]{MR2234120}, applied to the
trivial character of $\Iw$.
\end{proof}

We next record the Plancherel formula with its normalization constants.
By \cite[Theorem~VIII.1.1]{MR1989693}, there are constants $\Cps,\Cst>0$ such that for all
$\varphi\in\Hk(G)$
\begin{equation}\label{eq:plancherelgl2}
\begin{aligned}
(\varphi/dg)(1)={}&\int_{\Gsc}\tr\pi(\varphi)\,d\mu_G(\pi)
+\Cst\int_{\Xc}\tr(\St\otimes\chi)(\varphi)\,d\chi\\
&+\frac{\Cps}{2}\int_{\Xc\times\Xc}\mu(\chi_1\chi_2^{-1})
\tr\pi(\chi_1,\chi_2)(\varphi)\,d\chi_1\,d\chi_2,
\end{aligned}
\end{equation}
where $\mu$ is Harish-Chandra's $\mu$-function of the unitary principal series.
Up to a positive constant, $\mu$ is the inverse of the scalar by
which the composite of the two standard intertwining operators between
$\pi(\chi_1,\chi_2)$ and $\pi(\chi_2,\chi_1)$ acts
\cite[\S\S IV.3 and V.2]{MR1989693}. By \cite[Proposition~4.5.10]{MR1431508},
this scalar is a product of Tate $\gamma$-factors of $\chi=\chi_1\chi_2^{-1}$ and $\chi^{-1}$. Hence, by
Tate's local functional equation \cite[\S23]{MR2234120}, $\mu$ depends only on $\chi$ and is given,
for a constant $\Cmu>0$ depending on the normalization of measures, by
\begin{equation}\label{eq:mu}
\mu(\chi)=\Cmu\,q^{\cond(\chi)}\,
\frac{L(1,\chi)L(1,\chi^{-1})}{L(0,\chi)L(0,\chi^{-1})}
=\begin{cases}
\Cmu\,\dfrac{\abs{1-z}^2}{\abs{1-q^{-1}z}^2}
&\text{if $\cond(\chi)=0$ and $z=\chi(\varpi)$,}\\[1.5ex]
\Cmu\,q^{\cond(\chi)}&\text{if $\cond(\chi)\ge1$.}
\end{cases}
\end{equation}
Taking $\varphi=e_{K_n}$ in \eqref{eq:plancherelgl2}, and using
$\tr\pi(e_{K_n})=\dim\pi^{K_n}$ and $(e_{K_n}/dg)(1)=\vol(K_n)^{-1}$, we get the decomposition
\begin{equation}\label{eq:threeparts}
1=\scd{K_n}+\Pst+\Pps
\end{equation}
of the normalized Plancherel mass of the $K_n$-fixed vectors into its
supercuspidal, Steinberg and principal series parts, where
\begin{align*}
\Pst&=\vol(K_n)\Cst\int_{\Xc}\dim(\St\otimes\chi)^{K_n}\,d\chi,\\
\Pps&=\vol(K_n)\frac{\Cps}{2}\int_{\Xc\times\Xc}\mu(\chi_1\chi_2^{-1})
\dim\pi(\chi_1,\chi_2)^{K_n}\,d\chi_1\,d\chi_2.
\end{align*}

\begin{proposition}[Exact Plancherel masses for $\GL_2$]\label{prop:gl2exact}
For every non-archimedean local field $F$ and every $n\ge1$,
\[
\Pps=\frac12,\qquad
\Pst=\frac{q^{n-1}(q+1)-1}{q(q+1)q^{3n-3}},\qquad
\scd{K_n}=\frac12-q^{1-2n}+\frac1{(q+1)q^{3n-2}}.
\]
In particular $\lim_n\scd{K_n}=\frac12=\elld[\GL_2]{\Mat_2(\OO)}$, and
$q^{-2n}<\frac12-\scd{K_n}<q^{1-2n}$ for every $n\ge1$.
\end{proposition}

\begin{proof}
We first determine $\Cps\Cmu$ and $\Cst$ by applying
\eqref{eq:plancherelgl2} to $e_K$ and $e_{\Iw}$.
Write $\chi(\varpi)=z=e^{i\theta}$
for $\chi\in\Xc_0$. Expanding $\abs{1-z}^2=2-z-\bar z$ and
$\abs{1-q^{-1}z}^{-2}=\sum_{j,l\ge0}q^{-j-l}z^j\bar z^l$ and taking
the constant term,
\begin{equation}\label{eq:I0}
\int_{\Xc_0}\mu(\chi)\,d\chi
=\frac{\Cmu}{2\pi}\int_0^{2\pi}\frac{\abs{1-e^{i\theta}}^2}{\abs{1-q^{-1}e^{i\theta}}^2}\,d\theta
=\Cmu\Bigl(2\sum_{j\ge0}q^{-2j}-2\sum_{j\ge0}q^{-2j-1}\Bigr)
=\frac{2q}{q+1}\,\Cmu.
\end{equation}
By \cref{lem:fixed}, only the unramified principal series have $K$-fixed
vectors, of dimension $1$, so $\varphi=e_K$ gives, by Fubini and the
invariance of $d\chi_1$,
\[
1=\frac{\Cps}{2}\int_{\Xc_0\times\Xc_0}\mu(\chi_1\chi_2^{-1})\,d\chi_1\,d\chi_2
=\frac{\Cps}{2}\int_{\Xc_0}\mu(\chi)\,d\chi=\frac{q}{q+1}\,\Cps\Cmu,
\]
i.e., $\Cps\Cmu=(q+1)/q$. Similarly, $\varphi=e_{\Iw}$ gives
\[
q+1=\frac{\Cps}{2}\cdot2\int_{\Xc_0}\mu(\chi)\,d\chi+\Cst\int_{\Xc_0}d\chi
=2+\Cst,
\]
i.e., $\Cst=q-1$.

To evaluate the principal-series mass, put
$S_m=\int_{\Xc_m}\mu(\chi)\,d\chi$ for $m\ge1$. We claim that
\begin{equation}\label{eq:Sm}
S_m=\frac{(q-1)q^{2m}}{q+1}\,\Cmu.
\end{equation}
By \eqref{eq:mu}, \eqref{eq:I0} and \eqref{eq:massXm}, the set
$\Xc_1\setminus\Xc_0$ has measure $q-2$ and contributes $q(q-2)\Cmu$, so
$S_1=\bigl(\frac{2q}{q+1}+q(q-2)\bigr)\Cmu=\frac{(q-1)q^2}{q+1}\Cmu$; and
for $m\ge2$ the set $\Xc_m\setminus\Xc_{m-1}$ has measure
$(q-1)^2q^{m-2}$ and contributes $(q-1)^2q^{2m-2}\Cmu$, which is
$\frac{(q-1)q^{2m}}{q+1}\Cmu-\frac{(q-1)q^{2m-2}}{q+1}\Cmu$. This proves
\eqref{eq:Sm} by induction.

We can now evaluate the two nonsupercuspidal contributions.
By \cref{lem:fixed}, Fubini and
\eqref{eq:massXm}, \eqref{eq:volumes},
\begin{align*}
\Pps&=\vol(K_n)\frac{\Cps}{2}q^{n-1}(q+1)\int_{\Xc_n\times\Xc_n}\mu(\chi_1\chi_2^{-1})\,d\chi_1\,d\chi_2\\
&=\vol(K_n)\frac{\Cps\Cmu}2q^{n-1}(q+1)\cdot(q-1)q^{n-1}\cdot\frac{(q-1)q^{2n}}{q+1}
=\frac12,
\end{align*}
and
\[
\Pst=\vol(K_n)(q-1)\bigl(q^{n-1}(q+1)-1\bigr)(q-1)q^{n-1}
=\frac{q^{n-1}(q+1)-1}{q(q+1)q^{3n-3}}.
\]
The formula for $\scd{K_n}$ follows from \eqref{eq:threeparts}, and gives
$\frac12-\scd{K_n}=q^{1-2n}-\frac1{(q+1)q^{3n-2}}$, which lies strictly
between $q^{-2n}$ and $q^{1-2n}$ since $q^2-1>q^{2-n}$. Finally
$\elld[\GL_2]{\Mat_2(\OO)}=\frac12$ by \cref{prop:gl2}.
\end{proof}

For $n=1$ this gives $\scd{K_1}=\frac{q-1}{2(q+1)}$, in agreement with the
$\frac12q(q-1)$ cuspidal representations of $\GL_2(k)$, each of dimension
$q-1$ \cite[6.4~Theorem, 11.5~Theorem and 14.3~Proposition]{MR2234120}.

\subsection{\texorpdfstring{$\PGL_2$}{PGL2} in characteristic two}

\begin{example}[$\PGL_2$ in characteristic $2$]\label{ex:pgl2}
Let $\chr F=2$ and $\bG=\PGL_2$. In $\mathfrak{pgl}_2=\mathfrak{gl}_2/F$
we have $[E_{12},E_{21}]=E_{11}-E_{22}\equiv0$, since
$E_{11}+E_{22}$ is the identity. Thus the pairing
\eqref{eq:levipairing} vanishes identically, the incidence map of the
Borel subgroup is everywhere critical, and $\elld[\PGL_2]{\dLat}=1$ for
every $\dLat$ by \cref{lem:one}. Thus \cref{rem:lower} gives
$1-\scd{K_n}\ll q^{-n/4}$.

One can also see the full elliptic density directly. The trace pairing identifies
$\mathfrak{pgl}_2^*$ with $\mathfrak{sl}_2=\{\begin{psmallmatrix}a&b\\c&a\end{psmallmatrix}\}$,
and such a matrix is non-elliptic if and only if it has an $F$-rational
eigenline, i.e., if and only if $bc$ is a square in $F$. The subgroup
$F^2=\{x^2:x\in F\}$ is a closed additive subgroup of $F$ of measure zero (the image of $F^2\cap\OO$ in $\OO/\pp^N$ has
$q^{\lceil N/2\rceil}$ elements), so $\Par{\gd}$ is null by Fubini. The same
phenomenon occurs for the \oddorth{} parts $\Godd$ of all groups (\cref{lem:split}).

On the spectral side, squaring changes the distribution of conductors:
many characters $\chi$ have $\chi^2$ of smaller conductor. The
following proposition computes this distribution and the resulting
principal-series mass explicitly.
\end{example}

\begin{proposition}[Exact masses for $\PGL_2$ in characteristic two]\label{prop:pgl2}
Let $\chr F=2$, $G=\PGL_2(F)$, and let
$K_n=\ker(\PGL_2(\OO)\to\PGL_2(\OO/\pp^n))$, $n\ge1$, be the tower
\eqref{eq:tower} of the smooth model $\PGL_2$ over $\OO$
(\cref{lem:model}). Then, with $\ell=\lfloor n/2\rfloor$,
\[
\Pps=\frac{q^{-\ell}}2,\qquad
\Pst=\frac{q^{\ell}\bigl(q^{n-1}(q+1)-1\bigr)}{q(q+1)q^{3n-3}},\qquad
1-\scd{K_n}=\frac{q^{-\ell}}2\bigl(1+O(q^{-n})\bigr).
\]
\end{proposition}

\begin{proof}
The representations of $G$ are the representations of $\GL_2(F)$ with
trivial central character, and $K_n$ is the image of $1+\pp^n\Mat_2(\OO)$,
so by the classification recalled above and \cref{lem:fixed}, $\Gtemp$
consists of $\Gsc$, of the twists $\St\otimes\chi$ with $\chi^2=1$, and
of the principal series $\pi(\chi,\chi^{-1})$, $\chi\in\Xc$ modulo
$\chi\mapsto\chi^{-1}$, with $K_n$-fixed vectors of dimensions
$q^{n-1}(q+1)-1$, resp.\ $q^{n-1}(q+1)$, if $\cond(\chi)\le n$ and $0$
otherwise. We normalize $dg$ by $\vol(\PGL_2(\OO))=1$, so that
$\vol(\Iw)=(q+1)^{-1}$ and
$\vol(K_n)^{-1}=\abs{\PGL_2(\OO/\pp^n)}=q(q^2-1)q^{3(n-1)}$. By \cite[Theorem~VIII.1.1]{MR1989693}
\[
(\varphi/dg)(1)=\int_{\Gsc}\tr\pi(\varphi)\,d\mu_G(\pi)
+c_{\St}\sum_{\chi^2=1}\tr(\St\otimes\chi)(\varphi)
+\frac{\Cps'}{2}\int_{\Xc}\mu(\chi^2)\tr\pi(\chi,\chi^{-1})(\varphi)\,d\chi,
\]
with $c_{\St}$ the scalar formal degree of $\St$, and of each
$\St\otimes\chi$, with respect to $dg$ (i.e., $\fdeg{\St}=c_{\St}\,dg$), and $\mu$ as in
\eqref{eq:mu} with $\Cmu$ replaced by a constant $\Cmu'$, since
$\pi(\chi,\chi^{-1})$, viewed as a representation of $\GL_2(F)$, has
$\mu$-function $\mu(\chi\cdot\chi)$; the constants $\Cps'$, $\Cmu'$ refer
to the present normalization and differ from those of \cref{sec:gl2exact}.
The constants are determined as before: $\varphi=e_K$ gives
$1=\frac{\Cps'}{2}\int_{\Xc_0}\mu(\chi^2)\,d\chi=\frac{\Cps'}{2}\int_{\Xc_0}\mu(\chi)\,d\chi=\frac{q}{q+1}\Cps'\Cmu'$,
since $z\mapsto z^2$ preserves the Haar measure of the circle, and
$\varphi=e_{\Iw}$ gives $q+1=2+2c_{\St}$, because exactly two characters with
$\chi^2=1$ are unramified ($\chi(\varpi)=\pm1$). Hence
$\Cps'\Cmu'=\frac{q+1}q$ and $c_{\St}=\frac{q-1}2$.

It remains to compute the distribution of $\cond(\chi^2)$ on $\Xc_n$.
Squaring is injective on $\OO^\times$, and $(1+y)^2=1+y^2$, so for
$0\le j\le n$, with the convention $1+\pp^0=\OO^\times$,
\[
(1+\pp^j)^2\cap(1+\pp^n)=(1+\pp^{\max(j,\lceil n/2\rceil)})^2,
\]
and the image of $(1+\pp^j)^2$ in $\OO^\times/(1+\pp^n)$ has
order $[1+\pp^j:1+\pp^{\lceil n/2\rceil}]=q^{\lceil n/2\rceil-j}$
for $1\le j\le\lceil n/2\rceil$, and order
$(q-1)q^{\lceil n/2\rceil-1}$ for $j=0$.

A character $\chi\in\Xc_n$ satisfies $\cond(\chi^2)\le j$ if and
only if $\chi$ is trivial on $(1+\pp^j)^2$, and the characters of
$\OO^\times/(1+\pp^n)$ trivial on a subgroup of order $N$ form a
set of cardinality $(q-1)q^{n-1}/N$. Writing $\Xc_n=\Xc_n^{\circ}\cdot\Xc_0$
with $\Xc_n^\circ$ the characters trivial on $\varpi$, so that
$\cond(\chi^2)$ depends only on the $\Xc_n^\circ$-component, we get for
the number $N_j$ of $\xi\in\Xc_n^\circ$ with $\cond(\xi^2)\le j$:
\[
N_0=q^\ell,\qquad N_j=(q-1)q^{\ell+j-1}\quad(1\le j\le\lceil n/2\rceil),
\]
and $N_{\lceil n/2\rceil}=(q-1)q^{n-1}=\#\Xc_n^\circ$. Hence, by
\eqref{eq:mu}, \eqref{eq:I0} and the invariance of $d\chi$ on $\Xc_0$
under $\chi\mapsto\chi^2$,
\[
\int_{\Xc_n}\mu(\chi^2)\,d\chi
=\Cmu'\Bigl(N_0\frac{2q}{q+1}+\sum_{j=1}^{\lceil n/2\rceil}q^j(N_j-N_{j-1})\Bigr).
\]
The sum is evaluated as in the proof of \eqref{eq:Sm}: the term $j=1$
is $q^{\ell+1}(q-2)$ and the terms $j\ge2$ are $(q-1)^2q^{\ell+2j-2}$, so
that the expression in parentheses equals $q^\ell$ times the value of
$S_{\lceil n/2\rceil}/\Cmu'$, namely $(q-1)q^{\ell+2\lceil n/2\rceil}/(q+1)=(q-1)q^{n+\lceil n/2\rceil}/(q+1)$.
Therefore
\[
\Pps=\vol(K_n)\frac{\Cps'\Cmu'}2q^{n-1}(q+1)\cdot\frac{(q-1)q^{n+\lceil n/2\rceil}}{q+1}
=\frac{q^{2n+\lceil n/2\rceil-1}}{2q^{3n-1}}=\frac{q^{-\ell}}2,
\]
and, since the characters with $\chi^2=1$ and $\cond(\chi)\le n$ number
$2N_0=2q^\ell$,
\[
\Pst=\vol(K_n)\frac{q-1}2\bigl(q^{n-1}(q+1)-1\bigr)2q^\ell
=\frac{q^\ell\bigl(q^{n-1}(q+1)-1\bigr)}{q(q+1)q^{3n-3}}
\le q^{\ell-2n+1}\le 2q^{1-n}\cdot\frac{q^{-\ell}}2,
\]
since $2\ell\le n$. The last assertion follows.
\end{proof}

The two hyperspecial examples compare the general bound with the
exact order of $\abs{\scd{K_n}-\elld{\dLat}}$, for fixed $F$ as $n\to\infty$
(the implied constants depend on $q$, and for $\PGL_2$ on the parity of
$n$):
\begin{center}
\begin{tabular}{lccc}
\toprule
Group & Limit & General bound & Exact order \\
\midrule
$\GL_2$ & $1/2$ & $O(q^{-n/2})$ & $q^{-2n}$ \\
$\PGL_2$, $\chr F=2$ & $1$ & $O(q^{-n/4})$ & $q^{-n/2}$ \\
\bottomrule
\end{tabular}
\end{center}
In both examples the supercuspidal density approaches its limit
strictly from below: $\scd{K_n}<\elld{\dLat}$ for every $n$.

\bibliographystyle{amsalphapure}

\providecommand{\bysame}{\leavevmode\hbox to3em{\hrulefill}\thinspace}
\providecommand{\MR}{\relax\ifhmode\unskip\space\fi MR }
% \MRhref is called by the amsart/book/proc definition of \MR.
\providecommand{\MRhref}[2]{%
  \href{http://www.ams.org/mathscinet-getitem?mr=#1}{#2}
}
\providecommand{\href}[2]{#2}
\begin{thebibliography}{TVVW18}

\bibitem[BH06]{MR2234120}
Colin~J. Bushnell and Guy Henniart, \emph{The local {L}anglands conjecture for
  {$\rm GL(2)$}}, Grundlehren der Mathematischen Wissenschaften [Fundamental
  Principles of Mathematical Sciences], vol. 335, Springer-Verlag, Berlin,
  2006. \MR{2234120}

\bibitem[BH07]{MR2303528}
\bysame, \emph{Counting the discrete series for {${\rm GL}(n)$}}, Bull. Lond.
  Math. Soc. \textbf{39} (2007), no.~1, 133--137. \MR{2303528}

\bibitem[BK98]{MR1643417}
Colin~J. Bushnell and Philip~C. Kutzko, \emph{Smooth representations of
  reductive {$p$}-adic groups: structure theory via types}, Proc. London Math.
  Soc. (3) \textbf{77} (1998), no.~3, 582--634. \MR{1643417}

\bibitem[BM83]{MR737927}
Walter Borho and Robert MacPherson, \emph{Partial resolutions of nilpotent
  varieties}, Analysis and topology on singular spaces, {II}, {III} ({L}uminy,
  1981), Ast\'{e}risque, vol. 101--102, Soc. Math. France, Paris, 1983,
  pp.~23--74. \MR{737927}

\bibitem[Bor91]{MR1102012}
Armand Borel, \emph{Linear algebraic groups}, second ed., Graduate Texts in
  Mathematics, vol. 126, Springer-Verlag, New York, 1991. \MR{1102012
  (92d:20001)}

\bibitem[BP16]{MR3463423}
Rapha\"el Beuzart-Plessis, \emph{A short proof of the existence of
  supercuspidal representations for all reductive {$p$}-adic groups}, Pacific
  J. Math. \textbf{282} (2016), no.~1, 27--34. \MR{3463423}

\bibitem[Bum97]{MR1431508}
Daniel Bump, \emph{Automorphic forms and representations}, Cambridge Studies in
  Advanced Mathematics, vol.~55, Cambridge University Press, Cambridge, 1997.
  \MR{1431508}

\bibitem[Cas95]{CassNotes}
Bill Casselman, \emph{Introduction to admissible representations of $p$-adic
  groups}, Unpublished notes (1995),
  \url{http://www.math.ubc.ca/~cass/research/pdf/p-adic-book.pdf}.

\bibitem[CG97]{MR1433132}
Neil Chriss and Victor Ginzburg, \emph{Representation theory and complex
  geometry}, Birkh\"auser Boston, Inc., Boston, MA, 1997. \MR{1433132
  (98i:22021)}

\bibitem[CGP15]{MR3362817}
Brian Conrad, Ofer Gabber, and Gopal Prasad, \emph{Pseudo-reductive groups},
  second ed., New Mathematical Monographs, vol.~26, Cambridge University Press,
  Cambridge, 2015. \MR{3362817}

\bibitem[Chi26]{Chi26}
Jingren Chi, \emph{Equivalued affine {Springer} fibers in mixed
  characteristic}, 2026, arXiv:2606.31083, version 2.

\bibitem[Con14]{MR3362641}
Brian Conrad, \emph{Reductive group schemes}, Autour des sch\'emas en groupes.
  {V}ol. {I}, Panor. Synth\`eses, vol. 42/43, Soc. Math. France, Paris, 2014,
  pp.~93--444. \MR{3362641}

\bibitem[DeB06]{MR2214792}
Stephen DeBacker, \emph{Parameterizing conjugacy classes of maximal unramified
  tori via {B}ruhat-{T}its theory}, Michigan Math. J. \textbf{54} (2006),
  no.~1, 157--178. \MR{2214792}

\bibitem[Del76]{MR0425033}
Pierre Deligne, \emph{Le support du caract\`ere d'une repr\'esentation
  supercuspidale}, C. R. Acad. Sci. Paris S\'er. A-B \textbf{283} (1976),
  no.~4, Aii, A155--A157. \MR{0425033}

\bibitem[Dil23]{MR4617942}
Peter Dillery, \emph{Rigid inner forms over local function fields}, Adv. Math.
  \textbf{430} (2023), Paper No. 109204, 100. \MR{4617942}

\bibitem[Fin21a]{MR4357723}
Jessica Fintzen, \emph{On the construction of tame supercuspidal
  representations}, Compos. Math. \textbf{157} (2021), no.~12, 2733--2746.
  \MR{4357723}

\bibitem[Fin21b]{MR4199732}
\bysame, \emph{Types for tame {$p$}-adic groups}, Ann. of Math. (2)
  \textbf{193} (2021), no.~1, 303--346. \MR{4199732}

\bibitem[FS21]{MR4325284}
Jessica Fintzen and Sug~Woo Shin, \emph{Congruences of algebraic automorphic
  forms and supercuspidal representations}, Camb. J. Math. \textbf{9} (2021),
  no.~2, 351--429. \MR{4325284}

\bibitem[GH26]{GH26}
Itay Glazer and Yotam~I. Hendel, \emph{A lower bound on the analytic
  log-canonical threshold over local fields of positive characteristic}, Proc.
  Amer. Math. Soc. (2026), Electronically published September 30, 2026; DOI
  10.1090/proc/17891; arXiv:2511.01270.

\bibitem[GKM06]{MR2209851}
Mark Goresky, Robert Kottwitz, and Robert MacPherson, \emph{Purity of
  equivalued affine {S}pringer fibers}, Represent. Theory \textbf{10} (2006),
  130--146 (electronic). \MR{2209851 (2007i:22025)}

\bibitem[Glo13]{MR3057121}
Helge Gl\"{o}ckner, \emph{Invariant manifolds for analytic dynamical systems
  over ultrametric fields}, Expo. Math. \textbf{31} (2013), no.~2, 116--150.
  \MR{3057121}

\bibitem[Gre66]{MR0207700}
Marvin~J. Greenberg, \emph{Rational points in {H}enselian discrete valuation
  rings}, Inst. Hautes \'Etudes Sci. Publ. Math. (1966), no.~31, 59--64.
  \MR{0207700 (34 \#7515)}

\bibitem[How74]{MR0342645}
Roger Howe, \emph{The {F}ourier transform and germs of characters (case of
  {${\rm Gl}_{n}$} over a {$p$}-adic field)}, Math. Ann. \textbf{208} (1974),
  305--322. \MR{0342645 (49 \#7391)}

\bibitem[How77]{MR0492088}
Roger~E. Howe, \emph{Some qualitative results on the representation theory of
  {${\rm Gl}_{n}$} over a {$p$}-adic field}, Pacific J. Math. \textbf{73}
  (1977), no.~2, 479--538. \MR{0492088 (58 \#11242)}

\bibitem[Igu00]{MR1743467}
Jun-ichi Igusa, \emph{An introduction to the theory of local zeta functions},
  AMS/IP Studies in Advanced Mathematics, vol.~14, American Mathematical
  Society, Providence, RI, 2000. \MR{1743467 (2001j:11112)}

\bibitem[Kal19]{MR4013740}
Tasho Kaletha, \emph{Regular supercuspidal representations}, J. Amer. Math.
  Soc. \textbf{32} (2019), no.~4, 1071--1170. \MR{4013740}

\bibitem[Kim04]{MR2075916}
Ju-Lee Kim, \emph{Dual blobs and {P}lancherel formulas}, Bull. Soc. Math.
  France \textbf{132} (2004), no.~1, 55--80. \MR{2075916 (2005f:22025)}

\bibitem[Kim07]{MR2276772}
\bysame, \emph{Supercuspidal representations: an exhaustion theorem}, J. Amer.
  Math. Soc. \textbf{20} (2007), no.~2, 273--320 (electronic). \MR{2276772
  (2008c:22014)}

\bibitem[KP23]{MR4520154}
Tasho Kaletha and Gopal Prasad, \emph{Bruhat-{T}its theory---a new approach},
  New Mathematical Monographs, vol.~44, Cambridge University Press, Cambridge,
  2023. \MR{4520154}

\bibitem[KST20]{MR4046074}
Ju-Lee Kim, Sug~Woo Shin, and Nicolas Templier, \emph{Asymptotic behavior of
  supercuspidal representations and {S}ato-{T}ate equidistribution for
  families}, Adv. Math. \textbf{362} (2020), 106955, 57. \MR{4046074}

\bibitem[Mil17]{MR3729270}
J.~S. Milne, \emph{Algebraic groups}, Cambridge Studies in Advanced
  Mathematics, vol. 170, Cambridge University Press, Cambridge, 2017, The
  theory of group schemes of finite type over a field. \MR{3729270}

\bibitem[Moy17]{MR3753921}
Allen Moy, \emph{Computations with {B}ernstein projectors of {${\rm SL}(2)$}},
  Representation theory, number theory, and invariant theory, Progr. Math.,
  vol. 323, Birkh\"auser/Springer, Cham, 2017, pp.~503--533. \MR{3753921}

\bibitem[MP94]{MR1253198}
Allen Moy and Gopal Prasad, \emph{Unrefined minimal {$K$}-types for {$p$}-adic
  groups}, Invent. Math. \textbf{116} (1994), no.~1-3, 393--408. \MR{1253198
  (95f:22023)}

\bibitem[MP96]{MR1371680}
\bysame, \emph{Jacquet functors and unrefined minimal {$K$}-types}, Comment.
  Math. Helv. \textbf{71} (1996), no.~1, 98--121. \MR{1371680 (97c:22021)}

\bibitem[Oes82]{MR656627}
Joseph Oesterl{\'e}, \emph{R\'eduction modulo {$p^{n}$} des sous-ensembles
  analytiques ferm\'es de {${\bf Z}^{N}_{p}$}}, Invent. Math. \textbf{66}
  (1982), no.~2, 325--341. \MR{656627 (83j:12014)}

\bibitem[PR94]{MR1278263}
Vladimir Platonov and Andrei Rapinchuk, \emph{Algebraic groups and number
  theory}, Pure and Applied Mathematics, vol. 139, Academic Press Inc., Boston,
  MA, 1994, Translated from the 1991 Russian original by Rachel Rowen.
  \MR{1278263 (95b:11039)}

\bibitem[Ser81]{MR644559}
Jean-Pierre Serre, \emph{Quelques applications du th\'eor\`eme de densit\'e de
  {C}hebotarev}, Inst. Hautes \'Etudes Sci. Publ. Math. (1981), no.~54,
  323--401. \MR{644559 (83k:12011)}

\bibitem[Ser06]{MR2179691}
\bysame, \emph{Lie algebras and {L}ie groups}, Lecture Notes in Mathematics,
  vol. 1500, Springer-Verlag, Berlin, 2006, 1964 lectures given at Harvard
  University, Corrected fifth printing of the second (1992) edition.
  \MR{2179691 (2006e:17001)}

\bibitem[SS70]{MR0268192}
T.~A. Springer and R.~Steinberg, \emph{Conjugacy classes}, Seminar on
  {A}lgebraic {G}roups and {R}elated {F}inite {G}roups ({T}he {I}nstitute for
  {A}dvanced {S}tudy, {P}rinceton, {N}.{J}., 1968/69), Lecture Notes in
  Mathematics, Vol. 131, Springer, Berlin, 1970, pp.~167--266. \MR{0268192 (42
  \#3091)}

\bibitem[ST26]{ST26}
Loren Spice and Cheng-Chiang Tsai, \emph{Jordan decompositions in {Lie}
  algebras and their duals}, 2026, arXiv:2601.07168.

\bibitem[SVP26]{SVP26}
Jiahe Shen and Roger Van~Peski, \emph{Eigenvalues of {$p$}-adic random
  matrices}, 2026, arXiv:2601.06283.

\bibitem[TVVW18]{MR3803728}
David Taylor, V.~S. Varadarajan, Jukka Virtanen, and David Weisbart,
  \emph{Temperedness of measures defined by polynomial equations over local
  fields}, Pacific J. Math. \textbf{296} (2018), no.~1, 227--256. \MR{3803728}

\bibitem[Wal03]{MR1989693}
J.-L. Waldspurger, \emph{La formule de {P}lancherel pour les groupes
  {$p$}-adiques (d'apr\`es {H}arish-{C}handra)}, J. Inst. Math. Jussieu
  \textbf{2} (2003), no.~2, 235--333. \MR{1989693}

\bibitem[Xue09]{MR2558787}
Ting Xue, \emph{Nilpotent orbits in the dual of classical {L}ie algebras in
  characteristic 2 and the {S}pringer correspondence}, Represent. Theory
  \textbf{13} (2009), 609--635. \MR{2558787}

\bibitem[Yu01]{MR1824988}
Jiu-Kang Yu, \emph{Construction of tame supercuspidal representations}, J.
  Amer. Math. Soc. \textbf{14} (2001), no.~3, 579--622 (electronic).
  \MR{1824988 (2002f:22033)}

\end{thebibliography}


\end{document}
